\documentclass[12pt]{extarticle} %
\usepackage{times}
\usepackage{url}            %
\usepackage{booktabs}       %
\usepackage{amsfonts}       %
\usepackage{nicefrac}       %
\usepackage{microtype}      %
\usepackage{mkolar_definitions}
\usepackage{xspace}
\usepackage{multirow}
\usepackage{amsmath}
\usepackage{algorithm}
\usepackage{algorithmic}
\usepackage{color}
\usepackage{authblk}
\usepackage{color, colortbl}
\usepackage{enumitem}
\usepackage{comment}
\usepackage{bm}
\usepackage{subcaption}
\usepackage[left=2cm,top=2cm,right=2cm]{geometry}


\usepackage[scaled=.9]{helvet}
\usepackage{thmtools}
\newtheorem{myexp}{Example}

\usepackage{url}
\usepackage{authblk}
\usepackage{csquotes}
\usepackage{tikz}

%\usepackage[dvipsnames]{xcolor}
\usepackage[colorlinks=true, linkcolor=blue, citecolor=blue]{hyperref}

%%%\usepackage[style=alphabetic,natbib=true]{biblatex}
\usepackage{natbib}
%\bibliographystyle{plainnat}
%\bibpunct{(}{)}{;}{a}{,}{,}

\newenvironment{packed_enum}{
  \begin{enumerate}
    \setlength{\itemsep}{1pt}
    \setlength{\parskip}{-1pt}
    \setlength{\parsep}{0pt}
}{\end{enumerate}}

\usepackage{centernot}
\newcommand{\bigCI}{\mathrel{\text{\scalebox{1.07}{$\perp\mkern-10mu\perp$}}}}
\newcommand{\nbigCI}{\centernot{\bigCI}}

\newcount\Comments  %
\Comments=1 %
\definecolor{darkgreen}{rgb}{0,0.5,0}
\definecolor{darkred}{rgb}{0.7,0,0}
\definecolor{teal}{rgb}{0.3,0.8,0.8}
\definecolor{orange}{rgb}{1.0,0.5,0.0}
\definecolor{purple}{rgb}{0.8,0.0,0.8}
\newcommand{\kibitz}[2]{\ifnum\Comments=1{\textcolor{#1}{\textsf{\footnotesize #2}}}\fi}
\usetikzlibrary{positioning}

\definecolor{Gray}{gray}{0.9}

\newcommand{\brank}{R_B} \newcommand{\mrank}{R_M}
\newcommand{\olive}{\textsc{Olive}}
\newcommand{\MLE}{\mathrm{MLE}}
\newcommand{\Prr}{\mathrm{Pr}}
\newcommand{\VM}{\mathrm{VM}}
\newcommand{\SIS}{\mathrm{SIS}}
\newcommand{\error}{\ensuremath{\textsc{Error}}}
\newcommand{\Bayes}{\text{Bayes}}
\newcommand{\masa}[1]{\noindent{\textcolor{purple}{\{{\bf Masa:} \em #1\}}}}
\newcommand{\wen}[1]{\noindent{\textcolor{red}{\{{\bf WS:} \em #1\}}}}
\newcommand{\ming}[1]{\noindent{\textcolor{ProcessBlue}{\{{\bf Ming:} \em #1\}}}}
\newcommand{\newedit}{\color{blue}}
\newcommand{\rI}{\mathrm{I}}
\newcommand{\nan}[1]{\noindent{\textcolor{olive}{\{{\bf Nan:} \em #1\}}}}

\newcommand{\zihao}[1]{\noindent{\textcolor{red}{\{{\bf Zihao:} \em #1\}}}}

\newcommand{\hlan}[1]{\noindent{\textcolor{pink}{\{{\bf Hui:} \em #1\}}}}


\newcommand{\defeq}{\triangleq}
\newcommand{\Otilde}{\tilde{O}}
\newcommand{\KL}{\mathrm{KL}}
\newcommand{\sa}{(s,a)}
\newcommand{\misp}{\epsilon_{\text{misspec}}}
\newcommand{\gproj}{g_{\text{proj}}}
\newcommand{\E}{\mathbb{E}}
\newcommand{\error}{\ensuremath{\textsc{Error}}}

\renewcommand{\arraystretch}{2.}
\newtheorem{assumption}[theorem]{Assumption}



\begin{document}





\section{MLE-based Iterated TIKHNOV Regularized Estimator}
 We consider a nonparametric IV problem: 
\begin{align}\label{eq:inverse-iv}
    \mathbb{E}[Y-f(X)|Z]=0. 
\end{align}
We aim to introduce a new oracle-efficient method that can incorporate any function approximation. While there are many minimax estimation methods \citep{DikkalaNishanth2020MEoC} that can incorporate any function approximation, (a) these methods are hard to use in practice due to minimax optimization when using neural networks, (b) it is unclear how to perform cross-validation to choose function classes.  

Here is our procedure. 
\begin{itemize}    
    \item Introduce a function class: $\Gcal \subset [\mathcal{X}\times \mathcal{Z} \to \mathbb{R}]$
    \begin{align*}
        \hat g= \argmin_{g \in \mathcal{G}}\EE_n[-\log g(X|Z)]
    \end{align*}
     
    \item Introduce a function class: $\Fcal \subset [\mathcal{X} \to \mathbb{R}]$. We iterate according to the following procedure: 
    \begin{align}\label{eq:main-iter}
         \hat f_m =\argmin_{f \in \Fcal} 
       \EE_n[ - 2Y \hat \EE[f(X)|Z] +   \{\hat \EE[f(X)|Z]\}^2  ]   
         +\alpha \EE_n[\{f-\hat f_{m-1}\}^2(X)],
    \end{align}
     where $\int f(x)\hat g(x|Z)\mathrm{d}x = \mathrm{\hat \EE}[f(X)|Z]$.
\end{itemize}
For the boundedness of the function classes $\Fcal$ and $\Gcal$, we make the following assumption. \begin{assumption}
    We assume that the $\Fcal$ and $\Gcal$ are bounded by an absolute constant $C$ in $L_\infty$ norm.
\end{assumption}
Our methods rely on just standard ERM oracles. Besides, cross-validation is easy. 
Note in the latter optimization, we are aiming to solve 
\begin{align*}
        f^{\star}_{\alpha} &= \argmin_{f} \EE[\{h^{\star}(Z)  - (\Tcal f)(Z) \}^2 +\alpha f^2(X) ] 
\end{align*}
where $h^{\star}(Z)=\EE[Y|Z]$ and $\Tcal(f)(Z)=\EE[f(X)|Z]$. 

\begin{remark}
    Deep IV \citep{hartford2017deep} does not have any analysis. No explicit TIKHNOV regularization.  No iteration. 
\end{remark}
\begin{assumption}[*Smooth Function Class]\label{ass:smooth-func}
    We assume that the function space $\Fcal -f_0 $ satisfies
a smoothness assumption that implies a relationship between the $L_\infty$ and $L_2$ norms, i.e. for some $\gamma \leq 1$ we have $$
\|f - f_\alpha^*\|_\infty \leq O(\|f - f_\alpha^*\|_2^\gamma).
$$
\end{assumption}
Such a relationship is known for instance to hold for Sobolev spaces and more generally
for reproducing kernel Hilbert spaces (RKHS) with a polynomial eigendecay.
We also introduce the $\beta$- source condition, which is commonly used in the literature on inverse problems and captures how strongly we can identify the function $f_0$ by the data we observed. 
\begin{assumption}[Full Support]\label{ass:full-supp}
    We assume that the ground truth is lower bounded, i.e. has a full support, $$
    g^*(x\mid z) \geq c_0 >0.
    $$
\end{assumption}
The next assumption might also be helpful. \begin{assumption}\label{ass:bound-2-4}
    For every $f \in \Fcal$, we have $\|f-f_\alpha^*\|_4 \leq C_{2,4}\|f-f_\alpha^*\|_2$
\end{assumption}
\begin{definition}[$\beta$- source condition]\label{def:source-cond}
    There exists $w_0\in\bar{\Fcal}$ such that $f_0$ satisfies $$
    f_0 = (\Tcal^*\Tcal)^{\beta/2} w_0.
    $$
\end{definition}
Essentially, this assumption is claiming $
f_0 \in \Rcal(\Tcal^*\Tcal)^{\beta/2}$. Intuitively, when the parameter $\beta$ is large, the function $ f_0$ exhibits greater smoothness, in the sense that the $L_2$ -inner product of $f_0$ with eigenfunctions that have smaller eigenvalues relative to the operator $\Tcal$ tend to decay faster.

\section{Theory}

Let $\delta^2_n = O(\log(N_{[]}(\Fcal, \|\cdot\|_\infty, 1/n)/\zeta)/n)$ . Hereafter, every norm is defined w.r.t. the random variable $(X,Y,Z)$. 

First, we define the Hellinger distance between two distributions with probability density $p(x)$ and $q(x)$ respectively. The Hellinger distance between the two distributions is defined by: $$
H(P, Q)^2 = \int \big(\sqrt{p(x)} - \sqrt{q(x)}\big)^2 d(x).
$$
By invoking MLE guarantee (\citet[Theorem 21]{agarwal2020flambe}, \citet[Theorem 3.3]{ge2023provable}), we have 
\begin{lemma}
With probability $1-\zeta$, we have 
\begin{align*}
    \EE[H^2(\hat g(\cdot |Z),g^{\star}(\cdot \mid z))]\leq \delta^2_n. 
\end{align*}
Here the expectation is taken on the random variable $Z$.
\end{lemma}
Therefore, letting  $\hat \EE[f(X)|Z]=\hat \Tcal f(Z)$ and $\EE[f(X)|Z]=\Tcal f(Z)$, we have 
\begin{lemma}\label{lem:appro}
Under Assumption \ref{ass:full-supp}, for all $f\in\Fcal  $, we have  $ \|(\hat \Tcal-\Tcal)f\|_1\leq   2(1/c_0 +1)\delta_n \cdot\|f\|_2$ and $\|(\hat \Tcal-\Tcal)f\|_2 \leq 2(C/c_0 +1)^{1/2}\|f\|_{\infty}^{1/2}\|f\|_{2}^{1/2} \delta_n^{1/2}$
\end{lemma}
\begin{proof}
\begin{align*}
   \|(\hat \Tcal-\Tcal)f\|_1  &= \EE_{z \sim g^{\star}(z)}\left[\bigg|\EE_{x \sim g^{\star}(x|z)}\left [\frac{\hat g(x|z)}{g^{\star}(x|z)}f(x) - f(x) \right] \bigg| \right]\\
    &\leq \EE_{z \sim g^{\star}(z),x \sim g^{\star}(x|z)}\left[ \bigg|\frac{\hat g(x|z)}{g^{\star}(x|z)}f(x) - f(x)  \bigg| \right]\\
    &\leq \EE_{z \sim g^{\star}(z),x \sim g^{\star}(x|z)}\left[\sqrt{\frac{\hat g(x|z)}{g^{\star}(x|z)}}|f(x)| \bigg|\sqrt{\frac{\hat g(x|z)}{g^{\star}(x|z)}} - 1 \bigg|\right] \\
    &\qquad\qquad+ \EE_{z \sim g^{\star}(z),x \sim g^{\star}(x|z)}\left[|f(x)|\bigg|\sqrt{\frac{\hat g(x|z)}{g^{\star}(x|z)}} - 1  \bigg|\right] \\
    &\leq \EE\bigg[\frac{\hat g(x|z)}{g(^{\star}(x|z)}f^2(x) \bigg]^{1/2} \times \EE[2H^2(\hat g(\cdot|z),g^{\star}(\cdot |z))  ]\\
    &\qquad\qquad+\EE\bigg[\frac{\hat g(x|z)}{g(^{\star}(x|z)}f^2(x) \bigg]^{1/2} \times \EE[2H^2(\hat g(\cdot|z),g^{\star}(\cdot |z))  ] \tag{CS inequality} \\
    &\leq  2\{C/c_0+1\} \EE[f^2(x) ]^{1/2} \times \EE[2H^2(\hat g(\cdot|z),g^{\star}(\cdot |z))  ]^{1/2}\\
    &\leq 2(C/c_0 +1)\cdot\delta_n \cdot\|f\|_2 .
\end{align*} 
The second inequality comes directly from $$
\|(\hat{\Tcal}-  \Tcal)f\|_2 \leq \sqrt{\|(\hat{\Tcal}-  \Tcal)f\|_1} \|f\|_\infty^{1/2}.
$$
\end{proof}
\begin{remark}
    When $f$ belongs to $\{\Gcal, \Fcal\}$, we have $\|(\hat{\Tcal} - \Tcal)f\|_2 \leq O(\|f\|_2^{1/2} \delta_n^{1/2})$. Specifically, under Assumption \ref{ass:smooth-func}, we have $\|(\hat{\Tcal} - \Tcal)f\|_2 \leq O(\|f\|_2^{(1+\gamma)/2} \cdot\delta_n^{1/2})$.
\end{remark}
\begin{lemma}\label{lem:approx-2}
    Under Assumption \ref{ass:full-supp}, \ref{ass:bound-2-4}, for all $f\in\Fcal - f_0$, we have $\|(\hat \Tcal-\Tcal)f\|_2 \leq (C_{2,4}C)^{1/2}\cdot (C/c_0 +1)\|f\|_{2}\cdot \delta_n^{1/2}$.
\end{lemma}
\begin{proof}
\zihao{A brief proof with Cauchy Schwarz}
    \begin{align*}
    & \|(\hat \Tcal - \Tcal)f\|_2^2\\
    & = \EE_{z\sim g^*(z)}\bigg[\bigg|\EE_{x\sim g^*(x\mid z)}\bigg[\bigg(\frac{\hat{g}(x\mid z)}{g^*(x\mid z)} - 1\bigg)f(x)\bigg]\bigg|^2\bigg]\\
    &\leq \EE[f^4(x)]^{1/2}\EE[\{\hat g(x|z)/g^{\star}(x|z)-1\}^4]^{1/2}\\
    &\leq C_{2,4}\cdot C(C/c_0 +1)^2\cdot\|f\|_2^2\cdot \EE_{z\sim g^*(z)}[H(g(\cdot|z),g^{\star}(\cdot|z) )^2]^{1/2}\\
    &\leq C_{2,4}\cdot C(C/c_0 +1)^2\cdot\|f\|_2^2\cdot \delta_n
\end{align*}
\end{proof}
\zihao{We can get a tighter bound if we assume $\|\hat{g}/g^* - 1\|_4 \leq C_{2,4}\|\hat{g}/g^* - 1\|_2$}
\masa{Can we improve in this way?
Suppose $g^{\star}$ is lower-bounded by $c_0$. 
{\small 
\begin{align*}
    &\|(\hat \Tcal-\Tcal)f\|_1 \\
    &= \EE_{z \sim g^{\star}(z)}\left[|\EE_{x \sim g^{\star}(x|z)}\left [\frac{\hat g(x|z)}{g^{\star}(x|z)}f(x) - f(x) \right] | \right]\\
    &\leq \EE_{z \sim g^{\star}(z),x \sim g^{\star}(x|z)}\left[ |\frac{\hat g(x|z)}{g^{\star}(x|z)}f(x) - f(x)  | \right]\\
    &\leq \EE_{z \sim g^{\star}(z),x \sim g^{\star}(x|z)}\left[\sqrt{\frac{\hat g(x|z)}{g^{\star}(x|z)}}|f(x)| |\sqrt{\frac{\hat g(x|z)}{g^{\star}(x|z)}} - 1 |\right] + \EE_{z \sim g^{\star}(z),x \sim g^{\star}(x|z)}\left[|f(x)||\sqrt{\frac{\hat g(x|z)}{g^{\star}(x|z)}} - 1  |\right] \\
    &\leq \EE[\frac{\hat g(x|z)}{g(^{\star}(x|z)}f^2(x) ]^{1/2} \times \EE[2H^2(\hat g(\cdot|z),g^{\star}(\cdot |z))  ] +\EE[\frac{\hat g(x|z)}{g(^{\star}(x|z)}f^2(x) ]^{1/2} \times \EE[2H^2(\hat g(\cdot|z),g^{\star}(\cdot |z))  ] \tag{CS inequality} \\
    &\leq  2\{1/c_0+1\} \EE[f^2(x) ]^{1/2} \times \EE[2H^2(\hat g(\cdot|z),g^{\star}(\cdot |z))  ]^{1/2}. 
\end{align*} 
} 
Then, it looks like we can get the faster rate even if we do not assume Assumption 1. Then, we can use 
\begin{align*}
    \|(\hat \Tcal-\Tcal)f\|_2\leq \|f\|_{\infty}\sqrt{    \|(\hat \Tcal-\Tcal)f\|_1 }. 
\end{align*}
We might be able to directly upper-bounded $ \|(\hat \Tcal-\Tcal)f\|_2$ as well. Another way is
\begin{align*}
    & \|(\hat \Tcal - \Tcal)f\|_2\\
    &\leq \EE[f^4(x)]^{1/2}\EE[\{\hat g(x|z)/g^{\star}(x|z)-1\}^4]^{1/2}
\end{align*}
Under the assumption $\|f-f_{\alpha}\|_4 \leq C_{2,4}\|f-f_{\alpha}\|_2$ (this is commonly assumed. Check Assumption 8 in Orthogonal statistical machine learning, can we say a good thing? 
}

\masa{Another idea. Let's use the following loss function:
\begin{align*}
    \argmin_{g \in \Gcal}\EE_n[0.5\int g^2(x|Z)d\mu(x)]-\EE_n[g(x|z)].  
\end{align*}
where $\mu$ is the Lebsgue measure. Note each $g$ does not need to be normalized in this case. 
As in the Wainright book (Chapter 14), we can easily prove $$\EE_{z \sim g^{\star}(z)}\left[\{\int |\hat g(x|z)-g^{\star}(x|z) \mathrm{d}\mu(x)\}^2\right]=O(1/n)$$
Then, under the assumption $g^{\star}(x|z)>c_0$, we can get a result
\begin{align*}
    \|(\hat \Tcal-\Tcal)f\|_2\leq ?? 
\end{align*}
}

\iffalse 
Similarly, by invoking regression guarantee and letting $\EE[Y|Z]=h^{\star}(Z)$, 
\begin{lemma}
With probability $1-\delta$, we have 
\begin{align*}
        | \EE_n[Y \hat h(Z) - 0.5\hat h(Z)^2] -  \EE_n[Y h^{\star}(Z) - 0.5 h^{\star}(Z)^2] |
        \leq \delta^2_n.  
    \end{align*}
\end{lemma}






\fi 
We also need the following concentration inequality lemma which is used throughout the proofs. 
\begin{lemma}[Localized Concentration, \cite{foster2019orthogonal}]\label{lem:local-concen}
For any $f \in \mathcal{F}:=\times_{i=1}^d \mathcal{F}_i$ be a multivalued outcome function, that is almost surely absolutely bounded by a constant. Let $\ell(Z ; f(X)) \in \mathbb{R}$ be a loss function that is $O(1)$-Lipschitz in $h(X)$, with respect to the $\ell_2$ norm. Let $\delta_n=\Omega\left(\sqrt{\frac{d \log \log (n)+\log (1 / \zeta)}{n}}\right)$ be an upper bound on the critical radius of $\operatorname{star}\left(\mathcal{F}_i\right)$ for $i \in[d]$. Then for any fixed $h_* \in \mathcal{F}$, w.p. $1-\zeta$ :
$$
\forall f \in \mathcal{F}:\left|\left(\mathbb{E}_n-\mathbb{E}\right)\left[\ell(Z ; f(X))-\ell\left(Z ; f_0(X)\right)\right]\right|=O\left(d \delta_n \sum_{i=1}^d\left\|f_i-f_{i, 0}\right\|_{L_2}+d \delta_n^2\right)
$$
If the loss is linear in $f(X)$, i.e. $\ell\left(Z ; f(X)+f^{\prime}(X)\right)=\ell(Z ; f(X))+\ell\left(Z ; f^{\prime}(X)\right)$ and $\ell(Z ; \alpha f(X))=$ $\alpha \ell(Z ; f(X))$ for any scalar $\alpha$, then it suffices that we take $\delta_n=\Omega\left(\sqrt{\frac{\log (1 / \zeta)}{n}}\right)$ that upper bounds the critical radius of $\operatorname{star}\left(\mathcal{F}_i\right)$ for $i \in[d]$.
\end{lemma}
\begin{proof}
    For a detailed proof, please refer to \cite{foster2019orthogonal}.
\end{proof}
\subsection{Analysis for Non-iterative Version}

First, consider the case when $m=1$. According to the iteration in \eqref{eq:main-iter}, we want to analyze 
\begin{align}\label{eq:single-iter}
         \hat f  =\argmin_{f \in \Fcal} 
       \EE_n[ - 2 Y \hat \EE[f(X)|Z] +  \{\hat \EE[f(X)|Z]\}^2  ]   
         +\alpha \EE_n[f^2(X)]. 
    \end{align}
By Definition \ref{def:source-cond}, our goal is to prove 
\begin{align*}
    \|\hat f-f_0 \|^2_2\leq  \delta_n / \alpha^2+ \alpha^{\min(\beta,2)},
\end{align*}
where $f_0$ is the least norm solution in the inverse nonparametric IV problem in  \eqref{eq:inverse-iv}. 
\iffalse 
When $m\geq 2$, this can be extended to 
\begin{align*}
    \|\hat f_m-f_0 \|^2_2\leq  \delta^2_n/\alpha^2 +  \alpha^{\min(\beta,2m)} +     \|\hat f_{m-1}-f_{0} \|_2   \|\hat f_m-f_{0} \|_2
\end{align*}
\fi 
By setting $\alpha = O\big(\delta_n^{\frac{1}{\min\{\beta, 2\}+2}}\big)$ , we have $$
\|\hat{f} - f_0 \|_2^2 \leq \delta_n^{\frac{\min(\beta,2)}{\min(\beta,2) + 2}}.
$$
With Assumption \ref{ass:smooth-func}, we can provide another bound: $$
    \|\hat \hat{f}-f_0 \|^2_2\leq  \delta_n^{2/(3-\gamma)} / \alpha^{4 / (3-\gamma)} + \alpha^{\min(\beta,2)}
$$
In this case, set $\alpha = \delta_n^{\frac{2}{(3-\gamma)\min\{\beta, 2\}+4}}$, we have $$
\|\hat{f} - f_0\|_2^2 \leq \delta_n^{\frac{2\min\{\beta,2\}}{(3-\gamma)\min\{\beta, 2\}+4}}.
$$
Let 
\begin{align}\label{eq:defs}
    L(k) &:=  - 2 Y k(Z) +  \{k(Z)\}^2  \\
    f^{\star}_{\alpha} &:= \argmin_{f\in\Fcal} \EE[\{h^{\star}(Z)  - (\Tcal f)(Z) \}^2 +\alpha f^2(X) ]\nonumber. 
\end{align}
As in \citet{LiaoLuofeng2020PENE}, it is well-known that 
\begin{align*}
    \|f^{*}_{\alpha}-f_0\|^2_2 \leq \alpha^{\min(\beta,2)} \|w_0\|_2^2. 
\end{align*}
By the optimality of $f_\alpha^*$ in \eqref{eq:single-iter}, we have 
\begin{align*}
    & \alpha \|\hat{f} - f^{\star}_{\alpha} \|^2_2 + \|\Tcal(\hat{f}-f^{\star}_{\alpha})\|^2_2  \\
    &\qquad\qquad\leq \EE[L(\Tcal \hat{f}) ] - \EE[L (\Tcal f^{\star}_{\alpha})] + \alpha\{\EE[\hat{f}^2(Z)] - \EE[f^{\star}_{\alpha}(Z)^2]\},
\end{align*}
Recall that \begin{align*}
    &\EE[L(\Tcal \hat{f}) ] - \EE[L (\Tcal f^{\star}_{\alpha})] + \alpha\{\EE[\hat{f}^2(Z)] - \EE[f^{\star}_{\alpha}(Z)^2]\} =\\
    &\qquad\EE[ - 2 Y  \Tcal (f-f^{\star}_{\alpha})(Z)  + (\Tcal f)^2(Z) -  (\Tcal f^{\star}_{\alpha})^2(Z) ] +  \alpha\{\EE[f^2(Z)] - \EE[f^{\star}_{\alpha}(Z)^2]\},
\end{align*}


we have 
\begin{align}\label{eq:single-quadratic}
    & \alpha \|f - f^{\star}_{\alpha} \|^2_2 + \|\Tcal(f-f^{\star}_{\alpha})\|^2_2  \nonumber\\ 
    &\qquad= \EE[ - 2 Y  \Tcal (f-f^{\star}_{\alpha})(Z)  + (\Tcal f)^2(Z) -  (\Tcal f^{\star}_{\alpha})^2(Z) ] +  \alpha\{\EE[f^2(Z)] - \EE[f^{\star}_{\alpha}(Z)^2]\}\nonumber\\
    &\qquad= \EE[ - 2 Y  \hat \Tcal (f-f^{\star}_{\alpha})(Z) + (\hat \Tcal f)^2(Z) - (\hat \Tcal f^{\star}_{\alpha})^2(Z) ]   +  C_1 \times \EE[|(\hat \Tcal - \Tcal) (f-f^{\star}_{\alpha})(Z)|] \nonumber\\ 
    &\qquad\qquad+  \alpha\{\EE[f^2(Z)] - \EE[f^{\star}_{\alpha}(Z)^2]\} \nonumber\\
    &\qquad\leq \mathrm{Emp} + \mathrm{Loss} + C_1 \times \EE[|(\hat \Tcal - \Tcal) (f-f^{\star}_{\alpha})(Z)|]  \nonumber\\ 
    &\qquad= \mathrm{Emp} + \mathrm{Loss} + \|(\hat \Tcal-\Tcal) (  f - f^{\star}_{\alpha})\|_2
\end{align}
where 
\begin{align*}
    &\mathrm{Emp} = |(\EE_n-\EE)  [L(\hat \Tcal f)- L(\hat \Tcal f^{\star}_{\alpha})+\alpha(f^2(Z)-f^{\star}_{\alpha}(Z)^2)]|,  \\ 
    & \mathrm{Loss} =  \EE_n[ -2Y  \hat \Tcal (f-f^{\star}_{\alpha})(Z)  + (\hat \Tcal f)^2(Z) - (\hat \Tcal f^{\star}_{\alpha})^2(Z) + \alpha\{f^2(Z)- f^{\star}_{\alpha}(Z)^2 \}]. 
\end{align*} 
Here, using Lemma \ref{lem:local-concen}, the term $\mathrm{Emp} $ is upper-bounded as follows: 
\begin{align}\label{eq:single-emp}    
  \mathrm{Emp}  &\leq \delta_n\{ \alpha \|f-f^{\star}_{\alpha}\|_2+ \|\hat \Tcal(f-f^{\star}_{\alpha})\|_2+   \delta_n   \} \nonumber\\ 
&\leq \delta_n\{ \alpha \|f-f^{\star}_{\alpha}\|_2 +  \|\Tcal(f-f^{\star}_{\alpha})\|_2 + \|(\hat \Tcal-\Tcal)(f-f^{\star}_{\alpha})\|_2+   \delta_n   \}. 
\end{align} 
In the first inequality, we use the following variance bound: 
\begin{align*}
    &\EE[\{L(\hat \Tcal f)- L(\hat \Tcal f^{\star}_{\alpha})\}^2]\\
    &\leq \EE\left [ \left \{ -2 Y ( \hat \Tcal f - \hat \Tcal f^{\star}_{\alpha})(Z) + \{\hat \Tcal f-\hat \Tcal f^{\star}_{\alpha} \}(Z)\{ \hat \Tcal f + \hat \Tcal f^{\star}_{\alpha}\}(Z) \right \}^2 \right ] \\ 
    &\leq C_2 \times \|\hat \Tcal (  f - f^{\star}_{\alpha})\|^2_2. 
\end{align*}

Furthermore, recall that by our iteration in \eqref{eq:single-iter}, we have
\begin{align*}
     \EE_n[  - 2 Y  \Tcal (\hat f-f^{\star}_{\alpha})(Z)  + (\Tcal \hat f)^2(Z) - (\Tcal f^{\star}_{\alpha})^2(Z)+ \alpha\{\hat f^2(Z)- f^{\star}_{\alpha}(Z)^2 \} ]\leq 0. 
\end{align*}
Hence, we have 
\begin{align}\label{eq:single-loss}
    \mathrm{Loss}\leq 0. 
\end{align}

Combining everything, we have 
\begin{align*}
    &\alpha \|\hat f - f^{\star}_{\alpha} \|^2_2 + \|\Tcal(\hat f-f^{\star}_{\alpha})\|^2_2 \\
    & \leq \delta_n\{ \alpha \|\hat f-f^{\star}_{\alpha}\|_2 + \|\Tcal(f-f^{\star}_{\alpha})\|_2 +  \|(\hat \Tcal-\Tcal)(\hat f-f^{\star}_{\alpha})\|_2+   \delta_n   \}+  \|(\hat \Tcal-\Tcal)(\hat f-f^{\star}_{\alpha})\|_2 \\ 
     & \leq \delta_n\{ \alpha \|\hat f-f^{\star}_{\alpha}\|_{2} + \|\Tcal(f-f^{\star}_{\alpha})\|_2 +   \delta_n^{1/2} \|\hat{f} - f_\alpha^*\|_2+   \delta_n   \}+  \delta_n^{1/2} \|\hat{f} - f_\alpha^*\|_2\tag{From Lemma \eqref{lem:appro}} \\
     &\leq c_1 \delta^2_n  + c_2 \delta_n^{1/2} \|\hat{f} - f_\alpha^*\|_2
\end{align*}
\zihao{Is the rate too slow? Directly bounding $\|(\hat{\Tcal} - \Tcal)f
\|_2$ with our previous tricks give a bound of $O(\delta_n^{1/2})$ instead of $O(\delta_n)$, which might incur a problem here}
Here the constant $c_1$ and $c_2$ hide constants related to $C, c_0$. The first inequality comes from \eqref{eq:single-emp}. We implicitly use $\delta_n\leq 1, \alpha \leq 1$ in the last inequality. 
By Lemma \ref{lem:aux-1}, we have 
\begin{align*}
\|\hat f - f^{\star}_{\alpha} \|^2_2 \leq  O( (\delta_n / \alpha^2)  + \delta_n^2 / \alpha\big) = O(\delta_n / \alpha^2).
\end{align*}
\begin{remark}
    With the smoothness assumption i.e. Assumption \ref{ass:smooth-func}, we could have a bound $\|\hat f - f^{\star}_{\alpha} \|^2_2 \leq O(\delta_n^{1/2}/\alpha)^{2/(1.5 - \gamma/2)} = O(\delta_n^{2/(3-\gamma)} / \alpha^{4 / (3-\gamma)})$ simply by $$
    \alpha \|\hat f - f^{\star}_{\alpha} \|^2_2 + \|\Tcal(\hat f-f^{\star}_{\alpha})\|^2_2 \leq c_1 \delta_n^2 + c_2 \delta_n^{1/2}\|\hat{f} - f_\alpha^*\|_2^{(1+\gamma)/2}.
    $$
\end{remark}
\begin{lemma}\label{lem:aux-1}
    For $$
    x^2 \leq c_1 + c_2 x^\gamma,
    $$
    where $c_1, c_2>0$, $0\leq\gamma\leq 1$, we have $x \leq 2 \max\big\{\sqrt{c_1}, c_2^{1/(2-\gamma)}\big\}$. 
\end{lemma}
\begin{proof}
   Since $x^2 - c_2 x^\gamma -c_1$ is a convex function with negative intercept, we only need to prove that for $x_0 = 2 \max\{\sqrt{c_1}, c_2^{1/(2-\gamma)}\}$, we have $x_0^2 - c_2 x_0^\gamma -c_1 \geq 0$. For $\sqrt{c_1} \geq c_2^{1/(2-\gamma)}$, we have $$
   x_0^2 = 4c_1 \geq c_1 + 2^\gamma c_1 \geq c_1+ c_2 x_0^\gamma,
   $$
   similarly we have the same result when $\sqrt{c_1} \leq c_2^{1/(2-\gamma)}$,
   and we conclude the proof.
\end{proof}

\subsubsection{Under Mis-specification}


\subsection{Analysis for Iterative Version}
In the following, we iteratively use the notation \begin{align}\label{eq:iter-tikh}
    f_{m,*} = \argmin_{f\in\Fcal} \|\Tcal f - \Tcal f_0\|^2 + \alpha \|f - f_{m-1,*}\|^2.
\end{align}


The following lemma utilizes the source condition to  control the regularization bias caused by the
Tikhonov regularization term.  It is demonstrated that the regularization bias of the t-th population iterate, denoted as $f_{m,*} - f_0$, is significantly smaller than that of the non-iterated one, particularly when $\beta$ is large.
\begin{lemma}[Iterated Tikhonov Regularization Bias]\label{lem:iter-tikh-bias}
    If $f_0$ is the minimum $L_2$-norm solution to the linear inverse problem and satisfies the $\beta$-source condition, then the solution to the $t$-th iterate of Tikhonov
regularization $f_{m,*}$, defined in Equation \eqref{eq:iter-tikh}, with $f_{0,*}= 0$ , satisfies that
$$
\|f_{m,*} - f_0\|^2 \leq \|w_0\|^2 \alpha^{\min\{\beta, 2t\}}, \qquad\|\Tcal f_{m,*} - \Tcal f_0\|^2 \leq \|w_0\|^2 \alpha^{\min\{\beta+1, 2t\}}.
$$
\end{lemma}
Recall that our solution $\hat{f}_{m}$ satisfies $$
\hat{f}_m = \argmin_{f\in\Fcal} L(\hat{\Tcal}f) + \alpha \mathbb{E}_n [\{f - \hat{f}_{m-1}\}^2].
$$
We define \begin{align*}
L_m(\tau) &= \mathbb{E}[\mathbb{E}[f_0 - f_{m,*} - \tau (\hat{f}_m - f_{m,*})\mid Z]^2]+\alpha \|f_{m,*} + \tau (\hat{f}_m - f_{m,*}) - f_{m-1,*}\|^2,
\end{align*}
By definition, $L_m(\tau)$ is minimized by $\tau = 0$. Note that by strong convexity and property of quadratic function, we have $$
L_m(1) - L_m(0) = L'(0) + L''(0) \geq L''(0),
$$
Therefore \begin{align*}
    &\alpha \|\hat{f}_m - f_{m,*}\|^2 + \|\Tcal(\hat{f}_m - f_{m,*})\|^2\\
    & \leq\|\Tcal(f_0 - \hat{f}_m)\|^2 - \|\Tcal(f_0 - f_{m,*})\|^2 + \alpha\big(\|\hat{f}_m - f_{m-1,*}\|^2 - \|f_{m,*} - f_{m-1,*}\|^2\big)\\
    & =\mathbb{E}[L(\Tcal\hat{f}_m)] - \mathbb{E}[L(\Tcal f_{m,*})]+ \alpha\big(\|\hat{f}_m - f_{m-1,*}\|^2 - \|f_{m,*} - f_{m-1,*}\|^2\big),\\
    % &\qquad+ \alpha\big(\|\hat{f}_m - f_{m-1,*}\|^2 - \|f_{m,*} - f_{m-1,*}\|^2\big).
    % &\leq\mathbb{E}[-2Y\cdot\Tcal(\hat{f}_m - f_{m,*}) + (\Tcal\hat{f}_m)^2 - (\Tcal f_{m,*})^2] + \alpha\big(\|\hat{f}_m - f_{m-1,*}\|^2 - \|f_{m,*} - f_{m-1,*}\|^2\big)\\
    % &\leq \mathbb{E}[-2Y\cdot\hat{\Tcal}(\hat{f}_m - f_{m,*}) + (\hat{\Tcal}\hat{f}_m)^2 - (\hat{\Tcal} f_{m,*})^2] + \alpha\big(\|\hat{f}_m - f_{m-1,*}\|^2 - \|f_{m,*} - f_{m-1,*}\|^2\big) \\
    % &\qquad + c\cdot\mathbb{E}[|(\Tcal - \hat{\Tcal})(\hat{f}_m - f_{m,*})|]+ \alpha\big(\|\hat{f}_m - f_{m-1,*}\|^2 - \|f_{m,*} - f_{m-1,*}\|^2\big)\\
    % & = \mathbb{E}[L(\hat{\Tcal}\hat{f}_m)] - \mathbb{E}[L(\hat{\Tcal} f_{m,*})] + c\cdot\mathbb{E}[|(\Tcal - \hat{\Tcal})(\hat{f}_m - f_{m,*})|]\\
    % &\qquad+ \alpha\big(\|\hat{f}_m - f_{m-1,*}\|^2 - \|f_{m,*} - f_{m-1,*}\|^2\big).
\end{align*}
and thus we have \begin{align*}
&\alpha \|\hat{f}_m - f_{m,*}\|^2 + \|\Tcal(\hat{f}_m - f_{m,*})\|^2\\
    & \leq \mathbb{E}[L(\hat{\Tcal}\hat{f}_m)] - \mathbb{E}[L(\hat{\Tcal} f_{m,*})] + c\cdot\mathbb{E}[|(\Tcal - \hat{\Tcal})(\hat{f}_m - f_{m,*})|]\\
    &\qquad+ \alpha\big(\|\hat{f}_m - f_{m-1,*}\|^2 - \|f_{m,*} - f_{m-1,*}\|^2\big)\\
    & \leq |(\mathbb{E} - \mathbb{E}_n)(L(\hat{\Tcal}\hat{f}_m) - L(\hat{\Tcal} f_{m,*}) )| + \EE_n[L(\hat{\Tcal}\hat{f}_m) - L(\hat{\Tcal} f_{m,*})] \\
    &\qquad+ c\cdot\mathbb{E}[|(\Tcal - \hat{\Tcal})(\hat{f}_m - f_{m,*})|] +  \alpha\big(\|\hat{f}_m - f_{m-1,*}\|^2 - \|f_{m,*} - f_{m-1,*}\|^2\big)\\
    &\leq c_1 (\delta_n\|\hat{\Tcal}(\hat{f}_m - f_{m,*})\| + \delta_n^2)+ \EE_n[L(\hat{\Tcal}\hat{f}_m) - L(\hat{\Tcal} f_{m,*})] + c_2\cdot\delta_n \|\hat{f}_m - f_{m,*}\|_\infty \\
    &\qquad +  \alpha\big(\|\hat{f}_m - f_{m-1,*}\|^2 - \|f_{m,*} - f_{m-1,*}\|^2\big),\\
\end{align*}
Here the second inequality comes from triangular inequality, the third inequality comes from Lemma \ref{lem:appro} and Lemma \ref{lem:local-concen}.
Recall that by \eqref{eq:main-iter},  $$
 \EE_n[L(\hat{\Tcal}\hat{f}_m) - L(\hat{\Tcal} f_{m,*})] \leq  \alpha(\|f_{m,*} - \hat{f}_{m-1}\|^2_n - \|\hat{f}_m - \hat{f}_{m-1}\|^2_n),
$$
we have \begin{align*}
&\alpha \|\hat{f}_m - f_{m,*}\|^2 + \|\Tcal(\hat{f}_m - f_{m,*})\|^2\\
&\leq c_1 (\delta_n\|(\hat{f}_m - f_{m,*})\| + \delta_n^2) + \alpha(\|f_{m,*} - \hat{f}_{m-1}\|^2_n - \|\hat{f}_m - \hat{f}_{m-1}\|^2_n) + c_2\cdot\delta_n \|\hat{f}_m - f_{m,*}\|_\infty \\
&\qquad +  \alpha\big(\|\hat{f}_m - f_{m-1,*}\|^2 - \|f_{m,*} - f_{m-1,*}\|^2\big).
\end{align*}
By \cite{bennett2023source} we can prove that $$
(\|f_{m,*} - f_{m-1,*}\|^2_n - \|\hat{f}_m - f_{m-1,*}\|^2_n) +\big(\|\hat{f}_m - f_{m-1,*}\|^2 - \|f_{m,*} - f_{m-1,*}\|^2\big)
$$
is bounded by $$O(\delta_n^2 + \delta_n(\|\hat{f}_m - f_{m,*}\| + \|\hat{f}_{m-1} - f_{m-1,*}\|)) + 2\|\hat{f}_{m-1} - f_{m-1,*}\|\|\hat{f}_{m} - f_{m,*}\|.$$ 
Therefore, we have \begin{align*}
    &\alpha \|\hat{f}_m - f_{m,*}\|^2 \\
    &\qquad\leq O\bigg(\delta_n^2 + \delta_n\|\hat{f}_m - f_{m,*}\| + \delta_n\|\hat{f}_m - f_{m,*}\|_\infty + \alpha\delta_n(\|\hat{f}_m - f_{m,*}\| + \|\hat{f}_{m-1} - f_{m-1,*}\|))\bigg)\\
    &\qquad\qquad + 2\alpha\|\hat{f}_{m-1} - f_{m-1,*}\|\|\hat{f}_{m} - f_{m,*}\|.
\end{align*}
By applying AM-GM inequality and utilizing $\alpha\leq 1$, we have $$
\frac{\alpha}{8}\|\hat{f}_m - f_{m,*}\|^2 \leq O\big(\delta_n^2/\alpha + \delta_n^2 
+\alpha\delta_n^2\big) + c_2\cdot\delta_n \|\hat{f}_m - f_{m,*}\|_\infty + 2\alpha\|\hat{f}_{m-1} - f_{m-1,*}\|^2,
\label{eqn:strong-iterated}$$
therefore we have $$
\|\hat{f}_m - f_{m,*}\|^2 \leq O\big(\delta_n^2/\alpha^2 + \delta_n^2 /\alpha
\big) + O\bigg(\frac{\delta_n}{\alpha} \|\hat{f}_m - f_{m,*}\|_\infty\bigg)+ 16\|\hat{f}_{m-1} - f_{m-1,*}\|^2.
$$
By Assumption \ref{ass:smooth-func}, we have $\|\hat{f}_m - f_{m,*}\|_\infty \leq\|\hat{f}_m - f_{m,*}\|_2^\gamma $, which implies $$
\|\hat{f}_m - f_{m,*}\|^2 \leq O\big(\delta_n^2/\alpha^2 + \delta_n^2 /\alpha
\big) + O\bigg({\delta_n}/{\alpha} \cdot\|\hat{f}_m - f_{m,*}\|^\gamma\bigg)+ 16\|\hat{f}_{m-1} - f_{m-1,*}\|^2,
$$
by Lemma \ref{lem:aux-1}, we have \begin{align*}
    \|\hat{f}_m - f_{m,*}\|^2 &\leq 4\max\big\{O\big(\delta_n^2/\alpha^2 + 16\|\hat{f}_{m-1} - f_{m-1,*}\|^2 
\big), O\big((\delta_n/\alpha)^{2/(2-\gamma)}\big)\big\}\\
&\leq O(128^m \max\big\{\delta_n^2/\alpha^2, (\delta_n/\alpha)^{2/(2-\gamma)}\big\}),
\end{align*}
where the second inequality comes from induction.
Combine this result and Lemma \ref{lem:iter-tikh-bias}, we are now ready to prove that \begin{align}
    \|\hat{f}_m - f_0\|^2 \lesssim 128^m \max\big\{\delta_n^2/\alpha^2, (\delta_n/\alpha)^{2/(2-\gamma)}\big\}) + \|w_0\|^2 \alpha^{\min\{\beta, 2m\}}.
\end{align}
\section{MLE Analysis with Model Misspecification}
In this section, we analyze the case where the realizability assumption does not hold, i.e. $g^*(x\mid z)$ does not necessarily lie in $\Gcal$. 
\paragraph{Route 1. Similar to \cite{Zhang_2006,agarwal2020flambe}.} First, we define the misspecification error $\epsilon$ in $L_\infty$ norm by $\misp$, i.e. $\misp = \min_{g\in\Gcal} \|g - g^*\|_\infty$, and the  corresponding projection of $g^*$ on $\Gcal$ to be $\gproj$. i.e. $\gproj = \argmin_{g\in \Gcal} \|\log g - \log g^*\|_\infty$. Here the $L_\infty$ norm is taken w.r.t. $\Xcal \times \Zcal$. By MLE procedure, we have \begin{align}\label{ineq:mle-lower}
    \frac{1}{n}\sum_{i=1}^{n}\log\bigg(\frac{\hat{g}(x_i\mid z_i)}{g^*(x_i \mid z_i)}\bigg)
    &= \frac{1}{n}\sum_{i=1}^{n}\log\bigg(\frac{\hat{g}(x_i\mid z_i)}{\gproj(x_i\mid z_i)}\bigg)+\frac{1}{n}\sum_{i=1}^{n}\log\bigg(\frac{\gproj(x_i\mid z_i)}{g^*(x_i\mid z_i)}\bigg)\nonumber\\
    &\geq \frac{1}{n}\sum_{i=1}^{n} \log(g^*(x_i\mid z_i)) - \frac{1}{n}\sum_{i=1}^{n} \log(\gproj(x_i\mid z_i))\nonumber\\
    &\geq - \misp,
\end{align}
Here the first inequality comes from $\hat{g} = \argmax_{g\in \Gcal} \sum_{i=1}^n\log (g(x_i\mid z_i))$ and $\gproj \in \Gcal$.
Meanwhile, by Markov's inequality and Boole’s inequality, it holds with probability at least $1-\delta$ that for all $\bar{g}\in\Ncal_{[]}(\Gcal,\|\cdot\|_1,1/n)$, we have \begin{align*}
    \sum_{i=1}^{n} \frac{1}{2}\log\bigg(\frac{\bar{g}(x_i\mid z_i)}{g^*(x_i \mid z_i)}\bigg) &\leq n\log\bigg(\EE_{z\sim g^*(z), x\sim g^*(x\mid z)}\bigg[\exp\bigg(\frac{1}{2}\log\bigg(\frac{\bar{g}(x\mid z)}{g^*(x\mid z)}\bigg) \bigg)\bigg]\bigg)\\
    &\qquad+\log\bigg(\frac{N_{[]}(\Gcal,\|\cdot\|_1, 1/n)}{\delta}\bigg),
\end{align*}
specify $\bar{g}$ to be the upper bracket of $\hat{g}$, by \eqref{ineq:mle-lower}, we have \begin{align*}
    -n\cdot\misp &\leq n\log\bigg(\EE_{z\sim g^*(z), x\sim g^*(x\mid z)}\bigg[\exp\bigg(\frac{1}{2}\log\bigg(\frac{\bar{g}(x\mid z)}{g^*(x\mid z)}\bigg) \bigg)\bigg]\bigg)\\
    &\qquad+\log\bigg(\frac{N_{[]}(\Gcal,\|\cdot\|_1, 1/n)}{\delta}\bigg)\\
    &\leq n\cdot \log\bigg(\EE_{z\sim g^*(z), x\sim g^*(x\mid z)}\bigg[\sqrt{\frac{\bar{g}(x\mid z)}{g^*(x\mid z)}}\bigg]\bigg)+\log\bigg(\frac{N_{[]}(\Gcal,\|\cdot\|_1, 1/n)}{\delta}\bigg)\\
    &= n\cdot\log\bigg(\EE_{z\sim g^*(z)}\bigg[\int_{\Xcal} \sqrt{\bar{g}(x\mid z)\cdot g^*(x\mid z)}dx\bigg] \bigg) + \log\bigg(\frac{N_{[]}(\Gcal,\|\cdot\|_1,1/n)}{\delta}\bigg),
\end{align*}
 Utilizing the $\log x \leq x-1$, we have $$
1- \EE_{z\sim g^*(z)}\bigg[\int_{\Xcal} \sqrt{\bar{g}(x\mid z)\cdot g^*(x\mid z)}dx\bigg] \leq \frac{1}{n} \log\bigg(\frac{N_{[]}(\Gcal,\|\cdot\|_1,1/n)}{\delta}\bigg)+{\misp}.
$$
Therefore we can bound the Hellinger distance between $\bar{g}(x\mid z)$ and $g^*(x\mid z)$,
\begin{align}\label{hell-dist-1}
    \EE_{z\sim g^*(z)}[H^2(\hat{g}(\cdot\mid z), g^*(\cdot \mid z))] &= \EE_{z\sim g^*(z)}\bigg[\int_{\Xcal}\big(\hat{g}(x\mid z)^{1/2} - g^*(x\mid z)^{1/2} \big)^2dx\bigg]\nonumber\\
    &\leq 2\EE_{z\sim g^*(z)}\bigg[1-\int_{\Xcal}\sqrt{\bar{g}(x\mid z)\cdot g^*(x\mid z)}\bigg] + \frac{\sqrt{C/c_0}}{n}\nonumber\\
    &\leq \frac{2}{n}\log\bigg(\frac{N_{[]}(\Gcal,\|\cdot\|_1,1/n)}{\delta}\bigg)+\frac{1}{n}+{2\misp},
\end{align}
here the first inequality comes from the fact that $\bar{g}$ is a $1/n$-upper bracketing of $\hat{g}$, and \begin{align*}
    \int_\Xcal \big|\sqrt{\bar{g}(x|z)\cdot g^*(x|z)} - \sqrt{\hat{g}(x|z)\cdot g^*(x|z)} \big|dx &= \int_{\Xcal} (\sqrt{\bar{g}(x|z)} - \sqrt{\hat{g}(x|z)})\sqrt{g^*(x|z)}dx\\
    &\leq \sqrt{\int_\Xcal g^*(x\mid z)dz}\cdot \sqrt{\int_\Xcal (\sqrt{\bar{g}(x|z)} - \sqrt{\hat{g}(x|z)})^2 dx}\\
    &\leq 1/n.
\end{align*}. 
\paragraph{Route 2. Critical Radius approach }
We aim to generalize the techniques in Chapt. 14 in \cite{wainwright2019high}. First, we define the population version of localized Rademacher complexity for function class $\Hcal^* := \operatorname{star}((\Hcal - h^*) \cup \{0\})$, where $\Hcal$ is bounded by a constant $b$ in $\|\cdot\|_\infty$, and will be specified later:
$$
\bar{R}_n (\delta; \Hcal^*) := \E_{z_i\sim g^*(z), x_i\sim g^*(x\mid z_i)} \E_\epsilon\bigg[\sup_{h\in \Hcal^*, \|h\|_2 \leq \delta} \bigg|\frac{1}{n}\sum_{i=1}^n \epsilon_i h(x_i, z_i)\bigg|\bigg],
$$
here $\epsilon_i$ are i.i.d. Rademacher random variables. The critical radius $\delta_n$ of function class $\Hcal^*$ is any solution such that \begin{align*}
    \delta^2 \geq c/n \text{ and }\bar{R}_n (\delta; \Hcal^*)\leq \delta^2/b.
\end{align*}
Such critical radius can be easily calculated for a large number of function classes. For example, we can use $$
\frac{64}{\sqrt{n}}\int_{\delta^2/2b}^{\delta}\sqrt{\log N_n(t, \Bcal(\delta,\Hcal^*))}dt \leq \frac{\delta^2}{b}
$$
to calculate $\delta_n$, where $\Bcal(\delta,\Hcal^*) := \{h\in\Hcal^*\mid \|h\|_2\leq \delta\}$, $N_n$ is the empirical covering number conditioned on $\{(x_i,z_i)\}_{i\in[n]}$.
For a cost function $\Lcal : \RR \rightarrow \RR$, we define $\Lcal_h(x,z):= \Lcal(h(x,z))$.
Next, we introduce the following lemma: 
\begin{lemma}[Theorem 14.20 in \cite{wainwright2019high}]\label{lem:2cond}
    With probability at least $1 - c_1 \exp(-c_2 n\delta_n^2/b)$, either of the following events hold: \begin{itemize}
        \item[(1)] $\|h - h^*\|_2\leq \delta_n$;
        \item[(2)] $|\PP_n (\Lcal_h - \Lcal_{h^*} ) -\PP (\Lcal_h - \Lcal_{h^*} ) | \leq 10L\delta_n\| h - h^*\|_2$.
    \end{itemize}
\end{lemma}
If  (1) in Lemma \ref{lem:2cond} holds, then our desired result is achieved. Therefore we condition on (2). Consider the case that $\PP_n(\Lcal_h - \Lcal_{h^*}) \leq \epsilon$, then we have \begin{align*}
    |\PP (\Lcal_h - \Lcal_{h^*} )| \leq 10L\delta_n\| h - h^*\|_2 + \epsilon.
\end{align*}
We make the following definition. \begin{definition}\label{def:strong-convexity}
        We say $\Lcal_h$ is $\gamma$-strongly convexity at $h^*$ if $$
\E_{z\sim g^*(z), x\sim g^*(x\mid z)}\big[\Lcal_h(x,z) - \Lcal_{h^*}(x,z) - \partial\Lcal_{h^*}(x,z) (h - h^*)(x,z) \big] \geq \frac{\gamma}{2} \|h - h^*\|^2_2
$$ for all $h\in\Hcal$.
\end{definition} 
When $\Lcal_h$ is $\gamma$-strongly convex at $h^*$, we have $$
\frac{\gamma}{2} \| h - h^*\|_2^2 \leq 10L \delta_n \|h - h^*\|_2 + \epsilon,
$$
solving this inequality and we got \begin{align}\label{eq:diff}
\| h - h^* \|_2 \leq \frac{20L\delta_n}{\gamma} + \sqrt{\frac{2 \epsilon}{\gamma}}.
\end{align}
Now we specify $\Hcal = \bigg\{ \sqrt{\frac{g+ g^*}{2g^*}}\biggl\vert g\in \Gcal \bigg\}$, and $\Lcal_h = -\log h(x)$ for $h\in \Hcal$. For any $h\in\Hcal$ we have $$
\|h\|_\infty = \bigg\|\sqrt{\frac{f + f^*}{ 2f^*}}\bigg\|_\infty \leq \sqrt{\frac{1}{2}\bigg(1+ \frac{C}{c_0}\bigg)},
$$
and $|\log h (x) - \log h'(x)|\leq \sqrt{2}|h(x) - h'(x)|$ since $\|h\|_\infty \geq 1/\sqrt{2}$. Note that $$
\| h - h^*\|_2^2 = \E_{z\sim g(z)}\bigg[H^2\bigg(\frac{g+ g^*}{2}\mid g^*\bigg)\bigg],
$$
and since $H^2(f_1\mid f_2) \leq 2 \operatorname{KL}(f_1\mid f_2)$, we have $\|h- h^*\|_2^2 \leq \PP(\Lcal_h - \Lcal_{h^*})$, thus $\Lcal$ is $2$-strongly convex at $h^*$. Note that \begin{align*}
\PP_n \Lcal_h - \PP_n \Lcal_{h^*} &= \frac{1}{2n} \sum_{i=1}^n\log\bigg(\frac{2g^*(x_i|z_i)}{(\hat{g}+g^*)(x_i|z_i)}\bigg)\\
&\leq \frac{1}{2n}\sum_{i=1}^n\log\bigg(\frac{g^*(x_i|z_i)}{\hat{g}(x_i|z_i)}\bigg).
\end{align*}
Since $\hat{g} = \argmax_{g\in\Gcal} \log(g(x_i|z_i))$, we have  $$
\PP_n \Lcal_h - \PP_n \Lcal_{h^*} \leq \min_{g\in \Gcal}\frac{1}{2n}\sum_{i=1}^n\log\bigg(\frac{g^*(x_i|z_i)}{{g}(x_i|z_i)}\bigg)\leq \min_{g\in\Gcal}\big\|\log(x|z) - \log g^*(x|z)\big\|_\infty.
$$
By applying Lemma \ref{lem:2cond} and specify $h = \sqrt{\frac{ \hat{g}+ g^*}{2g^*}}$, we have the following holds for any $\delta_n$ that satisfies the critical inequality and $\delta_n^2\geq (1+\frac{C}{c_0})\frac{1}{n}$: $$
\E_{z\sim g(z)}\bigg[H^2\bigg(\frac{\hat{g}+ g^*}{2}\mid g^*\bigg)\bigg] \leq \delta_n^2 + \misp
$$
with probability at least $1 - c_1\exp(c_2 \frac{c_0}{C+c_0}n\delta_n^2)$. Here $\misp := \min_{g\in\Gcal}\|\log(g(x|z)/g^*(x|z))\|_\infty$.
Using Lemma 4.1 in \cite{van1993hellinger}, we can prove that $$
\E_{z\sim g(z)}\bigg[H^2(\hat{g}(\cdot|z)\mid g^*(\cdot|z))\bigg] \leq \delta_n^2 + \misp
$$
\paragraph{Route 3. Bracketing Entropy \cite{bilodeau2023minimax}}
\section{Some Ideas for Improving MLE Rate}
\paragraph{Density estimation through projection.}Let's use the following loss function:
\begin{align*}
    \argmin_{g \in \Gcal}\EE_n[0.5\int g^2(x|Z)d\mu(x)]-\EE_n[g(x|z)].  
\end{align*}
where $\mu$ is the Lebsgue measure. Note each $g$ does not need to be normalized in this case. 
As in the Wainright book (Chapter 14), we can easily prove $$\EE_{z \sim g^{\star}(z)}\left[\{\int |\hat g(x|z)-g^{\star}(x|z)| \mathrm{d}\mu(x)\}^2\right]=O(\delta_n^2+\inf_{g\in \Gcal}\EE_{z \sim g^{\star}(z)}\left[\{\int | g(x|z)-g^{\star}(x|z)| \mathrm{d}\mu(x)\}^2\right])$$
with probability at least $1 - c_1 \exp(c_2n\delta_n^2)$. Consider the case of $\misp = 0$ then easily we have \begin{align*}
    \|(\Tcal - \hat{\Tcal}) f\|_2^2 &= \E_{z\sim g^*(z)}\bigg[\bigg(\int_{\Xcal} \{\hat{g}(x|z) - g^*(x| z)\}f(x)d\mu(x)\bigg)^2\bigg]\\
    &\leq \delta_n^2 \|f\|_\infty^2 .
\end{align*}
Therefore in our analysis (again, take the non-iterative case as an example), we have $$
\alpha \| \hat{f} - f^*_\alpha\|_2^2 \leq c_1 \delta_n^2 + \delta_n \| \hat{f} - f^*_\alpha\|_2 + \delta_n \|\hat{f} - f_\alpha^*\|_\infty.
$$
With Assumption \ref{ass:smooth-func}, we can still prove that $$
\|\hat{f} - f_\alpha^*\|_2^2 \leq (\delta_n/\alpha)^{\frac{2}{2-\gamma}}.
$$
Therefore \begin{align*}
    \|\hat f-f_0 \|^2_2\leq  (\delta_n/\alpha)^{\frac{2}{2-\gamma}}+ \alpha^{\min(\beta,2)}.
\end{align*}
By selecting $\alpha = O(\delta_n^{\frac{2}{2+ (2 - \gamma)\min\{\beta,2\}}})$, we have $$
\|\hat f-f_0 \|^2_2\leq \delta_n^{\frac{2\min\{\beta,2\}}{2+ (2 - \gamma)\min\{\beta,2\}}}.
$$
However, note that such a method is computationally inefficient. 
\section{Empirical Results}
In this paper, we adapt the experiment in \cite{kallus2022causal}, i.e. the double IV approach with negative control variables. We have the following conditional moment equality for variables $W,A,S,Q$:$$
\E[Y - h(W,A,S)|Q,A,S] = 0,
$$
where $h$ is the bridge function, $A$ is the treatment, $W$ is negative control outcome, $Q$ is negative control action, and $S$ is the covariate. We set $(Q,A,S)$ to be the instrumental variable, i.e. variable $Z$ in the theoretical part of this paper, and $(W,A,S)$ as the variable $X$ in the main paper. Our first step is to estimate the conditional density $p(w,a,s|q,a,s)$, and the second step is the iterative Tikhnov regression for $h$. Our target is to estimate $\E[Y(1)]$. In this paper, dimension of W = 1, dimension of S and Q = 15.
% \zihao{Note that $x= (w,a,s)$ and $z = (q,a,s)$, therefore a number of $\sigma_k(z)$ could converge to $0$ due to the same components in $X$ and $Z$. This could cause numerical instability when approximating with GMM-- the NN frequently output \texttt{Nan}. Shall we directly learn $p(w|q,a,s)$? }
\subsection{Density Estimation}
\paragraph{Function Class.} In this paragraph, we report the method of MLE. Our current architecture adopts the following conditional Gaussian Mixture Model: $$
p(x|z) = \sum_{k = 1}^K  \frac{1}{(2\pi)^{d/2}\sigma_k^d}\cdot \pi_k(z)\exp(-\|x - \mu_k(z)\|_2^2 / 2\sigma_k^2(z)),
$$
and we approximate both $\sigma_k(z), \pi_k(z)$ and $\mu_k(z)$ with a three-layer neural network. Since $X$ is continuous, we set NLL (negative log-likelihood) to be our loss function. Note that NLL not necessarily converges to 0 even when the sample number goes to infinity, which could add difficulty to the tuning process, therefore we only compare the results under different method, including MLE, density estimation through projection, and the linear version MLE in package \texttt{EconML} . We report the curve of NLL on the validation set. 
% \paragraph{Comparison between different Component Number $K$}
% First, we fix the sample size = 1600, and plot the testing loss curve of different $K$ with respect to the number of epochs in training. We train everything for 200 epochs to ensure the training curve converges. Due to the training instability we described, we use Adam to be the optimizer and lr = 1e-6 when $K>=40$, and lr = 1e-5 when $K<40$.

% \begin{figure}[h]  
%   \centering
%   \includegraphics[width=0.6\textwidth]{Figures/K_curves.png}  
%   \caption{The comparison between different component number $K$. Here we present the curves of NLL on the validation set v.s. the number of epochs in training. The training and validation set are split by a proportion of 8:2.}
%   \label{fig:example}
% \end{figure}

\paragraph{Regression Setup.}
In the second-stage regression, we use the following loss function: $$
\E_n[Y - \hat{\E}[f(X)|Z]]^2 + \alpha \E_n[(f - \hat{f}_{m-1})(X)^2].
$$
We use a 4-layer neural network as the function approximator of $f$. To estimate the conditional expectation of $X$ given $Z$, we sample $n^2$ samples from the pre-trained conditional GMM for each $X$, and take the mean of $f(X_{sampled})$ as the approximation of $\hat{\E}[f(X)|Z]$.
\paragraph{Density Estimation through Projection.} Since we have GMM as the function approximator for the ground truth density model, our loss function has the following form: \begin{align*}
    &0.5 \EE_n\bigg[\sum_{j,k = 1}^K \frac{\pi_k(Z) \pi_j(Z)}{(2\pi)^{d/2}}\cdot\frac{1}{(\sigma_j^2(Z)+\sigma_k^2(Z))^{d/2}}\exp\bigg\{-\frac{\|\mu_j(Z) - \mu_k(Z)\|_2^2}{2(\sigma_k^2(Z) + \sigma_j^2(Z))}\bigg\}\bigg]\\
    &\qquad\qquad - \E_n \bigg[\sum_{k=1}^K\frac{\pi_k(Z)}{(2\pi)^{d/2}\sigma_k^d(Z)}  \exp(- \|X - \mu_k(Z)\|_2^2 / 2\sigma_k^2(Z))\bigg]
\end{align*}
\zihao{Empirically not working due to instability. I think it's due to numerical issues such as truncation: the densities are small when the component numbers are big. Debugging. }
\begin{table}[!htbp]
    \begin{center}
    \caption{Normalized MSE for $\EE[Y(1)]$: $n_1$ = 1000 for the training set, and $n_2$ = 250 for the validation set. We used 10 random trials to calculate the results. The confidence interval is calculated by 2 times the standard deviation. }
    \begin{tabular}{c|lllll}
        \toprule
                & MLE-IV$(\alpha = 0.1)$ & KernelIV & DeepIV &DFIV& AGMM\\
        \midrule
        $t$ &0.019 \pm 0.008	&0.030\pm0.006	&0.095\pm 0.015 &0.160\pm 0.043 & 0.056 \pm 0.003 \\
        $t^3+10t$&0.003\pm 0.004& 	0.023\pm0.004 &0.008\pm 0.002 &0.181\pm0.012 &\\
        sigmoid$(t)$ &0.013\pm 0.004 &0.078\pm 0.008 &0.049\pm0.006 &0.198\pm 0.011 & \\
        $\exp(t) $ &0.006\pm 0.004 &0.020 \pm 0.007 &0.016\pm 0.021 &0.191\pm0.084 & \\
        \bottomrule
    \end{tabular}\label{tb:Fitness_formations}
    \end{center}
\end{table}
\begin{figure}[!htbp]  
  \centering
  \includegraphics[width=0.8\textwidth]{Figures/L_curve_t.jpg}  
  \caption{L-curve of MLE-IV. The y-axis represents $\|\hat{f}(X)\|_n$ and the x-axis represents $\|\E[Y - \hat{f}(X)|Z]\|_n$. The alpha list = {0, 1e-5,  1e-4, 1e-3, 1e-2, 0.1,0.2, 0.4, 1}. A very slight slope change at $\alpha = 0.1$.}
  \label{fig:example}
\end{figure}


\zihao{TODO:
\begin{itemize}
    \item Tune the baselines; add experiments for iterative case; (try out the projection density estimation)
    \item Experiments for model selection part
\end{itemize}
}
In MLE- IV and DeepIV experiment, we choose a 4-layer neural network to be the function approximator for both the Gaussian mixture model and bridge function $f$. We use Adam as the optimizer for both networks. The learning rate =0.001 for both networks, and we set weight decay = 0.0001 for the Gaussian mixture model, and weight decay = 0.001 for the bridge function approximator.
In the KIV experiment, we adopt the parameter-tuning method in \cite{singh2019kernel}. In DFIV experiments, we use the original training parameters in the experiment of \cite{xu2020learning}. The dimensions of both features for covariate and instrumental variables are set to 32. For the network architecture, we adapt a 4-layer neural network which has the same hidden layer structure as MLEIV network, except the output dimension of both networks in DFIV are set to be 10 for future linear regression step. However, the results don't seem to be satisfying-- the network seems to be predicting $\E[Y|A = 0]$, instead of a meaningful outcome.

\section{Model Selection}
In this section, we aim to prove excess risk guarantees with respect to the best model selection by constructing a convex combination of the candidate models. It is natural to adopt a two-step approach: first, we select the best convex combination for the likelihood model $\hat{g}(X|Z)$, the candidate models $f_j(X)$ are subsequently trained with this likelihood function over different function classes, and finally, the best convex combination of the candidate models using $\hat{g}(X|Z)$ as the nuisance function.

First let us focus on the model selection process for $f(X)$, given $\hat{g}(X|Z)$. 
Define the loss function:
\begin{align}
    \ell_{f,g}(Y,Z,X) = \underbrace{\left(Y - \int f(x)g(x|Z)dx\right)^2}_{\ell^{0}_{f,g}} + \underbrace{\alpha f(X)^2}_{r_f} \label{Eqn:q_agg_loss}
\end{align}
The goal is to find the a convex combination of the $M$ candidate models $\{f_1, \dots, f_M\}$ such that the loss achieves a $\frac{\log(M)}{n}$ rate with respect to the model selection. For $\theta\in\Theta = \{\theta | \sum_j \theta_j = 1, \theta_j\geq 0 \forall j\} $, denote $f_{\theta} = \sum_j \theta_j f_j$.For any convex combination $\theta$ over a set of candidate functions $\{f_1, \ldots, f_M\}$, we define the notation:
\begin{align}
    \ell_{\theta,g}(Y,Z,X) :=~& \ell_{f_{\theta}, g}(Y,Z,X) &
    R(\theta,g) :=~& P\ell_{\theta,g}(Y,Z,X)
\end{align}

\subsection{Model Selection Overview}
First we define some optimal aggregates in the following sense:
\begin{align*}
    j^*_{\alpha} :=~& \argmin_{j = 1, \dots, M}  R(h_{j},g_0) &
    j^* :=~& \argmin_{j = 1, \dots, M} \|h_0 - h_j\|^2\\
    \theta^*_{\alpha} :=~& \argmin_{\theta \in \Theta} R(h_{\theta},g_0) &
    \theta^* :=~& \argmin_{\theta \in \Theta} \|h_0 - h_{\theta}\|^2\\
    h_{\alpha}^* :=~& \argmin R(h,g_0) &
    h_{\alpha, \Hcal}^* :=~& \argmin_{h \in \Hcal} R(h,g_0)
\end{align*}
\begin{theorem} [Oracle Inequalities for Excess Risk] \label{thm:oracle_model_selection}
    Assume that $Y$ is almost surely bounded by $b$, each candidate model $h_j$ is almost uniformly bounded in $[-b,b]$ almost surely. Moreover, assume that $g^*$ and all $g\in \mathcal{G}$ is bounded away from 0, i.e. exist $c_0>0$ such that $g^*, g\geq c_0$. Let $\error(\hat{g}) = 2(\frac{1}{c_0}+1)^2\E[H(\hat{g},g_0)^2]^\frac{1}{2}$. Then, using the loss function described above, for any $\delta>0$, with probability $>1-\delta$:\\
    The output of \textbf{Best-ERM}, $\hat{\theta}$, satisfies:
    \begin{align*}
        R(h_{\hat{\theta}},g_0)~& -R(h^*,g_0)\\
        ~&\leq O\left(\frac{\log(M/\delta)}{n} + \error(\hat{g})^2\right)\\
        ~& +  \frac{1}{\alpha}\left(\min_{j=1}^M R(h_j, g_0) - R(h^*,g_0) \right)
    \end{align*} for any $h^*$ that satisfies the first order condition for the regularized loss, for instance $h_{\alpha}^*$ and $h_{\theta_{\alpha}^*}$.\\ 
    The output of \textbf{Convex-ERM}, $\hat{\theta}$, satisfies:
    \begin{align*}
        R(h_{\hat{\theta}},g_0) - ~&\min_{\theta\in\Theta}R(h_\theta,g_0)\\
        ~&\leq O\left(\frac{M+\log(1/\delta)}{n} + \error(\hat{g})^2\right)
    \end{align*}
     Furthermore, for any $\mu \leq \min\{\frac{\alpha(1-\nu)}{7}, \frac{1-\nu}{3}, \frac{\alpha\nu}{5}\}$, the output of \textbf{Q-aggregation}, $\hat{\theta}$, satisfies: 
    \begin{align*}
        R(f_{\hat{\theta}}, g_0) -~&\min_{j=1}^M R(f_j, g_0)\\
        ~&\leq O\left(\frac{\log(M/\delta)}{n} + \frac{1}{\mu}\error(\hat{g})\right)\\
    \end{align*}
\end{theorem}

Note: the $\error(\hat{g})$ term will be on the order of $\delta_n$. Since there will be lower order terms in the other parts of the proof, having $\error(\hat{g})$ is not as problematic.

Now we want to convert the excess risk bound to a strong metric bound on the MSE of the learned aggregate. For instance, consider the Q-aggregation aggregate:
\begin{align*}
    \|h_{\hat{\theta}} - h_0\|^2 \leq ~& 2\|h_{\hat{\theta}} - h_{\alpha}^*\|^2 + 2\|h_{\alpha}^*-h_0\|^2\\
    \leq~& \frac{2}{\alpha}\left(R(h_{\hat{\theta}},g_0) - R(h_{\alpha}^*,g_0)\right) + O\left(\alpha^{\min\{2,\beta\}}\right)\\
    =~& \frac{2}{\alpha}\left(R(h_{\hat{\theta}},g_0) -R(h_{\alpha,\Hcal}^*,g_0) + R(h_{\alpha, \Hcal}^*,g_0)-R(h_{\alpha}^*,g_0)\right) + O\left(\alpha^{\min\{2,\beta\}}\right)\\
    \leq~& \frac{2}{\alpha} \left(R(h_{\hat{\theta}},g_0) - R(h_{j_{\alpha}^*},g_0) + R(h_{j_{\alpha}^*},g_0) -R(h_{\alpha,\Hcal_{j}}^*,g_0) + R(h_{\alpha, \Hcal_{j}}^*,g_0)-R(h_{\alpha}^*,g_0)\right) \tag{for any $j$}\\
    ~&+O\left(\alpha^{\min\{2,\beta\}}\right)\\
    \leq ~&\frac{2}{\alpha} \left(R(h_{\hat{\theta}},g_0) - R(h_{j_{\alpha}^*},g_0) + R(h_{j},g_0) -R(h_{\alpha,\Hcal_{j}}^*,g_0) + R(h_{\alpha, \Hcal_{j}}^*,g_0)-R(h_{\alpha}^*,g_0)\right)\\
    ~&+O\left(\alpha^{\min\{2,\beta\}}\right)\\
    \leq ~& O\big(\frac{\log(M)}{\alpha n} + \frac{\delta_{n,\Gcal}}{\alpha} + \alpha^{\min\{2,\beta\}} + \frac{\delta_{n,j}}{\alpha^2}\big) \tag{since $R(h_{j},g_0) -R(h_{\alpha,\Hcal_{j}}^*,g_0) \leq O\left(\frac{\delta_{n,j}}{\alpha}\right)$}\\
    ~&+\frac{2}{\alpha}\left(R(h_{\alpha, \Hcal_{j}}^*,g_0)-R(h_{\alpha}^*,g_0)\right)\\
    \leq ~& O\big(\alpha^{\min\{2,\beta\}} + \frac{\delta_{n,j}}{\alpha^2}\big) + \frac{2}{\alpha}\left(R(h,g_0)-R(h_{\alpha}^*,g_0)\right) \tag{for any $h\in \Hcal_{j_{\alpha}^*}$}\\ 
\end{align*}
For any function class $\Hcal$, we have:
\begin{align*}
    R(h,g_0)-R(h_{\alpha}^*,g_0) = ~& \|\Tcal(h-h_{\alpha}^*)\|^2 + \alpha \|h-h_{\alpha}^*\|^2\\
    \leq ~& 2\|\Tcal(h-h_0)\|^2 + 2\|\Tcal(h_{\alpha}^*-h_0)\|^2 + 2\alpha\|h-h_0\|^2 + 2\alpha \|h_{\alpha}^* - h_0\|^2\\
    \leq~& 4\|h-h_0\|^2 + O\big(\|w_0\|^2\alpha^{\min\{\beta+1,2\}}\big)\\
\end{align*}
Hence, we can choose $h$ that attains $\min_{\Hcal}\|h-h_0\| = \epsilon_{\Hcal}$. Combining, we get that: 
\begin{align*}
    \|h_{\hat{\theta}} - h_0\|^2 \leq \min_{j}O\big(\alpha^{\min\{\beta+1,2\}-1} + \frac{\delta_{n,j} + \epsilon_\Gcal}{\alpha^2} + \frac{1}{\alpha} \epsilon_{\Hcal}\big)
\end{align*}
\begin{lemma}
    Under the same assumptions as Theorem \ref{thm:oracle_model_selection}, we have that for any $\delta>0$, with probability $>1-\delta$:\\
    The output of \textbf{Best-ERM}, $\hat{\theta}$, satisfies:\\
    \begin{align*}
        \|h_{\hat{\theta}}~&-h_0\|^2 \leq 6 \|h_{j^*}-h_0\|^2 + O\left(\|h_{\alpha}^*-h_0\|^2 \right)\\
        ~&+ \frac{1}{\alpha} O\left(\frac{\log(M/\delta)}{n} + \|h_{j_{\alpha}^*} - h_{\alpha}^*\|^2 +\frac{1}{\alpha} \|\Tcal (h_{j_{\alpha}^*} - h_{\alpha}^*)\|^2 + \error(\hat{g})^2\right)
    \end{align*}
    The output of \textbf{Convex-ERM}, $\hat{\theta}$, satisfies:\\
    \begin{align*}
        \|h_{\hat{\theta}}~&-h_0\|^2 \leq 6 \|h_{\theta^*}-h_0\|^2 +16 \|h_{j_{\alpha}^*} - h_{\alpha}^*\|^2+ O\left(\|h_{\alpha}^*-h_0\|^2 \right)\\
        ~&+ \frac{1}{\alpha} O\left(\frac{M +\log(1/\delta)}{n} + \|\Tcal (h_{j_{\alpha}^*} - h_{\alpha}^*)\|^2 + \error(\hat{g})^2\right)
    \end{align*}
    The output of \textbf{Q-aggregation}, $\hat{\theta}$, satisfies:
        \begin{align*}
        \|h_{\hat{\theta}}~&-h_0\|^2 \leq 6 \|h_{j^*}-h_0\|^2 + 8\|h_{j_{\alpha}^*} - h_{\alpha}^*\|^2  + O\left(\|h_{\alpha}^*-h_0\|^2 \right)\\
        ~&+ \frac{1}{\alpha} O\left(\frac{\log(M/\delta)}{n} + \|\Tcal (h_{j_{\alpha}^*} - h_{\alpha}^*)\|^2 + \error(\hat{g})\right)
    \end{align*}
\end{lemma}
Since we need to bound $\|(\Tcal-\hat{\Tcal})(h)\|_1$ by $\error(\hat{g})\|h\|_2$, we will consider the MLE approach for estimating $\hat{g}$.
Applying previous results for the iterated thinknov regularized algorithm which is used to trained the candidate models, we get that when we choose $\alpha = \delta_n^{\frac{1}{\min\{2m,\beta\} + 2}}$:
\begin{align*}
    \|h_{j^*_{\alpha}}-h_{\alpha}^*\|^2 =~& O\left(16^{2m}\delta_n^{\frac{\min\{2m,\beta\}}{\min\{2m,\beta\}+2}}\right)\\
    \|\Tcal (h_{j^*_{\alpha}}-h_{m,*})\|^2 \leq~& O\left(\frac{\delta_n^2}{\alpha} + \delta_n + \alpha \|\hat{h}_{m-1}-h_{m-1,*}\|^2 \right)\\
    =~&O\left(16^{2m}\delta_n^{\frac{\min\{2m,\beta\}+1}{\min\{2m,\beta\}+2}}\right)\\
    \error(\hat{g}) =~& O(\delta_n)
\end{align*}
Moreover, since
\begin{align*}
    \alpha \|h_{\alpha}^* - h_{m,*}\|^2 + \|\Tcal (h_{\alpha}^* - h_{m,*})\|^2 = R(h_{m,*},g_0) - R(h_{\alpha}^*,g_0)
\end{align*}

\subsection{Best-ERM for $f(x)$}
In this section, we will invoke Corollary 12 from \cite{lan2023causal} to provide a bound for the Best-ERM approach. First, we need to verify that the assumptions for the corollary holds. 

\subsection{Convex-ERM for $f(x)$}
\begin{theorem}[Convex aggregation via ERM]\label{convex-erm}
Assume that $Y$ is almost surely bounded by $b$, each candidate model $f_j$ is almost uniformly bounded in $[-b,b]$ almost surely. Moreover, assume that $g^*$ and all $g\in \mathcal{G}$ is bounded away from 0, i.e. exist $c_0>0$ such that $g^*, g\geq c_0$. Then, for any $\delta>0$, with probability$`>1-\delta$, the output of ERM using the loss function described above, $\hat{\theta}$, satisfies:
\begin{align*}
    R(\hat{\theta},g_0) - \min_{\theta\in\Theta}R(\theta,g_0) \leq O\left(\frac{M+\log(1/\delta)}{n} + \error(\hat{g})^2\right)
\end{align*} 
where $\error(\hat{g}) = 2(\frac{1}{c_0}+1)^2\E[H(\hat{g},g_0)^2]^\frac{1}{2}$.
\end{theorem}
Let $\theta^*_\alpha:=\argmin_{\theta}R(\theta, g_0)$. By strong convexity, we can rewrite this as:
\begin{align*}
\|f_{\hat{\theta}} - f_{\theta^*_\alpha}\|^2 \leq \frac{1}{\alpha} O\left(\frac{M+\log(1/\delta)}{n} + \error(\hat{g})^2\right)
\end{align*}

To get a MSE regret bound:
\begin{align*}
    \|f_0-f_{\hat{\theta}}\|^2 \leq ~& 2\|f_0 -f_{\theta^*}\|^2 + 4 \|f_{\theta^*}-f_{\hat{\theta}}\|^2\\
    \leq ~& 6\|f_0 -f_{\theta^*}\|^2 + 4\|f_{\hat{\theta}}-f_0\|^2\\
    \leq ~&  6\|f_0 -f_{\theta^*}\|^2  + 8\|f_{\hat{\theta}}-f_{\alpha}^*\|^2 + 8\|f_0-f_{\alpha}^*\|^2\\
    \leq ~& 6\|f_0 -f_{\theta^*}\|^2 + 16 \|f_{\hat{\theta}}-f_{\theta_{\alpha}^*}\|^2 + 16 \|f_{\theta_{\alpha}^*}-f_{\alpha}^*\|^2 + O\left(\alpha^{\max\{\beta,2\}}\|w_0\|\right)\\
\end{align*}
% Let $\theta^*$ be the optimal convex aggregate when no regularization is used, then we get:
% \begin{align*}
%     \|f_{\hat{\theta}} - f_{\theta^*}\|^2 \leq~& 2\|f_{\hat{\theta}} - f_{\theta^*_\alpha}\|^2 + 2\|f_{\theta^*_\alpha}-f_{\theta^*} \|^2\\
%     \leq~&\frac{1}{\alpha} O\left(\frac{M+\log(1/\delta)}{n} + \error(\hat{g})^2\right) +2\|f_{\theta^*_\alpha}-f_{\theta^*} \|^2
% \end{align*}
% then we get:
% \begin{align*}
%     \|f_{\hat{\theta}} - f_0\|^2 - \min_j \|f_j-f_0\|^2 \leq \|f_{\hat{\theta}} - f_0\|^2
% \end{align*}
% By strong convexity, we can bound the last term:
% \begin{align*}
%     \|f_{\theta_{\alpha}^*}-f_{\theta^*}\|^2 \leq~& \frac{1}{\alpha}\left(\|\Tcal(f_0-f_{\theta^*})\|^2 + \alpha\|f_{\theta^*}\|^2 -\|\Tcal(f_0-f_{\theta^{*}_{\alpha}})\|^2 - \alpha\|f_{\theta^{*}_{\alpha}}\|^2 -\|\Tcal(f_{\theta_{\alpha}^*}-f_{\theta^*})\|^2 \right)\\
%     \leq~&\|f_{\theta^*}\|^2 - \|f_{\theta^{*}_{\alpha}}\|^2 -\frac{1}{\alpha}\|\Tcal(f_{\theta_{\alpha}^*}-f_{\theta^*})\|^2
% \end{align*}
% Let $f_{\alpha}^*$ be defined in Eqn \ref{eq:defs}, we can bound the last term:
% \begin{align*}
%       \|f_{\theta_{\alpha}^*}-f_{\theta^*}\|^2 \leq~& 2\|f_{\theta_{\alpha}^*} - f_{\alpha}^*\|^2 + 2\|f_{\theta^*}-f_{\alpha}^*\|^2\\
%       \leq~&2\|f_{\theta_{\alpha}^*} - f_{\alpha}^*\|^2 + 4\|f_{\theta^*}-f_0\|^2 + 4\|f_{\alpha}^*-f_0\|^2\\
%       \leq~&2\|f_{\theta_{\alpha}^*} - f_{\alpha}^*\|^2 + 4\|f_{\theta^*}-f_0\|^2 + O\left(\alpha^{\max\{\beta,2\}}\|w_0\|^2\right)
% \end{align*}
% To get a MSE regret bound:
% \begin{align*}
%     \|f_0-f_{\hat{\theta}}\|^2 \leq~& 2\|f_0-f_{\theta_{\alpha}^*}\|^2 +  2\|f_{\theta_{\alpha}^*}-f_{\hat{\theta}}\|^2\\
%     \leq ~& 4\|f_0-f_{\theta^*}\|^2 + 4\|f_{{\theta_{\alpha}^*}} - f_{\theta^*}\|^2 +\frac{1}{\alpha} O\left(\frac{M+\log(1/\delta)}{n} + \error(\hat{g})^2\right) \label{eqn:trig_split}\\
%      \leq~& 4\|f_0-f_{\theta^*}\|^2 + 8\|f_{\theta_{\alpha}^*} - f_{\alpha}^*\|^2+8\|f_{\theta^*}-f_{\alpha}^*\|^2+\frac{1}{\alpha} O\left(\frac{M+\log(1/\delta)}{n} + \error(\hat{g})^2\right)\\
%     \leq~& 20\|f_0-f_{\theta^*}\|^2 + 8\|f_{\theta_{\alpha}^*} - f_{\alpha}^*\|^2+\frac{1}{\alpha} O\left(\frac{M+\log(1/\delta)}{n} + \error(\hat{g})^2\right) + O\left(\alpha^{\max\{\beta,2\}}\|w_0\|^2\right)
% \end{align*}

Now, we wish to bound the estimation error term: $\|f_{\theta_{\alpha}^*} - f_{\alpha}^*\|^2$.Let $j^*_\alpha=\argmin_j \alpha \|f_j\|^2 + \|\Tcal (f_j-f_0)\|$.
By optimality of $f_{\alpha}^*$ (i.e. assuming that $f_{\alpha}^*$ satisfies the first order condition),
\begin{align*}
    \alpha \|f_{\alpha}^*-f_{\theta_{\alpha}^*}\|^2 + \|\Tcal(f_{\alpha}^*-f_{\theta_{\alpha}^*})\|^2 =~& R(f_{\theta_{\alpha}^*},g_0) - R(f_{\alpha}^*,g_0)\\
    \leq~& R(f_{j^*_\alpha},g_0) - R(f_{\alpha}^*,g_0)\\
    =~&\alpha \|f_{j^*_\alpha}-f_{\alpha}^*\|^2 + \|\Tcal (f_{j^*_\alpha}-f_{\alpha}^*)\|^2\\
    % \leq~& (1+\alpha)\|f_{j^*_\alpha}-f_{\alpha}^*\|^2
\end{align*}
Thus, we get that 
Note that, if the same $\alpha$ is used as in learning the candidate models, we can bound $\|f_{\theta_{\alpha}^*} - f_{\alpha}^*\|^2 \leq \|f_{j^*_\alpha} - f_{\alpha}^*\|^2 + \frac{1}{\alpha}\|\Tcal (f_{j^*_\alpha} - f_{\alpha}^*)\|^2$.
Combining, we get:
\begin{align*}
    \|f_0-f_{\hat{\theta}}\|^2 \leq~& 6\|f_0-f_{\theta^*}\|^2 + 16 \|f_{j^*_\alpha} - f_{\alpha}^*\|^2  + O\left(\alpha^{\max\{\beta,2\}}\|w_0\|^2\right)\\
    ~& +\frac{1}{\alpha} O\left(\frac{M+\log(1/\delta)}{n} + \error(\hat{g})^2 + \|\Tcal (f_{j^*_\alpha} - f_{\alpha}^*)\|^2\right)
\end{align*}
\subsubsection{Proof for Convex Aggregation}
Before proving the oracle inequality for convex aggregation via ERM, we will first state the auxiliary lemma that bounds the empirical process.
\begin{assumption}[Separability Condition]\label{Separability}
Let $G$ be a real-valued function class defined on $\mathcal{Z}$. There exist $G_0\subset G$ such that $G_0$ is countable and for any $g \in G$, there exist a sequence $(g_k)_k$ such that for any $z\in \mathcal{Z}$, $(g_k(z))_k \rightarrow g(z)$ as $k \rightarrow \infty$.
\end{assumption}
\begin{lemma}[\cite{bartlett2006empirical}]\label{concentration}
    There exist an absolute constant $c_0>0$ such that the following holds. Let $G$ be a real-valued function class defined on $\mathcal{Z}$ such that Assumption \ref{Separability} holds. Further assume that for all $g\in G$, $Pg^2\leq B_1 Pg$ and $\sup_{g\in G}\|g\|_{L_\infty}= B_2$ for constants $B_1$ and $B_2$. Let $\lambda_n$ be such that:
    \begin{align*}
        \E\left[\sup_{Pg\leq\lambda_n, g\in \text{Star}(G)}|(P-P_n)g|\right]\leq \frac{1}{8}\lambda_n
    \end{align*}
    For every $\delta>0$, with probability greater than $1-\delta$, for every $g\in G$, we have:
    \begin{align*}
        |(P-P_n)g|\leq \frac{1}{2}\max\{Pg,\rho_n(\delta)\}
    \end{align*}where $\rho_n(\delta)=\max \{\lambda_n, \frac{c_0(B_1+B_2)\log(1/\delta)}{n}\}$
\end{lemma}
\begin{proof}[Proof of Theorem \ref{convex-erm}]
Let $\mathcal{L}_{\Theta,g} = \{\mathcal{L}_{\theta,g}:\theta\in \Theta\}$, where $\mathcal{L}_{\theta,g} :=\ell_{\theta,g} - \ell_{\theta^*_\alpha,g} $. In addition, we can decompose $\mathcal{L}_{\theta,g}$ into $\mathcal{L}_{\theta,g}^0:=\ell^0_{\theta,g} - \ell_{\theta^*_\alpha,g}^0$ and  $\mathcal{L}_{\theta}^r:=r_{\theta}-r_{\theta^*_\alpha}$, such that $\mathcal{L}_{\theta,g} = \mathcal{L}_{\theta,g}^0+\mathcal{L}_{\theta}^r$.
\begin{align*}
    R(\hat{\theta},g_0) - R(\theta^*_{\alpha}, g_0) =~& P\mathcal{L}_{\hat{\theta},g_0}\\
    =~& P\mathcal{L}_{\hat{\theta},\hat{g}} - (P\mathcal{L}_{\hat{\theta},\hat{g}}-P\mathcal{L}_{\hat{\theta},g_0})\\
    \leq~& P\mathcal{L}_{\hat{\theta},\hat{g}} + O(\error(\hat{g})) \tag{See Section \ref{proof_q_agg}}\\
    =~& P\mathcal{L}^0_{\hat{\theta},\hat{g}} + P\mathcal{L}^r_{\hat{\theta}} + O(\error(\hat{g}))\\
    \leq~& 2 P_n\mathcal{L}^0_{\hat{\theta},\hat{g}} + 2P_n\mathcal{L}^r_{\hat{\theta}} + \rho^0_n(\delta) + \rho^r_n(\delta) +O(\error(\hat{g})) \tag{by Lemma \ref{concentration}}\\
    =~&2 P_n \mathcal{L}_{\hat{\theta},\hat{g}} + \rho^0_n(\delta) + \rho^r_n(\delta) +O(\error(\hat{g}))\\
    \leq~& \rho^0_n(\delta) + \rho^r_n(\delta) +O(\error(\hat{g}))
\end{align*}
It remains to show that $\rho^0_n(\delta) $ and $ \rho^r_n(\delta)$ are on the order of $\frac{M}{n}$. First, we observe that both function classes $\mathcal{L}^{0}_{\Theta,\hat{g}}$ and $\mathcal{L}^{r}_{\Theta}$ satisfies the separability condition. Since $Y$ and $f$ are bounded by $C_Y$ and $C_f$ respectively, we have that $P(\mathcal{L}^{0}_{\theta,\hat{g}})^2\leq B_1 P\mathcal{L}^{0}_{\theta,\hat{g}}$ and $P(\mathcal{L}^{r}_{\theta})^2\leq B_2 P\mathcal{L}^{r}_{\theta}$ for constants $B_1$ and $B_2$.
Let $\mathcal{L}_{\Theta}$ denote either $\mathcal{L}^{0}_{\Theta,\hat{g}}$ and $\mathcal{L}^{r}_{\Theta}$. Then by the peeling argument in \cite{bartlett2012}, for any $\lambda>0$:
\begin{align*} 
\left\{f\in\text{Star}(\mathcal{L}_{\Theta}), P(f)\leq \lambda\right\}\subset \cup_{i=0}^{\infty}\left\{\alpha f: 0\leq \alpha \leq 2^{-i}, f\in \mathcal{L}_{\Theta}, Pf\leq 2^{i+1}\lambda\right\}
\end{align*}
which implies
\begin{align*}
\E\left[\sup_{P\ell\leq\lambda, f \in \mathcal{L}_{\Theta}}|(P-P_n)f|\right]\leq \sum_{i=0}^{\infty}2^{-i}\E\left[\sup_{Pf\leq 2^{i+1}\lambda, f \in \mathcal{L}_{\Theta}}|(P-P_n)f|\right]
\end{align*}
Thus, if $\lambda_n/8$ upper bounds the RHS, then it also satisfies the condition in Lemma \ref{concentration}. 
Let $\{e_1, \dots, e_{M'}\}$, $M'\leq M$, be an orthogonal basis of the linear subspace, $S$, spanned by the function class $\mathcal{F}$, then we have:
\begin{align*}
    \E\left[\sup_{Pf\leq\lambda, f \in \mathcal{L}_{\Theta}}|(P-P_n)f|\right] \leq~& \max\{1,\alpha\}8b \E\left[\sup_{Pf^2\leq\lambda , f \in S}\left|\frac{1}{n}\sum_{i=1}^n \epsilon_i f(X_i)\right|\right] \tag{By Symmetrization}\\
    \leq~& \max\{1,\alpha\}8b \E\left[\sup_{\beta\in \mathbb{R}^{M'} , \|\beta\|_2 \leq \sqrt{\lambda}}\left|\frac{1}{n}\sum_{i=1}^n \epsilon_i \left(\sum_{j=1}^{M'}\beta_j e_j(X_i)\right)\right|\right]\\
    \leq~& \max\{1,\alpha\}8b\sqrt{\lambda}\E\left(\sum_{j=1}^{M'}\left(\frac{1}{n}\sum_{i=1}^{n}\epsilon_i e_j(X_i)\right)^2\right)^{\frac{1}{2}} \tag{By the Cauchy-Schwarz inequality}\\
    \leq~& \max\{1,\alpha\}8b\sqrt{\frac{M'\lambda}{n}}
\end{align*}
Thus, by the pealing argument, we showed that $\lambda_n = O\left(\frac{M}{n}\right)$.
\end{proof}


\subsection{Q-aggregation for $f(x)$}
We consider the modified loss, which penalizes the weight put on each model, based on the individual model performance:
\begin{align*}
    \tilde{\ell}_{\theta,g}(Y,Z,X) =~& (1-\nu) \ell_{\theta,g}(Y,Z,X) + \nu \sum_{j} \theta_j \ell_{e_j,g}(Y,Z,X) &
    \tilde{R}(\theta,g) =~& P\tilde{\ell}_{\theta,g}(Z)
\end{align*}
We also define with:
\begin{align*}
    K(\theta) :=~&\sum_{j} \theta_j \log(1/\pi_j)
\end{align*} where $\pi_j$ is the prior for each candidate model.
Then the Q-aggregation ensemble is defined as:
\begin{align*}
    \hat{\theta} =~& \arg\min_{\theta} P_n\tilde{\ell}_{\theta, \hat{g}}(Y,Z,X) + \frac{\beta}{n} K(\theta)
\end{align*}

% Simply using ERM yields:
% \begin{align*}
%     \|f_{\hat{\theta}}-f_{\theta*}^{\alpha}\|^2 \leq~& 2 \|f_{\hat{\theta}}-f_{\theta*}^{\alpha}\|^2 + 2\|f_{\theta*}^{\alpha} - f_0\|^2\\
%     \leq~& O\left(\left(\frac{\delta_n}{\alpha}\right)^{\frac{2}{2-\gamma}}+\|f_{\theta*}^{\alpha} - f_0\|^2\right)
% \end{align*}
% Thus, the MSE regret with respect to the best convex combination (with the regularizer) is:
% \begin{align*}
%     \|f_{\hat{\theta}}-f_{\theta*}^{\alpha}\|^2 -\|f_{\theta*}^{\alpha} - f_0\|^2 \leq  O\left(\left(\frac{\delta_n}{\alpha}\right)^{\frac{2}{2-\gamma}}+\|f_{\theta*}^{\alpha} - f_0\|^2\right)
% \end{align*}
% Since we don't need to pick $\alpha$ according to the source condition, we get that:
% \begin{align*}
%     \|f_{\hat{\theta}}-f_{\theta*}^{\alpha}\|^2 -\|f_{\theta*}^{\alpha} - f_0\|^2 \leq  O\left(\left(\frac{M\log(n)}{n}\right)^{\frac{1}{2-\gamma}}+\|f_{\theta*}^{\alpha} - f_0\|^2\right)
% \end{align*}

\begin{theorem}[Oracle Inequality for Q-aggregation]\label{thm:q_agg}
Assume that $Y$ is almost surely bounded by $C_Y$, each candidate model $f_j$ is almost uniformly bounded in $[-C_f,C_f]$ almost surely. Moreover, assume that $g^*$ and all $g\in \mathcal{G}$ is bounded away from 0, i.e. exist $c_0>0$ such that $g^*, g\geq c_0$. Then, for any $\mu \leq \min\{\frac{\alpha(1-\nu)}{7}, \frac{1-\nu}{3}, \frac{\alpha\nu}{5}\}$  and $\beta \geq \{ \frac{128C_{Y,f}(1-\nu)^2}{\mu}, 16\sqrt{6}(1-\nu)C_f \sqrt{C_{Y,f}},\frac{4\nu(4\nu C_{Y,f} + \mu C_{\ell})}{\mu}\}$, where $C_{Y,f} = (C_Y+C_f)^2+\alpha^2 C_f^2$ and $C_{\ell} = (C_Y+C_f)^2+\alpha C_f^2$, w.p. $1-\delta$:
\begin{align*}
    R(\hat{\theta}, g_0) \leq~& \min_{j=1}^M\left( R(e_j, g_0)  +\frac{\beta}{n}\log(1/\pi_j)\right) + \widetilde{O}\left(  \frac{\log(1/\delta)}{n} + \frac{1}{\mu}\error(\hat{g})^{2}\right) \\
\end{align*} where $\error(\hat{g}) =2(\frac{1}{c_0} + 1)^2\mathbb{E}\left[H(\hat{g},g_0)^2\right]^{\frac{1}{2}}$.
\end{theorem}
Consider the case where we do not incorporate any priors on the models, then:
\begin{align*}
    R(\hat{\theta}, g_0) -\min_{j=1}^M R(e_j, g_0)\leq~& \widetilde{O}\left(\frac{\log(M/\delta)}{n} + \frac{1}{\mu}\error(\hat{g})^{2}\right)\\
\end{align*}
Moreover, by strong convexity,
\begin{align*}
    \alpha\|f_{\hat{\theta}} - f_{\alpha}^*\|^2 \leq~&  R(f_{\hat{\theta}}, g_0)  -  R(f_{\alpha}^*, g_0)\\
    \leq ~& \widetilde{O}\left(\frac{\log(M/\delta)}{n} + \frac{1}{\mu}\error(\hat{g})^{2}\right) + R(f_{j_{\alpha}^*}, g_0)  -  R(f_{\alpha}^*, g_0)\\
    \leq ~& \widetilde{O}\left(\frac{\log(M/\delta)}{n} + \frac{1}{\mu}\error(\hat{g})^{2}\right) + \alpha \|f_{j_{\alpha}^*}-f_{\alpha}^*\|^2 +\|\Tcal (f_{j_{\alpha}^*}-f_{\alpha}^*)\|^2
\end{align*} where $j^*_{\alpha} = \arg\min_j R(e_j,g_0)$.
Further, let $j^*$ denote $\argmin \|f_j-f_0\|^2$, we have:
\begin{align*}
    \|f_{\hat{\theta}}-f_0\|^2 \leq~&  2\|f_0 - f_{j^*}\|^2 + 2\|f_{j^*}-f_{\hat{\theta}}\|^2\\
    \leq~& 6\|f_0 - f_{j^*}\|^2 + 4 \|f_0-f_{\hat{\theta}}\|^2\\
    \leq~& 6\|f_0 - f_{j^*}\|^2  + 8 \|f_0-f_{\alpha}^*\|^2 + 8\|f_{\alpha}^*-f_{\hat{\theta}}\|^2\\
    \leq ~& 6\|f_0 - f_{j^*}\|^2 + 8\|f_{j_{\alpha}^*} - f_{\alpha}^*\|^2 + O\left(\alpha^{\max\{2, \beta\}}\|w_0\|^2\right) \\
    +~&\frac{1}{\alpha} O\left(\frac{\log(M/\delta)}{n} + \|\Tcal (f_{j_{\alpha}^*} - f_{\alpha}^*)\|^2 + \error(\hat{g})^2\right)\\
\end{align*}


% \begin{align*}
%     R(\hat{\theta}, g_0) -\min_{j=1}^M R(e_j, g_0) =~& \|\Tcal(f_0-f_{\hat{\theta}})\|^2 -\|\Tcal(f_0-f_{j_{m}})\|^2 +\alpha(\|f_{\hat{\theta}}\|^2-\|f_{j_{m}}\|^2)\\
%     \geq~& \|\Tcal(f_{\hat{\theta}}-f_{j^*})\|^2 + \alpha \|f_{\hat{\theta}}-f_{j^*}\|^2 + \E[D_{f}\ell_{f_{jm},g_0}[f_{\hat{\theta}}-f_{j^*}]]\\ =~&\|\Tcal(f_{\hat{\theta}}-f_{j^*})\|^2 + \alpha \|f_{\hat{\theta}}-f_{j^*}\|^2+2\E[\Tcal(f_0-f_{j^*})\Tcal(f_{\hat{\theta}}-f_{j^*}) + \alpha f_{j^*}(f_{\hat{\theta}}-f_{j^*})]\\
% \end{align*}
% Combining, we get:
% \begin{align*}
%     \alpha \|f_{\hat{\theta}}-f_{j^*}\|^2 \leq~& \widetilde{O}\left(\frac{\log(M/\delta)}{n} + \frac{1}{\mu}\error(\hat{g})^{2}\right) + 2\|\Tcal(f_0-f_{j^*})\|\|\Tcal(f_{\hat{\theta}}-f_{j^*})\| + 2\alpha\|f_{j^*}\|\|f_{\hat{\theta}}-f_{j^*}\|\\
%     \leq~&\widetilde{O}\left(\frac{\log(M/\delta)}{n} + \frac{1}{\mu}\error(\hat{g})^{2}\right) + 2\left(\|\Tcal(f_0-f_{j^*})\| + \alpha C_f\right)\|f_{\hat{\theta}}-f_{j^*}\|\\
% \end{align*}
% Thus,
% \begin{align*}
%         \|f_{\hat{\theta}}-f_{j^*}\|^2 \leq~& \frac{1}{\alpha}\widetilde{O}\left(\frac{\log(M/\delta)}{n} + \frac{1}{\mu}\error(\hat{g})^{2} +\left(\|\Tcal(f_0-f_{j^*})\| + \alpha\|f_{j^*}\|\right)^2\right)
% \end{align*}


\subsection{Proof of Oracle Inequality for Q-aggregation}\label{sec:Q-agg}
We consider the modified loss, which penalizes the weight put on each model, based on the individual model performance:
\begin{align*}
    \tilde{\ell}_{\theta,g}(Y,Z,X) =~& (1-\nu) \ell_{\theta,g}(Y,Z,X) + \nu \sum_{j} \theta_j \ell_{e_j,g}(Y,Z,X), &
    \tilde{R}(\theta,g) =~& \E\tilde{\ell}_{\theta,g}(Z)
\end{align*}
We also define with:
\begin{align*}
    K(\theta) :=~&\sum_{j} \theta_j \log(1/\pi_j),
\end{align*} where $\pi_j$ is the prior for each candidate model.
Then the Q-aggregation ensemble is defined as:
\begin{align*}
    \hat{\theta} =~& \arg\min_{\theta} P_n\tilde{\ell}_{\theta, \hat{g}}(Y,Z,X) + \frac{\beta}{n} K(\theta),
\end{align*}

\begin{theorem}[Oracle Inequality for Q-aggregation]\label{thm:q_agg}
Assume that $Y$ is almost surely bounded by $C_Y$, each candidate model $f_j$ is almost uniformly bounded in $[-C_f,C_f]$ almost surely. Moreover, assume that $g_0$ and all $g\in \mathcal{G}$ is bounded away from 0, i.e. exist $c_0>0$ such that $g_0, g\geq c_0$. Then, for any $\mu \leq \min\{\frac{\alpha(1-\nu)}{7}, \frac{1-\nu}{3}, \frac{\alpha\nu}{5}\}$  and $\beta \geq \{ \frac{128C_{Y,f}(1-\nu)^2}{\mu}, 16\sqrt{6}(1-\nu)C_f \sqrt{C_{Y,f}},\frac{4\nu(4\nu C_{Y,f} + \mu C_{\ell})}{\mu}\}$, where $C_{Y,f} = (C_Y+C_f)^2+\alpha^2 C_f^2$ and $C_{\ell} = (C_Y+C_f)^2+\alpha C_f^2$, w.p. $1-\delta$:
\begin{align*}
    R(\hat{\theta}, g_0) \leq~& \min_{j=1}^M\left( R(e_j, g_0)  +\frac{\beta}{n}\log(1/\pi_j)\right) + \widetilde{O}\left(  \frac{\log(1/\delta)}{n} + \frac{1}{\mu}\error(\hat{g})^2 + \error(\hat{g})\right), \\
\end{align*} where $\error(\hat{g}) =2(\frac{1}{c_0} + 1)^2\mathbb{E}\left[H(\hat{g},g_0)^2\right]^{\frac{1}{2}}$.
Moreover, when we do not incorporate any priors on the models, then:
\begin{align*}
    R(\hat{\theta}, g_0) -\min_{j=1}^M R(e_j, g_0)\leq~& \widetilde{O}\left(\frac{\log(M/\delta)}{n} + \frac{1}{\mu}\error(\hat{g})^2 +\error(\hat{g})\right).
\end{align*}
\end{theorem}

Before diving into the proof, we first introduce some population quantities that will play a key role in our analysis. Define:
\begin{align*}
    V(\theta) :=~& \sum_{j} \theta_j \|f_j - f_{\theta}\|^2, 
\end{align*}
$$    H(\theta, g) := \tilde{R}(\theta, g) + \mu V(\theta) 
    + \frac{\beta}{n} K(\theta),$$
and define:
\begin{align*}
    \theta^* =~& \arg\min_{\theta} \tilde{R}(\theta, g_0) + \mu V(\theta) + \frac{\beta}{n} K(\theta) 
    ~=~ \arg\min_{\theta} H(\theta, g_0), \\
\end{align*}
Here we also verify the strong convexity of the loss function in expectation with respect to some distance $d(\cdot,\cdot)$, i.e.
\begin{align*}
    \E\left[\ell_{f', g_0} - \ell_{f, g_0}\right]  - \E D_{f}\ell_{f, g_0}[f'-f]\geq \frac{\sigma}{2}d(f,f')^2
\end{align*} for some $\sigma$.

Using the loss function described above:
\begin{align*}
    ~&\E\left(\ell_{f', g_0} - \ell_{f, g_0}\right)  - \E D_{f}\ell_{f, g_0}[f'-f] \\=~& \mathbb{E}\left[-2Y(\Tcal(f'-f)) + (\Tcal(f'))^2 - (\Tcal(f))^2 + 2(Y-\Tcal(f))(\Tcal(f'-f))\right]\\
    ~&+ \alpha \mathbb{E}\left[(f')^2 - f^2 -2f(f'-f) \right]\\
    =~& \mathbb{E}\left[(\Tcal(f') - \Tcal(f))^2\right] + \alpha \mathbb{E}[(f'-f)^2] \\
    =~& \|\Tcal(f') - \Tcal(f)\|_{2}^{2} + \alpha \|f'-f\|_{2}^{2}.
\end{align*}

Now we introduce the shorthand notation:
\begin{align*}
    \tilde{{\cal E}}(\hat{g}) :=~& \E\big[\tilde{\ell}_{\hat{\theta},g_0}(Z) - \tilde{\ell}_{\theta^*,g_0}(Z)\big] - \E\big[\tilde{\ell}_{\hat{\theta},\hat{g}}(Z) - \tilde{\ell}_{\theta^*,\hat{g}}(Z)\big]\\
    {\cal E}(\hat{g}) :=~& \E\big[{\ell}_{\hat{\theta},g_0}(Z) - {\ell}_{\theta^*,g_0}(Z)\big] - \E\big[{\ell}_{\hat{\theta},\hat{g}}(Z) - {\ell}_{\theta^*,\hat{g}}(Z)\big]
    % ~\leq~ \error(\hat{g})\, \|f_{\hat{\theta}} - f_{\theta^*}\|^{\gamma}
\end{align*}

\begin{lemma}\label{lemma:diffndiff}
    Suppose that $|Y|$ is almost surely bounded by some constant $C_Y$, and $|f(X)|$ is almost surely uniformly bounded by a constant $C_f$. Then, 
    \begin{align*}
    {\cal E}(\hat{g}) 
    \leq~& \left(|C_Y-C_f|\|\hat{f}-f^*\|_{2}+C_f\|\Tcal(\hat{f}-f^*)\|_2\right)\error(\hat{g}),\\
     \tilde{{\cal E}}(\hat{g}) \leq~&
     \left((2C_Y + C_f)\|\hat{f}-f^*\|_{2}+C_f\|\Tcal(\hat{f}-f^*)\|_2 + 2C_f^2\right)\error(\hat{g}).
    \end{align*} where $\error(\hat{g}) =2(\frac{1}{c_0} + 1)^2\mathbb{E}\left[H(\hat{g},g_0)^2\right]^{\frac{1}{2}}$.
\end{lemma}
\begin{proof}
    \begin{align*}
    {\cal E}(\hat{g})
    =~& \E [-2Y (\Tcal-\hat{\Tcal})(\hat{f}-f^*) + (\Tcal(\hat{f}))^2-(\Tcal(f^*))^2 -(\hat{\Tcal}(\hat{f}))^2+(\hat{\Tcal}(f^*))^2]\\
    =~&\E[ -2Y (\Tcal-\hat{\Tcal})(\hat{f}-f^*) + (\Tcal-\hat{\Tcal})(\hat{f}-f^*)\hat{\Tcal}(\hat{f}+f^*)]\\
    ~&+ \E[\Tcal(\hat{f}-f^*)(\Tcal - \hat{\Tcal})(\hat{f}+f^*)]\\
    \leq~& \E[2|C_Y-C_f||(\Tcal-\hat{\Tcal})(\hat{f}-f^*)| + | \Tcal(\hat{f}-f^*)(\Tcal-\hat{\Tcal})(\hat{f}+f^*)|]
    % \leq~&\mathbb{E} \left[2C_Y \left(\int (f-f')^2g_0dx\right)^{\frac{1}{2}}\left(\int\left(\frac{g}{g_0}-1\right)^2g_0dx\right)^{\frac{1}{2}}\right] \tag{Here we need the true likelihood to be bounded away from $0$} \\
    % \leq~&\mathbb{E} \left\{2C_Y \mathbb{E}[(f-f')^2|Z]^{\frac{1}{2}}\mathcal{\chi}^2(g(X|Z),g_0(X|Z))^{\frac{1}{2}}\right\}\\
    % \leq~& 2C_Y \mathbb{E} \left[\mathbb{E}[(f-f')^2|Z] \right]^{\frac{1}{2}}\mathbb{E} \left[\mathcal{\chi}^2(g(X|Z),g_0(X|Z))\right]^{\frac{1}{2}}\\
    % =~&  2C_Y\|f-f'\|_{2} \mathbb{E} \left[\mathcal{\chi}^2(g(X|Z),g_0(X|Z))\right]^{\frac{1}{2}}\\
\end{align*}
First, we analyze the $(\Tcal-\hat{\Tcal})(\hat{f}-f^*)$ term:
\begin{align*}
    \mathbb{E}[|(\Tcal-\hat{\Tcal})(\hat{f}-f^*)|]\leq~&\mathbb{E}\left[\int |(\hat{f}-f^*)(g-g_0)|dx\right]\\
    \leq~&\mathbb{E}\left[\left(\int (\hat{f}-f^*)^2g_0dx\right)^{\frac{1}{2}}\left(\int\left(\frac{g}{g_0}-1\right)^2g_0dx\right)^{\frac{1}{2}}\right]\tag{Here we need the true likelihood to be bounded away from $0$}\\
    =~& \mathbb{E}\left[\mathbb{E}[(\hat{f}-f^*)^2|Z]^{\frac{1}{2}}\mathcal{\chi}^2(g(X|Z),g_0(X|Z))^{\frac{1}{2}}\right]\\
    \leq~& \mathbb{E}\left[\mathbb{E}[(\hat{f}-f^*)^2|Z] \right]^{\frac{1}{2}}\mathbb{E} \left[\mathcal{\chi}^2(g(X|Z),g_0(X|Z))\right]^{\frac{1}{2}}\\
    =~& \|\hat{f}-f^*\|_{2} \mathbb{E} \left[\mathcal{\chi}^2(g(X|Z),g_0(X|Z))\right]^{\frac{1}{2}}.
\end{align*}

Next, we analyze the $\Tcal(\hat{f}-f^*)(\Tcal-\hat{\Tcal})(\hat{f}+f^*)$ term:
\begin{align*}
    \mathbb{E}\left[| \Tcal(\hat{f}-f^*)(\Tcal-\hat{\Tcal})(\hat{f}+f^*)|\right] \leq~& 
    \mathbb{E}\left[|\Tcal(\hat{f}-f^*)|\mathbb{E}[(\hat{f}+f^*)^2|Z]^{\frac{1}{2}}\mathcal{\chi}^2(g(X|Z),g_0(X|Z))^{\frac{1}{2}}\right]\\
    \leq~& 2C_f\mathbb{E}\left[|\Tcal(\hat{f}-f^*)|\mathcal{\chi}^2(g(X|Z),g_0(X|Z))^{\frac{1}{2}}\right]\\
    \leq~&2C_f\|\Tcal(\hat{f}-f^*)\|_2\mathbb{E}\left[\mathcal{\chi}^2(g(X|Z),g_0(X|Z))\right]^{\frac{1}{2}}.
\end{align*}

Since $g_o(X|Z)\geq c_0$, we have:
\begin{align*}
    \mathcal{\chi}^2(g(X|Z),g_0(X|Z)) =~&\int\left(\frac{g}{g_0}-1\right)^2g_0dx\\
    =~&\int\left(\frac{\sqrt{g}}{\sqrt{g_0}}+1\right)^2\left(\frac{\sqrt{g}}{\sqrt{g_0}}-1\right)^2g_0dx\\
    \leq~&\int\left(\frac{1}{c_0}+1\right)^2\left(\frac{\sqrt{g}}{\sqrt{g_0}}-1\right)^2g_0dx\\
    =~&\left(\frac{1}{c_0}+1\right)^2 H(g,g_0)^2.
\end{align*}

Thus, we get:
\begin{align*}
    {\cal E}(\hat{g}) \leq~& \left(2|C_Y-C_f|\|\hat{f}-f^*\|_{2}+2C_f\|_2\Tcal(\hat{f}-f^*)\|\right)(\frac{1}{c_0} + 1)^2\mathbb{E}\left[H(g,g_0)^2\right]^{\frac{1}{2}}\\
    \leq~& \left(|C_Y-C_f|\|\hat{f}-f^*\|_{2}+C_f\|\Tcal(\hat{f}-f^*)\|_2\right)\error(\hat{g}),
\end{align*} where $\error(\hat{g}) =2(\frac{1}{c_0} + 1)^2\mathbb{E}\left[H(\hat{g},g_0)^2\right]^{\frac{1}{2}}$.

For $\tilde{\cal E}(\hat{g})$, we have that:
\begin{align*}
    \tilde{\cal E}(\hat{g}) =~& (1-\nu){\cal E}(\hat{g})+\E[-2Y(\Tcal -\hat{\Tcal})(\hat{f}-f^*) + \sum_{j} (\hat{\theta}_j-\theta_j^*)((\Tcal f_j)^2 -(\hat{\Tcal}f_j)^2)]\\
    = ~& (1-\nu){\cal E}(\hat{g}) +\E[-2Y(\Tcal -\hat{\Tcal})(\hat{f}-f^*) + \sum_j (\hat{\theta}_j-\theta_j^*)( \Tcal f_j +\hat{\Tcal}f_j)(\Tcal f_j -\hat{\Tcal}f_j)]\\
    \leq~&  (1-\nu){\cal E}(\hat{g}) + C_Y \|\hat{f}-f^*\|\error(\hat{g}) + 2\sum_j (\hat{\theta}_j+\theta_j^*) C_f\E[|(\Tcal-\hat{\Tcal})f_j|]\\
    \leq~& (1-\nu){\cal E}(\hat{g}) + C_Y \|\hat{f}-f^*\|\error(\hat{g}) + 2 C_f^2\error(\hat{g}).
\end{align*}

\end{proof}

\begin{lemma}\label{lem:prel-strong-conv} For any $\theta,\theta'\in \Theta$:
\begin{align*}\label{eqn:strong-conv-R-lemma}
    R(\theta, g)-R(\theta', g) \geq \langle \nabla R(\theta', g), \theta - \theta' \rangle + \|\Tcal(f_{\theta} - f_{\theta'})\|_2^2+\alpha\|f_{\theta}-f_{\theta'}\|^2.
\end{align*}
\end{lemma}
\begin{proof}
    Follows from a similar calculation as above.
\end{proof}

\begin{lemma}\label{lem:v-versus-tilde} 
For any $\theta\in \Theta$:
\begin{equation}
    \tilde{R}(\theta, g) \geq R(\theta, g) + \nu V(\theta).
\end{equation}
\end{lemma}
\begin{proof}
By Lemma~\ref{lem:prel-strong-conv}, we have that:
\begin{align*}
    R(e_j, g) - R(\theta, g) \geq \langle \nabla R(\theta, g), e_j - \theta\rangle +\alpha\|f_j - f_\theta\|^2.
\end{align*}
Multiplying each equation by $\theta_j$ and summing across $j$, yields:
\begin{align*}
    R(\theta, g) \leq \sum_j \theta_j R(e_j, g) - \alpha V(\theta),
\end{align*}
where we used the fact that $\sum_{j}\theta_j=1$ and $\sum_{j} \theta_j \langle \nabla R(\theta, g), e_j - \theta\rangle=0$.

Thus we have the following:
\begin{align*}
    \tilde{R}(\theta, g) =~& (1-\nu) R(\theta, g) + \nu \sum_j \theta_j R(e_j, g)\\
    \geq~& (1-\nu) R(\theta, g) + \nu R(\theta, g) +\nu\alpha V(\theta)\\
    =~& R(\theta, g) + \nu\alpha V(\theta).
\end{align*}
\end{proof}
\begin{lemma}\label{lem:strong-conv-tilde}
Let $\theta_*=\argmin_{\theta\in \Theta} H(\theta, g_0)$, then for any $\theta\in \Theta$:
\begin{align}
    H(\theta, g_0) - H(\theta^*, g_0) \geq \left(1-\nu - \mu\right) \|\Tcal(f_{\theta} - f_{\theta^*})\|^2 + \alpha(1-\nu)\|f_{\theta} - f_{\theta^*}\|^2.
\end{align}
\end{lemma}
\begin{proof}
By Lemma~\ref{lem:prel-strong-conv} we have:
\begin{align}\label{eqn:strong-conv-R}
    R(\theta, g_0)-R(\theta', g_0) \geq \langle \nabla R(\theta', g_0), \theta - \theta' \rangle + \|\Tcal(f_{\theta}-f_{\theta'})\|^2+\alpha\|f_{\theta}-f_{\theta'}\|^2.
\end{align}
Now let: 
\begin{align*}
    q(Z;\theta, g) =~& \sum_{j=1}^{M}\theta_{j} \ell(Z;f_{e_{j}}(X), g)\\
    L_e(g)=~&[\ell(Z;f_{e_{1}}(X),g)\cdots \ell(Z;f_{e_{M}}(X), g)]^{T}\\
    Q(\theta, g) =~& \mathbb{E}\left[q(Z;\theta, g)\right].
\end{align*}
Then we have:
\begin{align*}
    q(Z;\theta, g)-q(Z; \theta', g) =~&\sum_{j=1}^{M}\ell(Z;f_{e_{j}}(X), g)(\theta_{j}-\theta'_{j})
    = \langle{L_{e}(g), \theta-\theta'}\rangle\\
    =~& \langle \nabla q(Z; \theta', g), \theta-\theta' \rangle.
\end{align*}
By taking the expectation over $Z$:
\begin{align}\label{eqn:strong-conv-G}
    Q(\theta, g) - Q(\theta', g) = \sum_{j=1}^{M}\theta_{j}R(e_j, g)-\sum_{j=1}^{M}\theta_{j}'R(e_j, g)=\langle \nabla Q(\theta', g), \theta-\theta' \rangle.
\end{align}
By adding Eq.~\eqref{eqn:strong-conv-R} and \eqref{eqn:strong-conv-G} we have that: 
\begin{align*}
    \tilde{R}(\theta, g_0)-\tilde{R}(\theta', g_0) =~& (1 - \nu) (R(\theta, g_0)-R(\theta', g_0)) + \nu (Q(\theta, g_0) - Q(\theta', g_0))\\
    \geq~& (1-\nu) \langle \nabla R(\theta',g_0), \theta - \theta' \rangle + \nu \langle \nabla Q(\theta',g_0), \theta-\theta' \rangle\\
    ~&+ (1-\nu)\left(\|\Tcal(f_{\theta}-f_{\theta'})\|^2+\alpha\|f_{\theta}-f_{\theta'}\|^2\right) \\
    \geq~& \langle \nabla \tilde{R}(\theta',g_0), \theta - \theta' \rangle+ (1-\nu)\left(\|\Tcal(f_{\theta}-f_{\theta'})\|^2+\alpha\|f_{\theta}-f_{\theta'}\|^2\right) .
\end{align*}
 We also have that for any $\theta, \theta'$:\footnote{Since the function $V(\theta)$ can be interpreted as the expected variance of $f_j(X)$ over the distribution of $j$. Thus with simple manipulations it can be shown to be equal to $V(\theta)=\sum_j \theta_j \|f_j\|^2- \|f_\theta\|^2$. Since this is a quadratic in $\theta$ function, with Hessian $\mathbb{E}[FF']$, where $F$ is the vector $(f_1(X), \ldots, f_j(X))$, the above equality follows by an exact second order Taylor expansion and the observation that $(\theta-\theta')^\top\mathbb{E}[FF'](\theta-\theta') = \|f_\theta - f_{\theta'}\|^2$}
\begin{align*}
    V(\theta)-V(\theta') = \langle \nabla V(\theta'), \theta - \theta' \rangle - \|f_{\theta}-f_{\theta'}\|_{2}^2
\end{align*}
which implies that $V(\theta)$ is at most 2-strongly concave. 
We also have:
\begin{align*}
    K(\theta) - K(\theta') = \sum_j (\theta_j - \theta'_j) \log(1/\pi_j) = \langle \nabla K(\theta'), \theta - \theta' \rangle,
\end{align*}

Adding the three inequalities and substituting $\theta'$ with $\theta^*$ we have:
\begin{align*}
        H(\theta, g_0) - H(\theta^*, g_0) \geq~& \langle \nabla H(\theta^*, g_0), \theta-\theta^* \rangle + \left(1-\nu \right) \|\Tcal(f_{\theta} - f_{\theta^*})\|^2\\
        ~&+ (\alpha(1-\nu)-\mu)\|f_{\theta} - f_{\theta^*}\|^2.
\end{align*}
Note that by the optimality of $\theta^*$ over the convex set $\Theta$ we have that $\langle \nabla H(\theta^*, g_0), \theta-\theta^* \rangle \geq 0$, which immediately yields the statement we want to prove.
\end{proof}

\begin{lemma}\label{lem:main}
Define the centered and offset empirical processes:
\begin{align*}
    Z_{n,1} =~& (1-\nu)(\E-\E_n) (\ell_{\hat{\theta}, \hat{g}}(Z) - \ell_{\theta^*, \hat{g}}(Z)) - \mu \sum_{j} \hat{\theta}_j \|f_j - f_{\theta^*}\|^2 {-\frac{1}{s}K(\hat{\theta})}\\
    Z_{n,2} =~& \nu(\E-\E_n) \sum_{j} (\hat{\theta}_j-\theta_j^*)\ell_{e_j, \hat{g}}(Z) - \mu \hat{\theta}^\top H \theta^* {-\frac{1}{s}K(\hat{\theta})}\\
    H_{ij} =~& \|f_i - f_j\|^2 ,
\end{align*}
where $s$ is a positive number. 
% Assume that:
% \begin{align*}
%     {\cal E}(\hat{g}) :=~& P(\tilde{\ell}_{\hat{\theta},g_0}(Z) - \tilde{\ell}_{\theta^*,g_0}(Z)) - P(\tilde{\ell}_{\hat{\theta},\hat{g}}(Z) - \tilde{\ell}_{\theta^*,\hat{g}}(Z)) \\
%     \leq ~& \left(C_1\|\hat{f}-f^*\|_{2} + C_2\|\Tcal(\hat{f}-f^*)\|_2\right)\error(\hat{g})
% \end{align*} for some constants $C_1$ ad $C_2$.
and $\ell(Z; f(X), g_0)$ is $\sigma$-strongly convex in expectation with respect to $f$. For {$\beta \geq \frac{4n}{s}$} and $\mu \leq \min\{\frac{\alpha(1-\nu)}{7}, \frac{1-\nu}{3}, \frac{\alpha\nu}{5}\}$:
\begin{align*}
    R(\hat{\theta}) \leq \min_{j} [R(e_j, g_0){+ \frac{\beta}{n} \log(1/\pi_j)}] + 2\, (Z_{n,1} + Z_{n,2}) + O\big(\frac{1}{\mu}\error(\hat{g})^2 + \error(\hat{g})\big) .\label{eqn:lemma9}
\end{align*}
\end{lemma}
\begin{proof}
\begin{align*}
    \tilde{R}(\hat{\theta}, g_0) - \tilde{R}(\theta^*, g_0) =~& 
    \E(\tilde{\ell}_{\hat{\theta},g_0}(Z) - \tilde{\ell}_{\theta^*,g_0}(Z))\\
    =~& 
    \E(\tilde{\ell}_{\hat{\theta},\hat{g}}(Z) - \tilde{\ell}_{\theta^*,\hat{g}}(Z)) + \tilde{\cal E}(\hat{g})\\
    \leq~& 
    (P - P_n)(\tilde{\ell}_{\hat{\theta}, \hat{g}}(Z) - \tilde{\ell}_{\theta^*, \hat{g}}(Z)) + \tilde{\cal E}(\hat{g}) 
    {-\frac{\beta}{n}(K(\hat{\theta})-K(\theta^*))}
    \tag{By optimality of $\hat{\theta}$, $P_n(\tilde{\ell}_{\hat{\theta}, \hat{g}}(Z) - \tilde{\ell}_{\theta^*, \hat{g}}(Z)) {+ \frac{\beta}{n}(K(\hat{\theta})-K(\theta^*))}\leq 0$}\\
    =~& \mu\sum_j \hat{\theta}_j \|f_j - f_{\theta^*}\|^2 + \mu \hat{\theta}^\top H\theta^*+ Z_{n,1} + Z_{n,2} + \tilde{\cal E}(\hat{g})\\
    ~& {-\frac{\beta}{n}(K(\hat{\theta})-K(\theta^*)) + \frac{2}{s}K(\hat{\theta})}.
    \tag{definition of $Z_{n,1}, Z_{n,2}$}
\end{align*}
By simple algebraic manipulations note that for any $\theta\in \Theta$, we have:
\begin{align*}
    \sum_j \hat{\theta}_j \|f_j - f_{\theta^*}\|^2 = V(\hat{\theta}) + \|f_{\hat{\theta}} - f_{\theta^*}\|^2 \\
    \hat{\theta}^\top H \theta^* = V(\hat{\theta}) + V(\theta^*) + \|f_{\hat{\theta}} - f_{\theta^*}\|^2.
\end{align*}
Letting $Z_n=Z_{n,1} + Z_{n,2}$, we have:
\begin{align*}
    \tilde{R}(\hat{\theta}, g_0) - \tilde{R}(\theta^*, g_0) \leq~& 2\mu V(\hat{\theta}) + \mu V(\theta^*) + 2\mu \|f_{\hat{\theta}} - f_{\theta^*}\|^2 + Z_n+ \tilde{\cal E}(\hat{g}) {-\frac{\beta}{n}(K(\hat{\theta})-K(\theta^*)) + \frac{2}{s}K(\hat{\theta})}.
    % \leq~& 2\mu V(\hat{\theta}) + \mu V(\theta^*) + 2\mu \|f_{\hat{\theta}} - f_{\theta^*}\|^2 + Z_n + \error(\hat{g})\, \|f_{\hat{\theta}} - f_{\theta^*}\|^{\gamma}\\
    % ~& {-\frac{\beta}{n}(K(\hat{\theta})-K(\theta^*)) + \frac{2}{s}K(\hat{\theta})}
\end{align*}
Since we have
\begin{align*}
    \tilde{\cal E}(\hat{g}) :=~& \E(\tilde{\ell}_{\hat{\theta},g_0}(Z) - \tilde{\ell}_{\theta^*,g_0}(Z)) - \E(\tilde{\ell}_{\hat{\theta},\hat{g}}(Z) - \tilde{\ell}_{\theta^*,\hat{g}}(Z)) \\
    \leq~& \left((2C_Y + C_f)\|\hat{f}-f^*\|_{2}+C_f\|\Tcal(\hat{f}-f^*)\|_2 + 2C_f^2\right)\error(\hat{g}) +  2C_f^2\error(\hat{g})\\
    \leq~& \mu\left(\|f_{\hat{\theta}}-f_{\theta^*}\|_{2}^2+\|\Tcal(f_{\hat{\theta}}-f_{\theta^*})\|_2^2\right)+\underbrace{\frac{1}{2\mu}\left((2C_Y+C_f)^2+C_f^2\right)\error(\hat{g})^2 + 2C_f^2\error(\hat{g})}_{{\cal A}(\hat{g})} \tag{By Young's inequality}.
\end{align*}
% By Young's inequality, for any $a\geq 0$ and $b\geq 0$, we have $a\,b\leq \frac{a^p}{p} + \frac{b^q}{q}$, such that $\frac{1}{p}+\frac{1}{q}=1$. Taking $a= \left(\frac{\gamma}{2\mu}\right)^{\gamma/2} \error(\hat{g})$ and $b=\left(\frac{2\mu}{\gamma}\right)^{\gamma/2}\|f_{\hat{\theta}}-f_{\theta^*}\|^\gamma$ and $q=2/\gamma$ and $p=\frac{1}{1-1/q}=\frac{1}{1-\gamma/2}=\frac{2}{2-\gamma}$
% \begin{align}
%     \error(\hat{g})\, \|f_{\hat{\theta}} - f_{\theta^*}\|^{\gamma} \leq \underbrace{\frac{2}{2-\gamma}\left(\frac{\gamma}{2\mu}\right)^{\frac{\gamma}{2-\gamma}} \error(\hat{g})^{\frac{2}{2-\gamma}}}_{{\cal A}(\hat{g})}+ \mu \|f_{\hat{\theta}} - f_{\theta^*}\|^2
% \end{align}
Thus we can derive:
\begin{align}
    \tilde{R}(\hat{\theta}, g_0) - \tilde{R}(\theta^*, g_0) \leq~& 2\mu V(\hat{\theta}) + \mu V(\theta^*) + 3\mu \|f_{\hat{\theta}} - f_{\theta^*}\|^2 + \mu\|\Tcal(f_{\hat{\theta}}-f_{\theta^*})\|^2\\\label{eqn:first-upper}
    ~&  + Z_n + {\cal A}(\hat{g})  -\frac{\beta}{n}(K(\hat{\theta})-K(\theta^*)) + \frac{2}{s}K(\hat{\theta}).
\end{align}
By Lemma~\ref{lem:strong-conv-tilde}, we have:
\begin{align*}
    \tilde{R}(\hat{\theta}, g_0) - \tilde{R}(\theta^*, g_0) \geq~& \mu (V(\theta^*) - V(\hat{\theta})) {+\frac{\beta}{n}(K(\theta^*)-K(\hat{\theta}))} +(1-\nu-\mu)\|\Tcal(f_{\hat{\theta}} - f_{\theta^*})\|^2\\
    ~&+ \left(\alpha(1-\nu) - \mu\right)\|f_{\hat{\theta}} - f_{\theta^*}\|^2 \\
    \geq~& \mu (V(\theta^*) - V(\hat{\theta})) {+\frac{\beta}{n}(K(\theta^*)-K(\hat{\theta}))} + 6\mu  \|f_{\hat{\theta}} - f_{\theta^*}\|^2 \\
    ~&+ 2\mu\|\Tcal(f_{\hat{\theta}} - f_{\theta^*})\|^2
    \tag{Since: $\mu \leq \min\{\frac{\alpha(1-\nu)}{7},\frac{1-\nu}{3}\}$}
\end{align*}
Combining the latter upper and lower bound on $\tilde{R}(\hat{\theta}, g_0) - \tilde{R}(\theta^*, g_0)$, we have:
\begin{align*}
    3\mu \|f_{\hat{\theta}} - f_{\theta^*}\|^2 +\mu \|\Tcal(f_{\hat{\theta}} - f_{\theta^*})\|^2\leq 3\mu V(\hat{\theta}) + Z_n {+ \frac{2}{s} K(\hat{\theta})} + {\cal A}(\hat{g}).
\end{align*}
Plugging the bound on $3\mu\|f_{\hat{\theta}} - f_{\theta^*}\|^2$ back into Equation~\eqref{eqn:first-upper} we have:
\begin{align*}
    \tilde{R}(\hat{\theta}, g_0) - \tilde{R}(\theta^*, g_0) \leq ~& 5\mu V(\hat{\theta}) + \mu V(\theta^*) + 2\, Z_n + 2\,{\cal A}(\hat{g}) \\
    ~& {+\frac{\beta}{n}(K(\theta^*)-K(\hat{\theta})) + \frac{4}{s}K(\hat{\theta})}.
\end{align*}
By Lemma~\ref{lem:v-versus-tilde}, we have:
\begin{align*}
    R(\hat{\theta}, g_0) \leq~& \tilde{R}(\theta^*, g_0) + \mu V(\theta^*)+ \left(5\mu - \alpha\nu\right) V(\hat{\theta}) + 2 Z_n + 2{\cal A}(\hat{g}) \\
    ~& {+\frac{\beta}{n}K(\theta^*) + \left(\frac{4}{s} - \frac{\beta}{n}\right)K(\hat{\theta})}\\
    \leq~& \tilde{R}(\theta^*, g_0) + \mu V(\theta^*) {+ \frac{\beta}{n}K(\theta^*)} + 2 Z_n + 2{\cal A}(\hat{g})   \tag{since $\mu\leq \frac{\nu\alpha}{5}$,  $V(\hat{\theta})>0$, {$\beta \geq \frac{4n}{s}$ and $K(\hat{\theta}\geq 0 )$}}\\
    =~& \min_{\theta\in \Theta} [\tilde{R}(\theta, g_0) + \mu V(\theta){+ \frac{\beta}{n}K(\theta^*)}] + 2 Z_n + 2{\cal A}(\hat{g}) \tag{Since $\theta_*$ minimizes $H(\theta, g_0)$}\\
    \leq~& \min_{j} [\tilde{R}(e_j, g_0) + \mu V(e_{j}) {+ \frac{\beta}{n}\log(1/\pi_j)}]+ 2 Z_n + 2{\cal A}(\hat{g}) \\
    =~& \min_{j} [R(e_j, g_0){+ \frac{\beta}{n} \log(1/\pi_j)}] + 2 Z_n +2{\cal A}(\hat{g}). \tag{Since $\tilde{R}(e_j, g_0)=R(e_j, g_0)$ and $V(e_j)=0$ for all $e_j$}
\end{align*}
\end{proof}

\begin{lemma}\label{lem:z1}
For any positive $s\leq \min\left\{ \frac{\mu n}{32C_{Y,f}(1-\nu)^2}, \frac{n}{4\sqrt{6} C_f (1-\nu) \sqrt{C_{Y,f}}}\right\}$, where $C_{Y,f} = ((C_Y+C_f)^2+\alpha^2 C_f^2)$:
\begin{align}
    {\Pr\left(Z_{n,1} \geq \frac{\log(1/\delta)}{s}\right) \leq \delta}
\end{align}
\end{lemma}
\begin{proof}
By Lemma~\ref{lem:exp-tail} it suffices to show that:
\begin{align*}
    {I:=\E\left[\exp\left\{s Z_{n,1}\right\}\right] \leq 1}.
\end{align*}
% First note that:
% \begin{align}
%     I = \E\left[\exp\left\{s Z_{n,1}\} \exp\{- \log(M)\right\}\right] = \frac{1}{M} \E\left[\exp\left\{s Z_{n,1}\right\}\right]
% \end{align}
First, by the mean value theorem, we observe that:
\begin{align*}
    |\ell_{\hat{\theta},\hat{g}}(Z,X) - \ell_{\hat{\theta},\hat{g}}(Z,X)|\leq 2(C_Y+C_f)|\hat{\Tcal}(f_{\hat{\theta}}-f_{\theta^*})| + 2\alpha C_f|f_{\hat{\theta}}-f_{\theta^*}|.
\end{align*}
Define:
\begin{align*}
    f_{\ell,\theta}(X,Z) = 2(C_Y+C_f)\hat{\Tcal}(f_{\theta}) + 2\alpha C_f f_{\theta}.
\end{align*}
By the symmetrization argument:
\begin{align*}
    I =~& \E\left[\exp\left\{s \left((1-\nu)(P-P_n)(\ell_{\hat{\theta}}(Z,X) - \ell_{\theta^*}(Z,X)) - \mu \sum_j \hat{\theta}_j \|f_j - f_{\theta^*}\|^2 {-\frac{1}{s}\sum_j\hat{\theta}_j \log(1/\pi_j)}\right)\right\}\right]\\
    \leq~& \E\left[\exp\left\{s \max_{\theta} \left((1-\nu)(P-P_n)(\ell_{\hat{\theta}}(Z,X) - \ell_{\theta^*}(Z,X)) - \mu \sum_j \theta_j \|f_j - f_{\theta^*}\|^2 {+\frac{1}{s}\sum_j\theta_j \log(\pi_j)} \right)\right\}\right]\\
    \leq~& \E\left[\exp\left\{s \max_{\theta} \left(2\,(1-\nu)P_{n,\epsilon}(\ell_{\hat{\theta}}(Z,X) - \ell_{\theta^*}(Z,X)) - \mu \sum_j \theta_j \|f_j - f_{\theta^*}\|^2 {+\frac{1}{s}\sum_j\theta_j\log(\pi_j)}\right)\right\}\right].\tag{By symmetrization (Lemma \ref{sym})}
    % \leq~& \E\left[\exp\left\{s \max_{\theta} \left(2\,(1-\nu) P_{n,\epsilon}(f_{\ell,\hat{\theta}}(X,Z) - f_{\ell,\theta^*}(X,Z)) - \mu \sum_j \theta_j \|f_j - f_{\theta^*}\|^2 +\frac{1}{s}\sum_j\theta_j \log(\pi_j)\right)\right\}\right] \tag{contr.}
\end{align*}
By a contraction argument (Lemma \ref{contr}), we have that:
\begin{align*}
    P_{n,\epsilon}(\ell_{\hat{\theta}}(Z,X) - \ell_{\theta^*}(Z,X)) =~& P_{n,\epsilon}\left((Y - \Tcal(f_{\theta}))^2 - (Y - \Tcal(f_{\theta^*}))^2 + \alpha  f_{\theta}^2 - \alpha f_{\theta^*}^2\right)\\
    \leq~& P_{n,\epsilon}\left(2(C_Y+C_f)\Tcal(f_{\theta}) - \Tcal(f_{\theta}^*) + \alpha f_{\theta}-f_{\theta^*}\right)\\
    =~&P_{n,\epsilon}(f_{\ell,\hat{\theta}}(X,Z) - f_{\ell,\theta^*}(X,Z)).
\end{align*}
Then, we get:
\begin{align*}
    I\leq~& \E\left[\exp\left\{s \max_{\theta} \left(2\,(1-\nu) P_{n,\epsilon}(f_{\ell,\hat{\theta}}(X,Z) - f_{\ell,\theta^*}(X,Z)) - \mu \sum_j \theta_j \|f_j - f_{\theta^*}\|^2 +\frac{1}{s}\sum_j\theta_j \log(\pi_j)\right)\right\}\right]
\end{align*}

Since the quantity:
\begin{align*}
   2\,(1-\nu) P_{n,\epsilon}(f_{\ell,\hat{\theta}}(X,Z) - f_{\ell,\theta^*}(X,Z)) - \mu \sum_j \theta_j \|f_j - f_{\theta^*}\|^2 +\frac{1}{s}\sum_j\theta_j \log(\pi_j)
\end{align*}
is a linear function of $\theta$ (recall that $f_{\theta}=\sum_j \theta_j f_j$), the maximum over $\theta$ is attained at one of $e_1, \ldots, e_M$. Thus:
\begin{align*}
    I \leq ~&\E\left[\exp\left\{s \max_{j} \left(2\,(1-\nu) P_{n,\epsilon}(f_{\ell,\hat{\theta}}(X,Z) - f_{\ell,\theta^*}(X,Z)) - \mu \|f_j - f_{\theta^*}\|^2{+\frac{1}{s}\log(\pi_j)}\right)\right\}\right]\\
    \leq~& \E\left[\max_{j}{\pi_j} \exp\left\{s \left(2\,(1-\nu) P_{n,\epsilon}(f_{\ell,\hat{\theta}}(X,Z) - f_{\ell,\theta^*}(X,Z)) - \mu \|f_j - f_{\theta^*}\|^2\right)\right\}\right]\\
    \leq~& \E\left[\sum_{j} {\pi_j} \exp\left\{s \left(2\,(1-\nu) P_{n,\epsilon}(f_{\ell,\hat{\theta}}(X,Z) - f_{\ell,\theta^*}(X,Z)) - \mu \|f_j - f_{\theta^*}\|^2\right)\right\}\right]\\
    =~&\sum_{j} {\pi_j} \E\left[ \exp\left\{s \left(2\,(1-\nu) P_{n,\epsilon}(f_{\ell,\hat{\theta}}(X,Z) - f_{\ell,\theta^*}(X,Z)) - \mu \|f_j - f_{\theta^*}\|^2\right)\right\}\right]\\
    =~& \sum_{j} {\pi_j}\E_{Z_{1:n},X_{1:n}}\E_\epsilon\left[ \exp\left\{s \left(2\,(1-\nu) P_{n,\epsilon}(f_{\ell,\hat{\theta}}(X,Z) - f_{\ell,\theta^*}(X,Z)) - \mu \|f_j - f_{\theta^*}\|^2\right)\right\}\right]\\
    =~& \sum_{j} {\pi_j}\E_{Z_{1:n},X_{1:n}}\left[ \E_\epsilon\left[ \exp\left\{s 2\,(1-\nu) P_{n,\epsilon}(f_{\ell,\hat{\theta}}(X,Z) - f_{\ell,\theta^*}(X,Z))\right\}\right] \exp\left\{- s \mu \|f_j - f_{\theta^*}\|^2\right\}\right].
\end{align*}
By Hoeffding's inequality in Lemma~\ref{lem:hoeffding}, we can bound the first term in the product:
\begin{align*}
    ~&\E_\epsilon\left[ \exp\left\{s 2\,(1-\nu) P_{n,\epsilon}(f_{\ell,\hat{\theta}}(X,Z) - f_{\ell,\theta^*}(X,Z))\right\}\right]\\
    ~&\leq \exp\left\{\frac{2 s^2}{n}\,(1-\nu)^2 P_n(f_{\ell,\hat{\theta}}(X,Z) - f_{\ell,\theta^*}(X,Z))^2\right\}\\
    ~&\leq \exp\left\{\frac{16 s^2}{n}\,(1-\nu)^2 P_n\left[(C_Y+C_f)^2(\Tcal(f_{\hat{\theta}})(Z)-\Tcal(f_{\theta^*}(Z))^2 + \alpha^2 C_f^2(f_{\hat{\theta}}(X)-f_{\theta^*}(X))^2\right]\right\}\\
    ~&\leq \exp\left\{\frac{16 s^2}{n}\,(1-\nu)^2 P_n((C_Y+C_f)^2+\alpha^2 C_f^2)(f_{\hat{\theta}}(X)-f_{\theta^*}(X))^2 \right\}.
\end{align*}
Let $C_{Y,f}:=((C_Y+C_f)^2+\alpha^2 C_f^2)$, then we have:
\begin{align*}
    I\leq~&\sum_{j} {\pi_j}\E_{X_{1:n}}\left[ \exp\left\{\frac{16s^2}{n}\,(1-\nu)^2C_{Y,f} P_n(f_{j}(X) - f_{\theta^*}(X))^2\right\}\, \exp\left\{- s \mu \|f_j - f_{\theta^*}\|^2\right\}\right]\\
    =~& \sum_{j} {\pi_j}\E_{X_{1:n}}\left[ \exp\left\{\frac{16s^2}{n}\,(1-\nu)^2 C_{Y,f}  P_n(f_{j}(X) - f_{\theta^*}(X))^2\right\}\, \exp\left\{- s \mu P (f_j(X) - f_{\theta^*}(X))^2\right\}\right]\\
    =~& \sum_{j} {\pi_j} \E_{X_{1:n}}\left[ \exp\left\{\frac{16s^2\,(1-\nu)^2C_{Y,f} }{n} \left(P_n - \frac{\mu\, n}{2 s(1-\nu)^2 C_{Y,f} }P\right)(f_{j}(X) - f_{\theta^*}(X))^2\right\}\right].
\end{align*}
Since $s$ is small enough such that: $\frac{\mu\, n}{16 s (1-\nu)^2 C_{Y,f}}\geq 2$ (which is satisfied by our assumptions on $s$), we have that:
\begin{align*}
    I \leq~& \sum_{j} \pi_j \E_{Z_{1:n}}\left[ \exp\left\{\frac{16s^2\,(1-\nu)^2 C_{Y,f}}{n} \left(P_n - 2P\right)(f_{j}(Z) - f_{\theta^*}(Z))^2\right\}\right]\\
    =~& \sum_{j} \pi_j\E_{Z_{1:n}}\left[ \exp\left\{\frac{16s^2\,(1-\nu)^2 C_{Y,f}}{n} \left(\left(P_n - P\right)(f_{j}(Z) - f_{\theta^*}(Z))^2 - P (f_{j}(Z) - f_{\theta^*}(Z))^2 \right)\right\}\right].
\end{align*}
Since $f_j(Z)\in [-C_f, C_f]$, we have that $f_j(X) - f_{\theta^*}(X)\in [-2C_f, 2C_f]$. Thus:
\begin{align*}
(f_j(X)-f_{\theta^*}(X))^2\geq \frac{(f_j(X)-f_{\theta^*}(X))^2}{4C_f^2} (f_j(X)-f_{\theta^*}(X))^2 = \frac{1}{4C_f^2} (f_j(X)-f_{\theta^*}(X))^4.
\end{align*}
Thus:
\begin{align*}
    I \leq~& \sum_{j} \pi_j \E_{Z_{1:n}}\left[ \exp\left\{\frac{16s^2\,(1-\nu)^2 C_{Y,f}}{n} \left(\left(P_n - P\right)(f_{j}(Z) - f_{\theta^*}(Z))^2 - \frac{1}{4C_f^2} P (f_{j}(Z) - f_{\theta^*}(Z))^4 \right)\right\}\right].
\end{align*}
It suffices to show that each of the summands is upper bounded by $1$. Then their weighted average is also upper bounded by $1$.
Applying Lemma~\ref{lem:bernstein} with:
\begin{align*}
    V_i =~& (f_j(Z_i) - f_{\theta^*}(Z_i))^2\nonumber\\
    c =~& 4C_f^2\nonumber\\
    \lambda =~& \frac{16s^2(1-\nu)^2C_{Y,f}}{n^2}\nonumber\\
    c_0 =~& \frac{1}{4C_f^2}.
\end{align*}
we have that as long as:
\begin{align*}
    \lambda := \frac{16s^2(1-\nu)^2C_{Y,f}}{n^2} \leq \frac{2c_0}{1+2c_0 c} =: \frac{2\cdot 1/4C_f^2}{1 + 2} = \frac{1}{6C_f^2}.
\end{align*}
Then:
\begin{align*}
    \E_{Z_{1:n}}\left[ \exp\left\{\frac{16s^2\,(1-\nu)^2 C_{Y,f}}{n} \left(\left(P_n - P\right)(f_{j}(Z) - f_{\theta^*}(Z))^2 - \frac{1}{4C_f^2} P (f_{j}(Z) - f_{\theta^*}(Z))^4 \right)\right\}\right]\leq 1
\end{align*}
Thus we need that:
\begin{align*}
    s \leq \frac{n}{4\sqrt{6} C_f (1-\nu) \sqrt{C_{Y,f}}},
\end{align*}
which is satisfied by our assumptions on $s$.
\end{proof}

\begin{lemma}\label{lem:z2}
For any positive $s\leq \frac{2 \mu n}{\nu (\nu C_{\ell,f} + 2 \mu C_{\ell})}$ where $C_{\ell,f} = 8((C_Y+C_f)^2+\alpha^2C_f^2)$ and $C_{\ell} = (C_Y+C_f)^2+\alpha C_f^2$:
\begin{align*}
    \Pr\left(Z_{n,2} \geq \frac{\log(1/\delta)}{s}\right) \leq \delta.
\end{align*}
\end{lemma}
\begin{proof}
It suffices to show that:
\begin{align*}
    I := \E\left[\exp\left\{s Z_{n,2})\right\}\right] \leq 1.
\end{align*}
% First note that:
% \begin{align}
%     I = \E\left[\exp\left\{s Z_{n,2}\right\} \exp\left\{-\log(M)\right\}\right] = \frac{1}{M}  \E\left[\exp\left\{s Z_{n,2}\right\}\right]
% \end{align}
By definition of $Z_{n,2}$:
\begin{align*}
    I =~& \E\left[\exp\left\{s \left( \nu(P-P_n) \sum_{j} (\hat{\theta}_j-\theta_j^*)\ell_{e_j}(Z) - \mu \hat{\theta}^\top H \theta^*  -\frac{1}{s}\sum_j\hat{\theta}_j \log(1/\pi_j) \right)\right\}\right]\\
    =~& \E\left[\exp\left\{s \left( \nu(P-P_n) \sum_{j} \hat{\theta}_j \ell_{e_j}(Z) + \frac{1}{s}\sum_j\hat{\theta}_j \log(\pi_j)\right)\right\}\right.\\
    ~&\quad \left.\exp\left\{s\left( - \nu (P-P_n)\sum_{j} \theta_j^*\ell_{e_j}(Z) - \mu \hat{\theta}^\top H \theta^* \right)\right\}\right].
\end{align*}
Note that:
\begin{align*}
    - \nu (P-P_n)\sum_{j} \theta_j^*\ell_{e_j}(Z) - \mu \hat{\theta}^\top H \theta^* =~& \sum_{j} \theta_j^* \left( -\nu (P-P_n) \ell_{e_j}(Z) - \mu \hat{\theta}^\top H_{\cdot, j}\right) \\
    =~& \sum_{j} \theta_j^* \underbrace{\left( -\nu (P-P_n) \ell_{e_j}(Z) - \mu \sum_{k} \hat{\theta}_k \|f_k - f_j\|^2\right)}_{\Phi_j}.
\end{align*}
Thus:
\begin{align*}
    I=~&  \E\left[\exp\left\{s \left( \nu(P-P_n) \sum_{j} \hat{\theta}_j \ell_{e_j}(Z)  +\frac{1}{s}\sum_j\hat{\theta}_j \log(\pi_j)\right)\right\} \exp\left\{\sum_{j}\theta_j^* s\Phi_j\right\}\right].
\end{align*}
By Jensen's inequality:
\begin{align*}
    \exp\left\{\sum_{j}\theta_j^* s\Phi_j\right\} \leq \sum_{j}\theta_j^* \exp\left\{s\Phi_j\right\}.
\end{align*}
Thus:
\begin{align*}
    I\leq~& \sum_{j}\theta_j^*  \E\left[\exp\left\{s \left( \nu(P-P_n) \sum_{k} \hat{\theta}_k \ell_{e_k}(Z){+\frac{1}{s}\sum_k \hat{\theta}_k \log(\pi_k)}\right)\right\} \exp\left\{ s\Phi_j\right\}\right]\\
    =~& \sum_{j}\theta_j^* \E\left[\exp\left\{s \left( \nu(P-P_n) \sum_{k} \hat{\theta}_k \ell_{e_k}(Z) + \Phi_j {+\frac{1}{s}\sum_k\hat{\theta}_k \log(\pi_k)}\right)\right\}\right]\\
    =~& \sum_{j}\theta_j^*  \E\left[\exp\left\{s \left( \nu(P-P_n) \left(\sum_{k} \hat{\theta}_k \ell_{e_k}(Z) - \ell_{e_j}(Z)\right) - \mu \sum_{k} \hat{\theta}_k \|f_k - f_j\|^2{+\frac{1}{s}\sum_k\hat{\theta}_k \log(\pi_k)}\right)\right\}\right]\\
    =~& \sum_{j}\theta_j^*  \E\left[\exp\left\{s \left( \nu(P-P_n) \sum_{k} \hat{\theta}_k \left(\ell_{e_k}(Z) - \ell_{e_j}(Z)\right) - \mu \sum_{k} \hat{\theta}_k \|f_k - f_j\|^2 + \frac{1}{s}\sum_k\hat{\theta}_k \log(\pi_k)\right)\right\}\right]\\
    \leq~& \sum_{j}\theta_j^* \E\left[\exp\left\{s \max_{\theta} \left( \nu(P-P_n) \sum_{k} \theta_k \left(\ell_{e_k}(Z) - \ell_{e_j}(Z)\right) - \mu \sum_{k} \theta_k \|f_k - f_j\|^2 + \frac{1}{s}\sum_k\theta_k 
    \log(\pi_k)\right)\right\}\right].
\end{align*}
Since the objective in the exponent is linear in $\theta$ it takes its value at one of the extreme points $e_1, \ldots, e_M$:
\begin{align*}
    I \leq~& \sum_{j}\theta_j^*  \E\left[\exp\left\{s \max_{k} \left( \nu(P-P_n) \left(\ell_{e_k}(Z) - \ell_{e_j}(Z)\right) - \mu \|f_k - f_j\|^2+\frac{1}{s} \log(\pi_k)\right)\right\}\right]\\
    =~& \sum_{j}\theta_j^*  \E\left[\max_{k}{\pi_k}\exp\left\{s  \left( \nu(P-P_n) \left(\ell_{e_k}(Z) - \ell_{e_j}(Z)\right) - \mu \|f_k - f_j\|^2\right)\right\}\right]\\
    \leq~& \sum_{j}\theta_j^* \E\left[\sum_{k}{\pi_k} \exp\left\{s  \left( \nu(P-P_n) \left(\ell_{e_k}(Z) - \ell_{e_j}(Z)\right) - \mu \|f_k - f_j\|^2\right)\right\}\right]\\
    \leq~& \sum_{j}\theta_j^* \sum_{k} {\pi_k}\E\left[ \exp\left\{s  \left( \nu(P-P_n) \left(\ell_{e_k}(Z) - \ell_{e_j}(Z)\right) - \mu \|f_k - f_j\|^2\right)\right\}\right].
\end{align*}
By a the mean-value theorem, we have:
\begin{align*}
    |\ell_{e_k}(Z,X) - \ell_{e_j}(Z,X)|\leq 2(C_Y+C_f)|\hat{\Tcal}(f_{k}-f_{j})| + 2\alpha C_f|f_{\hat{\theta}}-f_{\theta^*}|.
\end{align*}
Thus,
\begin{align*}
     \|\ell_{e_k} - \ell_{e_j}\|^2 \leq~& 8(C_Y+C_f)^2\|\hat{\Tcal}(f_{k}-f_{j})\|^2 + 8\alpha^2 C_f^2 \|f_k - f_j\|^2\\
     \leq~& 8\left((C_Y+C_f)^2+\alpha^2C_f^2\right) \|f_k - f_j\|^2.
\end{align*}
% By Lipschitzness of $\ell$, we have that:
% \begin{align*}
%     \|\ell_{e_k} - \ell_{e_j}\|^2 \leq C_b^2 \|f_k - f_j\|^2
% \end{align*}
Let $C_{\ell,f} = 8\left((C_Y+C_f)^2+\alpha^2C_f^2\right)$, thus:
\begin{align*}
    I \leq~& \sum_{j}\theta_j^* \sum_{k} {\pi_k}\E\left[ \exp\left\{s  \left( \nu(P-P_n) \left(\ell_{e_k}(Z) - \ell_{e_j}(Z)\right) - \frac{\mu}{C_{\ell,f}} \|\ell_{e_k} - \ell_{e_j}\|^2\right)\right\}\right] \\
    =~& \sum_{j}\theta_j^* \sum_{k} {\pi_k}\E\left[ \exp\left\{s  \left( \nu(P-P_n) \left(\ell_{e_k}(Z) - \ell_{e_j}(Z)\right) - \frac{\mu}{C_{\ell,f}} P(\ell_{e_k}(Z) - \ell_{e_j}(Z))^2\right)\right\}\right]\\
    =~& \sum_{j}\theta_j^* \sum_{k}{\pi_k} \E\left[ \exp\left\{s \nu \left((P-\PP_n) \left(\ell_{e_k}(Z) - \ell_{e_j}(Z)\right) - \frac{\mu}{\nu\, C_{\ell,f}} P(\ell_{e_k}(Z) - \ell_{e_j}(Z))^2\right)\right\}\right].
\end{align*}
Invoking Lemma~\ref{lem:bernstein} with:
\begin{align*}
    V_i =~& \ell_{e_j}(Z_i) - \ell_{e_k}(Z_i)\\
    c_0 =~& \frac{\mu}{\nu C_{\ell,f}}\\
    \lambda =~& \frac{s \nu}{n}\\
    c=~& C_{\ell} :=(C_Y+C_f)^2+\alpha C_f^2
\end{align*}
Hence we need:
\begin{align*}
    ~&\lambda := \frac{s\nu}{n} \leq \frac{2c_0}{1 + 2c_0c} =: \frac{2 \mu}{\nu C_{\ell,f} (1 + 2 \frac{\mu}{\nu C_{\ell,f}} C_{\ell})} = \frac{2 \mu}{(\nu C_{\ell,f} + 2 \mu C_{\ell})} \implies s \leq \frac{2 \mu n}{\nu (\nu C_{\ell,f} + 2 \mu C_{\ell})}
\end{align*}
\end{proof}


% \subsubsection{Proof of Theorem \ref{thm:q_agg} }\label{proof_q_agg}
% We introduce some population quantities that will play a key role in our analysis. Define:
% \begin{align*}
%     V(\theta) :=~& \sum_{j} \theta_j \|f_j - f_{\theta}\|^2 
% \end{align*}
% $$    H(\theta, g) := \tilde{R}(\theta, g) + \mu V(\theta) 
%     + \frac{\beta}{n} K(\theta)$$
% and define:
% \begin{align*}
%     \theta^* =~& \arg\min_{\theta} \tilde{R}(\theta, g_0) + \mu V(\theta) + \frac{\beta}{n} K(\theta) 
%     ~=~ \arg\min_{\theta} H(\theta, g_0) \\
% \end{align*}
% Then we introduce the shorthand notation and invoke the error decomposition assumption:
% \begin{align*}
%     \tilde{{\cal E}}(\hat{g}) :=~& (\tilde{\ell}_{\hat{\theta},g_0}(Z) - \tilde{\ell}_{\theta^*,g_0}(Z)) - (\tilde{\ell}_{\hat{\theta},\hat{g}}(Z) - \tilde{\ell}_{\theta^*,\hat{g}}(Z))\\
%     {\cal E}(\hat{g}) :=~& ({\ell}_{\hat{\theta},g_0}(Z) - {\ell}_{\theta^*,g_0}(Z)) - ({\ell}_{\hat{\theta},\hat{g}}(Z) - {\ell}_{\theta^*,\hat{g}}(Z))
%     % ~\leq~ \error(\hat{g})\, \|f_{\hat{\theta}} - f_{\theta^*}\|^{\gamma}
% \end{align*}

% Suppose that $|Y|$ is almost surely bounded by some constant $C_Y$, and $|f(X)|$ is almost surely uniformly bounded by a constant $C_f$.
% \hlan{This is the diff of diff of the original loss, not the modified loss.}
% \begin{align*}
%     {\cal E}(\hat{g})
%     =~& \E [-2Y (\Tcal-\hat{\Tcal})(\hat{f}-f^*) + (\Tcal(\hat{f}))^2-(\Tcal(f^*))^2 -(\hat{\Tcal}(\hat{f}))^2+(\hat{\Tcal}(f^*))^2]\\
%     =~&\E[ -2Y (\Tcal-\hat{\Tcal})(\hat{f}-f^*) + (\Tcal-\hat{\Tcal})(\hat{f}-f^*)\hat{\Tcal}(\hat{f}+f^*)]\\
%     ~&+ \E[\Tcal(\hat{f}-f^*)(\Tcal - \hat{\Tcal})(\hat{f}+f^*)]\\
%     \leq~& \E[2|C_Y-C_f||(\Tcal-\hat{\Tcal})(\hat{f}-f^*)| + | \Tcal(\hat{f}-f^*)(\Tcal-\hat{\Tcal})(\hat{f}+f^*)|]
%     % \leq~&\mathbb{E} \left[2C_Y \left(\int (f-f')^2g_0dx\right)^{\frac{1}{2}}\left(\int\left(\frac{g}{g_0}-1\right)^2g_0dx\right)^{\frac{1}{2}}\right] \tag{Here we need the true likelihood to be bounded away from $0$} \\
%     % \leq~&\mathbb{E} \left\{2C_Y \mathbb{E}[(f-f')^2|Z]^{\frac{1}{2}}\mathcal{\chi}^2(g(X|Z),g_0(X|Z))^{\frac{1}{2}}\right\}\\
%     % \leq~& 2C_Y \mathbb{E} \left[\mathbb{E}[(f-f')^2|Z] \right]^{\frac{1}{2}}\mathbb{E} \left[\mathcal{\chi}^2(g(X|Z),g_0(X|Z))\right]^{\frac{1}{2}}\\
%     % =~&  2C_Y\|f-f'\|_{2} \mathbb{E} \left[\mathcal{\chi}^2(g(X|Z),g_0(X|Z))\right]^{\frac{1}{2}}\\
% \end{align*}
% First, we analyze the $(\Tcal-\hat{\Tcal})(\hat{f}-f^*)$ term:
% \begin{align*}
%     \mathbb{E}[|(\Tcal-\hat{\Tcal})(\hat{f}-f^*)|]\leq~&\mathbb{E}\left[\int |(\hat{f}-f^*)(g-g_0)|dx\right]\\
%     \leq~&\mathbb{E}\left[\left(\int (\hat{f}-f^*)^2g_0dx\right)^{\frac{1}{2}}\left(\int\left(\frac{g}{g_0}-1\right)^2g_0dx\right)^{\frac{1}{2}}\right]\tag{Here we need the true likelihood to be bounded away from $0$}\\
%     =~& \mathbb{E}\left[\mathbb{E}[(\hat{f}-f^*)^2|Z]^{\frac{1}{2}}\mathcal{\chi}^2(g(X|Z),g_0(X|Z))^{\frac{1}{2}}\right]\\
%     \leq~& \mathbb{E}\left[\mathbb{E}[(\hat{f}-f^*)^2|Z] \right]^{\frac{1}{2}}\mathbb{E} \left[\mathcal{\chi}^2(g(X|Z),g_0(X|Z))\right]^{\frac{1}{2}}\\
%     =~& \|\hat{f}-f^*\|_{2} \mathbb{E} \left[\mathcal{\chi}^2(g(X|Z),g_0(X|Z))\right]^{\frac{1}{2}}\\
% \end{align*}

% Next, we analyze the $\Tcal(\hat{f}-f^*)(\Tcal-\hat{\Tcal})(\hat{f}+f^*)$ term:
% \begin{align*}
%     \mathbb{E}\left[| \Tcal(\hat{f}-f^*)(\Tcal-\hat{\Tcal})(\hat{f}+f^*)|\right] \leq~& 
%     \mathbb{E}\left[|\Tcal(\hat{f}-f^*)|\mathbb{E}[(\hat{f}+f^*)^2|Z]^{\frac{1}{2}}\mathcal{\chi}^2(g(X|Z),g_0(X|Z))^{\frac{1}{2}}\right]\\
%     \leq~& 2C_f\mathbb{E}\left[|\Tcal(\hat{f}-f^*)|\mathcal{\chi}^2(g(X|Z),g_0(X|Z))^{\frac{1}{2}}\right]\\
%     \leq~&2C_f\|\Tcal(\hat{f}-f^*)\|_2\mathbb{E}\left[\mathcal{\chi}^2(g(X|Z),g_0(X|Z))\right]^{\frac{1}{2}}
% \end{align*}

% Since $g_o(X|Z)\geq c_0$, we have:
% \begin{align*}
%     \mathcal{\chi}^2(g(X|Z),g_0(X|Z)) =~&\int\left(\frac{g}{g_0}-1\right)^2g_0dx\\
%     =~&\int\left(\frac{\sqrt{g}}{\sqrt{g_0}}+1\right)^2\left(\frac{\sqrt{g}}{\sqrt{g_0}}-1\right)^2g_0dx\\
%     \leq~&\int\left(\frac{1}{c_0}+1\right)^2\left(\frac{\sqrt{g}}{\sqrt{g_0}}-1\right)^2g_0dx\\
%     =~&\left(\frac{1}{c_0}+1\right)^2 H(g,g_0)^2\\
% \end{align*}

% Thus, we get:
% \begin{align*}
%     {\cal E}(\hat{g}) \leq~& \left(2|C_Y-C_f|\|\hat{f}-f^*\|_{2}+2C_f\|_2\Tcal(\hat{f}-f^*)\|\right)(\frac{1}{c_0} + 1)^2\mathbb{E}\left[H(g,g_0)^2\right]^{\frac{1}{2}}\\
%     \leq~& \left(|C_Y-C_f|\|\hat{f}-f^*\|_{2}+C_f\|\Tcal(\hat{f}-f^*)\|_2\right)\error(\hat{g})\\
% \end{align*} where $\error(\hat{g}) =2(\frac{1}{c_0} + 1)^2\mathbb{E}\left[H(\hat{g},g_0)^2\right]^{\frac{1}{2}}$.

% For $\tilde{\cal E}(\hat{g})$, we have that:
% \begin{align*}
%     \tilde{\cal E}(\hat{g}) =~& {\cal E}(\hat{g})+\E[-2Y(\Tcal -\hat{\Tcal})(\hat{f}-f^*) + \sum_{j} (\hat{\theta}_j-\theta_j^*)((\Tcal f_j)^2 -(\hat{\Tcal}f_j)^2)]\\
%     = ~& {\cal E}(\hat{g}) +\E[-2Y(\Tcal -\hat{\Tcal})(\hat{f}-f^*) + \sum_j (\hat{\theta}_j-\theta_j^*)( \Tcal f_j +\hat{\Tcal}f_j)(\Tcal f_j -\hat{\Tcal}f_j)]\\
%     % \leq~&  {\cal E}(\hat{g}) + \E[2C_Y|(\Tcal -\hat{\Tcal})(\hat{f}-f^*)| + \sum_j(\Tcal f_j +\hat{\Tcal}f_j)(\Tcal -\hat{\Tcal})(\hat{f}-f^*)]
% \end{align*}

% To be able to use Q-aggregation for model selection, we also need strong convexity of the loss function in expectation with respect to some distance $d(\cdot,\cdot)$, i.e.
% \begin{align*}
%     P\left(\ell_{f', g_0} - \ell_{f, g_0}\right)  - P D_{f}\ell_{f, g_0}[f'-f]\geq \frac{\sigma}{2}d(f,f')^2
% \end{align*} for some $\sigma$.

% Using the loss function described above:
% \begin{align*}
%     ~&P\left(\ell_{f', g_0} - \ell_{f, g_0}\right)  - P D_{f}\ell_{f, g_0}[f'-f] \\=~& \mathbb{E}\left[-2Y(\Tcal(f'-f)) + (\Tcal(f'))^2 - (\Tcal(f))^2 + 2(Y-\Tcal(f))(\Tcal(f'-f))\right]\\
%     ~&+ \alpha \mathbb{E}\left[(f')^2 - f^2 -2f(f'-f) \right]\\
%     =~& \mathbb{E}\left[(\Tcal(f') - \Tcal(f))^2\right] + \alpha \mathbb{E}[(f'-f)^2] \\
%     =~& \|\Tcal(f') - \Tcal(f)\|_{2}^{2} + \alpha \|f'-f\|_{2}^{2}
% \end{align*}

% \begin{lemma}\label{lem:prel-strong-conv} For any $\theta,\theta'\in \Theta$:
% \begin{align*}\label{eqn:strong-conv-R-lemma}
%     R(\theta, g)-R(\theta', g) \geq \langle \nabla R(\theta', g), \theta - \theta' \rangle + \|\Tcal(f_{\theta} - f_{\theta'})\|_2^2+\alpha\|f_{\theta}-f_{\theta'}\|^2
% \end{align*}
% \end{lemma}
% \begin{proof}
%     Follows from a similar calculation as above.
% \end{proof}

% \begin{lemma}\label{lem:v-versus-tilde} 
% For any $\theta\in \Theta$:
% \begin{equation}
%     \tilde{R}(\theta, g) \geq R(\theta, g) + \nu V(\theta)
% \end{equation}
% \end{lemma}
% \begin{proof}
% By Lemma~\ref{lem:prel-strong-conv}, we have that:
% \begin{align*}
%     R(e_j, g) - R(\theta, g) \geq \langle \nabla R(\theta, g), e_j - \theta\rangle +\alpha\|f_j - f_\theta\|^2
% \end{align*}
% Multiplying each equation by $\theta_j$ and summing across $j$, yields:
% \begin{align*}
%     R(\theta, g) \leq \sum_j \theta_j R(e_j, g) - \alpha V(\theta)
% \end{align*}
% where we used the fact that $\sum_{j}\theta_j=1$ and $\sum_{j} \theta_j \langle \nabla R(\theta, g), e_j - \theta\rangle=0$.

% Thus we have that:
% \begin{align*}
%     \tilde{R}(\theta, g) =~& (1-\nu) R(\theta, g) + \nu \sum_j \theta_j R(e_j, g)\\
%     \geq~& (1-\nu) R(\theta, g) + \nu R(\theta, g) +\nu\alpha V(\theta)\\
%     =~& R(\theta, g) + \nu\alpha V(\theta)
% \end{align*}
% \end{proof}
% \begin{lemma}\label{lem:strong-conv-tilde}
% Let $\theta_*=\argmin_{\theta\in \Theta} H(\theta, g_0)$, then for any $\theta\in \Theta$:
% \begin{align}
%     H(\theta, g_0) - H(\theta^*, g_0) \geq \left(1-\nu - \mu\right) \|\Tcal(f_{\theta} - f_{\theta^*})\|^2 + \alpha(1-\nu)\|f_{\theta} - f_{\theta^*}\|^2
% \end{align}
% \end{lemma}
% \begin{proof}
% By Lemma~\ref{lem:prel-strong-conv} we have:
% \begin{align}\label{eqn:strong-conv-R}
%     R(\theta, g_0)-R(\theta', g_0) \geq \langle \nabla R(\theta', g_0), \theta - \theta' \rangle + \|\Tcal(f_{\theta}-f_{\theta'})\|^2+\alpha\|f_{\theta}-f_{\theta'}\|^2
% \end{align}
% Now let: 
% \begin{align*}
%     q(Z;\theta, g) =~& \sum_{j=1}^{M}\theta_{j} \ell(Z;f_{e_{j}}(X), g)\\
%     L_e(g)=~&[\ell(Z;f_{e_{1}}(X),g)\cdots \ell(Z;f_{e_{M}}(X), g)]^{T}\\
%     Q(\theta, g) =~& \mathbb{E}\left[q(Z;\theta, g)\right]\\
% \end{align*}
% Then we have:
% \begin{align*}
%     q(Z;\theta, g)-q(Z; \theta', g) =~&\sum_{j=1}^{M}\ell(Z;f_{e_{j}}(X), g)(\theta_{j}-\theta'_{j})
%     = \langle{L_{e}(g), \theta-\theta'}\rangle\\
%     =~& \langle \nabla q(Z; \theta', g), \theta-\theta' \rangle
% \end{align*}
% By taking the expectation over $Z$:
% \begin{align}\label{eqn:strong-conv-G}
%     Q(\theta, g) - Q(\theta', g) = \sum_{j=1}^{M}\theta_{j}R(e_j, g)-\sum_{j=1}^{M}\theta_{j}'R(e_j, g)=\langle \nabla Q(\theta', g), \theta-\theta' \rangle
% \end{align}
% By adding Equations~\eqref{eqn:strong-conv-R} and \eqref{eqn:strong-conv-G} we have that: 
% \begin{align*}
%     \tilde{R}(\theta, g_0)-\tilde{R}(\theta', g_0) =~& (1 - \nu) (R(\theta, g_0)-R(\theta', g_0)) + \nu (Q(\theta, g_0) - Q(\theta', g_0))\\
%     \geq~& (1-\nu) \langle \nabla R(\theta',g_0), \theta - \theta' \rangle + \nu \langle \nabla Q(\theta',g_0), \theta-\theta' \rangle\\
%     ~&+ (1-\nu)\left(\|\Tcal(f_{\theta}-f_{\theta'})\|^2+\alpha\|f_{\theta}-f_{\theta'}\|^2\right) \\
%     \geq~& \langle \nabla \tilde{R}(\theta',g_0), \theta - \theta' \rangle+ (1-\nu)\left(\|\Tcal(f_{\theta}-f_{\theta'})\|^2+\alpha\|f_{\theta}-f_{\theta'}\|^2\right) 
% \end{align*}
%  We also have that for any $\theta, \theta'$:\footnote{Since the function $V(\theta)$ can be interpreted as the expected variance of $f_j(X)$ over the distribution of $j$. Thus with simple manipulations it can be shown to be equal to $V(\theta)=\sum_j \theta_j \|f_j\|^2- \|f_\theta\|^2$. Since this is a quadratic in $\theta$ function, with Hessian $\mathbb{E}[FF']$, where $F$ is the vector $(f_1(X), \ldots, f_j(X))$, the above equality follows by an exact second order Taylor expansion and the observation that $(\theta-\theta')^\top\mathbb{E}[FF'](\theta-\theta') = \|f_\theta - f_{\theta'}\|^2$}
% \begin{align*}
%     V(\theta)-V(\theta') = \langle \nabla V(\theta'), \theta - \theta' \rangle - \|f_{\theta}-f_{\theta'}\|_{2}^2
% \end{align*}
% which implies that $V(\theta)$ is at most 2-strongly concave. 
% We also have:
% \begin{align*}
%     K(\theta) - K(\theta') = \sum_j (\theta_j - \theta'_j) \log(1/\pi_j) = \langle \nabla K(\theta'), \theta - \theta' \rangle
% \end{align*}

% Adding the three inequalities and substituting $\theta'$ with $\theta^*$ we have:
% \begin{align*}
%         H(\theta, g_0) - H(\theta^*, g_0) \geq~& \langle \nabla H(\theta^*, g_0), \theta-\theta^* \rangle + \left(1-\nu \right) \|\Tcal(f_{\theta} - f_{\theta^*})\|^2\\
%         ~&+ (\alpha(1-\nu)-\mu)\|f_{\theta} - f_{\theta^*}\|^2
% \end{align*}
% Note that by the optimality of $\theta^*$ over the convex set $\Theta$ we have that $\langle \nabla H(\theta^*, g_0), \theta-\theta^* \rangle \geq 0$, which immediately yields the statement we want to prove.
% \end{proof}

% \begin{lemma}\label{lem:main}
% Define the centered and offset empirical processes:
% \begin{align*}
%     Z_{n,1} =~& (1-\nu)(P-P_n) (\ell_{\hat{\theta}, \hat{g}}(Z) - \ell_{\theta^*, \hat{g}}(Z)) - \mu \sum_{j} \hat{\theta}_j \|f_j - f_{\theta^*}\|^2 {-\frac{1}{s}K(\hat{\theta})}\\
%     Z_{n,2} =~& \nu(P-P_n) \sum_{j} (\hat{\theta}_j-\theta_j^*)\ell_{e_j, \hat{g}}(Z) - \mu \hat{\theta}^\top H \theta^* {-\frac{1}{s}K(\hat{\theta})}\\
%     H_{ij} =~& \|f_i - f_j\|^2 
% \end{align*}
% where $s$ is a positive number. 
% % Assume that:
% % \begin{align*}
% %     {\cal E}(\hat{g}) :=~& P(\tilde{\ell}_{\hat{\theta},g_0}(Z) - \tilde{\ell}_{\theta^*,g_0}(Z)) - P(\tilde{\ell}_{\hat{\theta},\hat{g}}(Z) - \tilde{\ell}_{\theta^*,\hat{g}}(Z)) \\
% %     \leq ~& \left(C_1\|\hat{f}-f^*\|_{2} + C_2\|\Tcal(\hat{f}-f^*)\|_2\right)\error(\hat{g})
% % \end{align*} for some constants $C_1$ ad $C_2$.
% and $\ell(Z; f(X), g_0)$ is $\sigma$-strongly convex in expectation with respect to $f$. For {$\beta \geq \frac{4n}{s}$} and $\mu \leq \min\{\frac{\alpha(1-\nu)}{7}, \frac{1-\nu}{3}, \frac{\alpha\nu}{5}\}$:
% \begin{align*}
%     R(\hat{\theta}) \leq \min_{j} [R(e_j, g_0){+ \frac{\beta}{n} \log(1/\pi_j)}] + 2\, (Z_{n,1} + Z_{n,2}) + \frac{1}{\mu}\error(\hat{g})^2 \label{eqn:lemma9}
% \end{align*}
% \end{lemma}
% \begin{proof}
% \begin{align*}
%     \tilde{R}(\hat{\theta}, g_0) - \tilde{R}(\theta^*, g_0) =~& 
%     P(\tilde{\ell}_{\hat{\theta},g_0}(Z) - \tilde{\ell}_{\theta^*,g_0}(Z))\\
%     =~& 
%     P(\tilde{\ell}_{\hat{\theta},\hat{g}}(Z) - \tilde{\ell}_{\theta^*,\hat{g}}(Z)) + \tilde{\cal E}(\hat{g})\\
%     \leq~& 
%     (P - P_n)(\tilde{\ell}_{\hat{\theta}, \hat{g}}(Z) - \tilde{\ell}_{\theta^*, \hat{g}}(Z)) + \tilde{\cal E}(\hat{g}) 
%     {-\frac{\beta}{n}(K(\hat{\theta})-K(\theta^*))}
%     \tag{By optimality of $\hat{\theta}$, $P_n(\tilde{\ell}_{\hat{\theta}, \hat{g}}(Z) - \tilde{\ell}_{\theta^*, \hat{g}}(Z)) {+ \frac{\beta}{n}(K(\hat{\theta})-K(\theta^*))}\leq 0$}\\
%     =~& \mu\sum_j \hat{\theta}_j \|f_j - f_{\theta^*}\|^2 + \mu \hat{\theta}^\top H\theta^*+ Z_{n,1} + Z_{n,2} + \tilde{\cal E}(\hat{g})\\
%     ~& {-\frac{\beta}{n}(K(\hat{\theta})-K(\theta^*)) + \frac{2}{s}K(\hat{\theta})}
%     \tag{definition of $Z_{n,1}, Z_{n,2}$}
% \end{align*}
% By simple algebraic manipulations note that for any $\theta\in \Theta$:
% \begin{align*}
%     \sum_j \hat{\theta}_j \|f_j - f_{\theta^*}\|^2 = V(\hat{\theta}) + \|f_{\hat{\theta}} - f_{\theta^*}\|^2 \\
%     \hat{\theta}^\top H \theta^* = V(\hat{\theta}) + V(\theta^*) + \|f_{\hat{\theta}} - f_{\theta^*}\|^2
% \end{align*}
% Letting $Z_n=Z_{n,1} + Z_{n,2}$, we have:
% \begin{align*}
%     \tilde{R}(\hat{\theta}, g_0) - \tilde{R}(\theta^*, g_0) \leq~& 2\mu V(\hat{\theta}) + \mu V(\theta^*) + 2\mu \|f_{\hat{\theta}} - f_{\theta^*}\|^2 + Z_n\\
%     ~&+ \tilde{\cal E}(\hat{g}) {-\frac{\beta}{n}(K(\hat{\theta})-K(\theta^*)) + \frac{2}{s}K(\hat{\theta})}\\
%     % \leq~& 2\mu V(\hat{\theta}) + \mu V(\theta^*) + 2\mu \|f_{\hat{\theta}} - f_{\theta^*}\|^2 + Z_n + \error(\hat{g})\, \|f_{\hat{\theta}} - f_{\theta^*}\|^{\gamma}\\
%     % ~& {-\frac{\beta}{n}(K(\hat{\theta})-K(\theta^*)) + \frac{2}{s}K(\hat{\theta})}
% \end{align*}
% Since we have
% \begin{align*}
%     \tilde{\cal E}(\hat{g}) :=~& P(\tilde{\ell}_{\hat{\theta},g_0}(Z) - \tilde{\ell}_{\theta^*,g_0}(Z)) - P(\tilde{\ell}_{\hat{\theta},\hat{g}}(Z) - \tilde{\ell}_{\theta^*,\hat{g}}(Z)) \\
%     \leq~& \left(|C_Y-C_f|\|f_{\hat{\theta}}-f^*\|_{2}+C_f\|\Tcal(f_{\hat{\theta}}-f_{\theta^*})\|_2\right)\error(\hat{g})\\
%     \leq~& \mu\left(\|f_{\hat{\theta}}-f_{\theta^*}\|_{2}^2+\|\Tcal(f_{\hat{\theta}}-f_{\theta^*})\|_2^2\right)+\underbrace{\frac{1}{2\mu}\left((C_Y-C_f)^2+C_f^2\right)\error(\hat{g})^2}_{{\cal A}(\hat{g})} \tag{By Young's inequality}\\
% \end{align*}
% % By Young's inequality, for any $a\geq 0$ and $b\geq 0$, we have $a\,b\leq \frac{a^p}{p} + \frac{b^q}{q}$, such that $\frac{1}{p}+\frac{1}{q}=1$. Taking $a= \left(\frac{\gamma}{2\mu}\right)^{\gamma/2} \error(\hat{g})$ and $b=\left(\frac{2\mu}{\gamma}\right)^{\gamma/2}\|f_{\hat{\theta}}-f_{\theta^*}\|^\gamma$ and $q=2/\gamma$ and $p=\frac{1}{1-1/q}=\frac{1}{1-\gamma/2}=\frac{2}{2-\gamma}$
% % \begin{align}
% %     \error(\hat{g})\, \|f_{\hat{\theta}} - f_{\theta^*}\|^{\gamma} \leq \underbrace{\frac{2}{2-\gamma}\left(\frac{\gamma}{2\mu}\right)^{\frac{\gamma}{2-\gamma}} \error(\hat{g})^{\frac{2}{2-\gamma}}}_{{\cal A}(\hat{g})}+ \mu \|f_{\hat{\theta}} - f_{\theta^*}\|^2
% % \end{align}
% Thus we can derive:
% \begin{align}
%     \tilde{R}(\hat{\theta}, g_0) - \tilde{R}(\theta^*, g_0) \leq~& 2\mu V(\hat{\theta}) + \mu V(\theta^*) + 3\mu \|f_{\hat{\theta}} - f_{\theta^*}\|^2 + \mu\|\Tcal(f_{\hat{\theta}}-f_{\theta^*})\|^2\\\label{eqn:first-upper}
%     ~&  + Z_n + {\cal A}(\hat{g})  -\frac{\beta}{n}(K(\hat{\theta})-K(\theta^*)) + \frac{2}{s}K(\hat{\theta})
% \end{align}
% By Lemma~\ref{lem:strong-conv-tilde}, we have:
% \begin{align*}
%     \tilde{R}(\hat{\theta}, g_0) - \tilde{R}(\theta^*, g_0) \geq~& \mu (V(\theta^*) - V(\hat{\theta})) {+\frac{\beta}{n}(K(\theta^*)-K(\hat{\theta}))} +(1-\nu-\mu)\|\Tcal(f_{\hat{\theta}} - f_{\theta^*})\|^2\\
%     ~&+ \left(\alpha(1-\nu) - \mu\right)\|f_{\hat{\theta}} - f_{\theta^*}\|^2 \\
%     \geq~& \mu (V(\theta^*) - V(\hat{\theta})) {+\frac{\beta}{n}(K(\theta^*)-K(\hat{\theta}))} + 6\mu  \|f_{\hat{\theta}} - f_{\theta^*}\|^2 \\
%     ~&+ 2\mu\|\Tcal(f_{\hat{\theta}} - f_{\theta^*})\|^2
%     \tag{Since: $\mu \leq \min\{\frac{\alpha(1-\nu)}{7},\frac{1-\nu}{3}\}$}
% \end{align*}
% Combining the latter upper and lower bound on $\tilde{R}(\hat{\theta}, g_0) - \tilde{R}(\theta^*, g_0)$, we have:
% \begin{align*}
%     3\mu \|f_{\hat{\theta}} - f_{\theta^*}\|^2 +\mu \|\Tcal(f_{\hat{\theta}} - f_{\theta^*})\|^2\leq 3\mu V(\hat{\theta}) + Z_n {+ \frac{2}{s} K(\hat{\theta})} + {\cal A}(\hat{g})
% \end{align*}
% Plugging the bound on $3\mu\|f_{\hat{\theta}} - f_{\theta^*}\|^2$ back into Equation~\eqref{eqn:first-upper} we have:
% \begin{align*}
%     \tilde{R}(\hat{\theta}, g_0) - \tilde{R}(\theta^*, g_0) \leq ~& 5\mu V(\hat{\theta}) + \mu V(\theta^*) + 2\, Z_n + 2\,{\cal A}(\hat{g}) \\
%     ~& {+\frac{\beta}{n}(K(\theta^*)-K(\hat{\theta})) + \frac{4}{s}K(\hat{\theta})}
% \end{align*}
% By Lemma~\ref{lem:v-versus-tilde}, we have:
% \begin{align*}
%     R(\hat{\theta}, g_0) \leq~& \tilde{R}(\theta^*, g_0) + \mu V(\theta^*)+ \left(5\mu - \alpha\nu\right) V(\hat{\theta}) + 2 Z_n + 2{\cal A}(\hat{g}) \\
%     ~& {+\frac{\beta}{n}K(\theta^*) + \left(\frac{4}{s} - \frac{\beta}{n}\right)K(\hat{\theta})}\\
%     \leq~& \tilde{R}(\theta^*, g_0) + \mu V(\theta^*) {+ \frac{\beta}{n}K(\theta^*)} + 2 Z_n + 2{\cal A}(\hat{g})   \tag{since $\mu\leq \frac{\nu\alpha}{5}$,  $V(\hat{\theta})>0$, {$\beta \geq \frac{4n}{s}$ and $K(\hat{\theta}\geq 0 )$}}\\
%     =~& \min_{\theta\in \Theta} [\tilde{R}(\theta, g_0) + \mu V(\theta){+ \frac{\beta}{n}K(\theta^*)}] + 2 Z_n + 2{\cal A}(\hat{g}) \tag{Since $\theta_*$ minimizes $H(\theta, g_0)$}\\
%     \leq~& \min_{j} [\tilde{R}(e_j, g_0) + \mu V(e_{j}) {+ \frac{\beta}{n}\log(1/\pi_j)}]+ 2 Z_n + 2{\cal A}(\hat{g}) \\
%     =~& \min_{j} [R(e_j, g_0){+ \frac{\beta}{n} \log(1/\pi_j)}] + 2 Z_n +2{\cal A}(\hat{g}) \tag{Since $\tilde{R}(e_j, g_0)=R(e_j, g_0)$ and $V(e_j)=0$ for all $e_j$}
% \end{align*}
% \end{proof}

% \begin{lemma}\label{lem:z1}
% For any positive $s\leq \min\left\{ \frac{\mu n}{32C_{Y,f}(1-\nu)^2}, \frac{n}{4\sqrt{6} C_f (1-\nu) \sqrt{C_{Y,f}}}\right\}$, where $C_{Y,f} = ((C_Y+C_f)^2+\alpha^2 C_f^2)$:
% \begin{align}
%     {\Pr\left(Z_{n,1} \geq \frac{\log(1/\delta)}{s}\right) \leq \delta}
% \end{align}
% \end{lemma}
% \begin{proof}
% By Lemma~\ref{lem:exp-tail} it suffices to show that:
% \begin{align}
%     {I:=\E\left[\exp\left\{s Z_{n,1}\right\}\right] \leq 1}
% \end{align}
% % First note that:
% % \begin{align}
% %     I = \E\left[\exp\left\{s Z_{n,1}\} \exp\{- \log(M)\right\}\right] = \frac{1}{M} \E\left[\exp\left\{s Z_{n,1}\right\}\right]
% % \end{align}
% First, by the mean value theorem, we observe that:
% \begin{align*}
%     |\ell_{\hat{\theta},\hat{g}}(Z,X) - \ell_{\hat{\theta},\hat{g}}(Z,X)|\leq 2(C_Y+C_f)|\hat{\Tcal}(f_{\hat{\theta}}-f_{\theta^*})| + 2\alpha C_f|f_{\hat{\theta}}-f_{\theta^*}|
% \end{align*}
% Define:
% \begin{align*}
%     f_{\ell,\theta}(X,Z) = 2(C_Y+C_f)\hat{\Tcal}(f_{\theta}) + 2\alpha C_f f_{\theta}
% \end{align*}
% By the symmetrization argument:
% \begin{align*}
%     I =~& \E\left[\exp\left\{s \left((1-\nu)(P-P_n)(\ell_{\hat{\theta}}(Z,X) - \ell_{\theta^*}(Z,X)) - \mu \sum_j \hat{\theta}_j \|f_j - f_{\theta^*}\|^2 {-\frac{1}{s}\sum_j\hat{\theta}_j \log(1/\pi_j)}\right)\right\}\right]\\
%     \leq~& \E\left[\exp\left\{s \max_{\theta} \left((1-\nu)(P-P_n)(\ell_{\hat{\theta}}(Z,X) - \ell_{\theta^*}(Z,X)) - \mu \sum_j \theta_j \|f_j - f_{\theta^*}\|^2 {+\frac{1}{s}\sum_j\theta_j \log(\pi_j)} \right)\right\}\right]\\
%     \leq~& \E\left[\exp\left\{s \max_{\theta} \left(2\,(1-\nu)P_{n,\epsilon}(\ell_{\hat{\theta}}(Z,X) - \ell_{\theta^*}(Z,X)) - \mu \sum_j \theta_j \|f_j - f_{\theta^*}\|^2 {+\frac{1}{s}\sum_j\theta_j\log(\pi_j)}\right)\right\}\right]\tag{By symmetrization (Lemma \ref{sym})}\\
%     % \leq~& \E\left[\exp\left\{s \max_{\theta} \left(2\,(1-\nu) P_{n,\epsilon}(f_{\ell,\hat{\theta}}(X,Z) - f_{\ell,\theta^*}(X,Z)) - \mu \sum_j \theta_j \|f_j - f_{\theta^*}\|^2 +\frac{1}{s}\sum_j\theta_j \log(\pi_j)\right)\right\}\right] \tag{contr.}
% \end{align*}
% By a contraction argument (Lemma \ref{contr}), we have that:
% \begin{align*}
%     P_{n,\epsilon}(\ell_{\hat{\theta}}(Z,X) - \ell_{\theta^*}(Z,X)) =~& P_{n,\epsilon}\left((Y - \Tcal(f_{\theta}))^2 - (Y - \Tcal(f_{\theta^*}))^2 + \alpha  f_{\theta}^2 - \alpha f_{\theta^*}^2\right)\\
%     \leq~& P_{n,\epsilon}\left(2(C_Y+C_f)\Tcal(f_{\theta}) - \Tcal(f_{\theta}^*) + \alpha f_{\theta}-f_{\theta^*}\right)\\
%     =~&P_{n,\epsilon}(f_{\ell,\hat{\theta}}(X,Z) - f_{\ell,\theta^*}(X,Z))
% \end{align*}
% Then, we get:
% \begin{align*}
%     I\leq~& \E\left[\exp\left\{s \max_{\theta} \left(2\,(1-\nu) P_{n,\epsilon}(f_{\ell,\hat{\theta}}(X,Z) - f_{\ell,\theta^*}(X,Z)) - \mu \sum_j \theta_j \|f_j - f_{\theta^*}\|^2 +\frac{1}{s}\sum_j\theta_j \log(\pi_j)\right)\right\}\right]
% \end{align*}

% Since the quantity:
% \begin{align*}
%    2\,(1-\nu) P_{n,\epsilon}(f_{\ell,\hat{\theta}}(X,Z) - f_{\ell,\theta^*}(X,Z)) - \mu \sum_j \theta_j \|f_j - f_{\theta^*}\|^2 +\frac{1}{s}\sum_j\theta_j \log(\pi_j)
% \end{align*}
% is a linear function of $\theta$ (recall that $f_{\theta}=\sum_j \theta_j f_j$), the maximum over $\theta$ is attained at one of $e_1, \ldots, e_M$. Thus:
% \begin{align*}
%     I \leq ~&\E\left[\exp\left\{s \max_{j} \left(2\,(1-\nu) P_{n,\epsilon}(f_{\ell,\hat{\theta}}(X,Z) - f_{\ell,\theta^*}(X,Z)) - \mu \|f_j - f_{\theta^*}\|^2{+\frac{1}{s}\log(\pi_j)}\right)\right\}\right]\\
%     \leq~& \E\left[\max_{j}{\pi_j} \exp\left\{s \left(2\,(1-\nu) P_{n,\epsilon}(f_{\ell,\hat{\theta}}(X,Z) - f_{\ell,\theta^*}(X,Z)) - \mu \|f_j - f_{\theta^*}\|^2\right)\right\}\right]\\
%     \leq~& \E\left[\sum_{j} {\pi_j} \exp\left\{s \left(2\,(1-\nu) P_{n,\epsilon}(f_{\ell,\hat{\theta}}(X,Z) - f_{\ell,\theta^*}(X,Z)) - \mu \|f_j - f_{\theta^*}\|^2\right)\right\}\right]\\
%     =~&\sum_{j} {\pi_j} \E\left[ \exp\left\{s \left(2\,(1-\nu) P_{n,\epsilon}(f_{\ell,\hat{\theta}}(X,Z) - f_{\ell,\theta^*}(X,Z)) - \mu \|f_j - f_{\theta^*}\|^2\right)\right\}\right]\\
%     =~& \sum_{j} {\pi_j}\E_{Z_{1:n},X_{1:n}}\E_\epsilon\left[ \exp\left\{s \left(2\,(1-\nu) P_{n,\epsilon}(f_{\ell,\hat{\theta}}(X,Z) - f_{\ell,\theta^*}(X,Z)) - \mu \|f_j - f_{\theta^*}\|^2\right)\right\}\right]\\
%     =~& \sum_{j} {\pi_j}\E_{Z_{1:n},X_{1:n}}\left[ \E_\epsilon\left[ \exp\left\{s 2\,(1-\nu) P_{n,\epsilon}(f_{\ell,\hat{\theta}}(X,Z) - f_{\ell,\theta^*}(X,Z))\right\}\right] \exp\left\{- s \mu \|f_j - f_{\theta^*}\|^2\right\}\right]
% \end{align*}
% By Hoeffding's inequality in Lemma~\ref{lem:hoeffding}, we can bound the first term in the product:
% \begin{align*}
%     ~&\E_\epsilon\left[ \exp\left\{s 2\,(1-\nu) P_{n,\epsilon}(f_{\ell,\hat{\theta}}(X,Z) - f_{\ell,\theta^*}(X,Z))\right\}\right]\\
%     ~&\leq \exp\left\{\frac{2 s^2}{n}\,(1-\nu)^2 P_n(f_{\ell,\hat{\theta}}(X,Z) - f_{\ell,\theta^*}(X,Z))^2\right\}\\
%     ~&\leq \exp\left\{\frac{16 s^2}{n}\,(1-\nu)^2 P_n\left[(C_Y+C_f)^2(\Tcal(f_{\hat{\theta}})(Z)-\Tcal(f_{\theta^*}(Z))^2 + \alpha^2 C_f^2(f_{\hat{\theta}}(X)-f_{\theta^*}(X))^2\right]\right\}\\
%     ~&\leq \exp\left\{\frac{16 s^2}{n}\,(1-\nu)^2 P_n((C_Y+C_f)^2+\alpha^2 C_f^2)(f_{\hat{\theta}}(X)-f_{\theta^*}(X))^2 \right\}\\
% \end{align*}
% Let $C_{Y,f}:=((C_Y+C_f)^2+\alpha^2 C_f^2)$, then we have:
% \begin{align*}
%     I\leq~&\sum_{j} {\pi_j}\E_{X_{1:n}}\left[ \exp\left\{\frac{16s^2}{n}\,(1-\nu)^2C_{Y,f} P_n(f_{j}(X) - f_{\theta^*}(X))^2\right\}\, \exp\left\{- s \mu \|f_j - f_{\theta^*}\|^2\right\}\right]\\
%     =~& \sum_{j} {\pi_j}\E_{X_{1:n}}\left[ \exp\left\{\frac{16s^2}{n}\,(1-\nu)^2 C_{Y,f}  P_n(f_{j}(X) - f_{\theta^*}(X))^2\right\}\, \exp\left\{- s \mu P (f_j(X) - f_{\theta^*}(X))^2\right\}\right]\\
%     =~& \sum_{j} {\pi_j} \E_{X_{1:n}}\left[ \exp\left\{\frac{16s^2\,(1-\nu)^2C_{Y,f} }{n} \left(P_n - \frac{\mu\, n}{2 s(1-\nu)^2 C_{Y,f} }P\right)(f_{j}(X) - f_{\theta^*}(X))^2\right\}\right]
% \end{align*}
% Since $s$ is small enough such that: $\frac{\mu\, n}{16 s (1-\nu)^2 C_{Y,f}}\geq 2$ (which is satisfied by our assumptions on $s$), we have that:
% \begin{align*}
%     I \leq~& \sum_{j} \pi_j \E_{Z_{1:n}}\left[ \exp\left\{\frac{16s^2\,(1-\nu)^2 C_{Y,f}}{n} \left(P_n - 2P\right)(f_{j}(Z) - f_{\theta^*}(Z))^2\right\}\right]\\
%     =~& \sum_{j} \pi_j\E_{Z_{1:n}}\left[ \exp\left\{\frac{16s^2\,(1-\nu)^2 C_{Y,f}}{n} \left(\left(P_n - P\right)(f_{j}(Z) - f_{\theta^*}(Z))^2 - P (f_{j}(Z) - f_{\theta^*}(Z))^2 \right)\right\}\right]
% \end{align*}
% Since $f_j(Z)\in [-C_f, C_f]$, we have that $f_j(X) - f_{\theta^*}(X)\in [-2C_f, 2C_f]$. Thus:
% \begin{align*}
% (f_j(X)-f_{\theta^*}(X))^2\geq \frac{(f_j(X)-f_{\theta^*}(X))^2}{4C_f^2} (f_j(X)-f_{\theta^*}(X))^2 = \frac{1}{4C_f^2} (f_j(X)-f_{\theta^*}(X))^4.
% \end{align*}
% Thus:
% \begin{align*}
%     I \leq~& \sum_{j} \pi_j \E_{Z_{1:n}}\left[ \exp\left\{\frac{16s^2\,(1-\nu)^2 C_{Y,f}}{n} \left(\left(P_n - P\right)(f_{j}(Z) - f_{\theta^*}(Z))^2 - \frac{1}{4C_f^2} P (f_{j}(Z) - f_{\theta^*}(Z))^4 \right)\right\}\right]
% \end{align*}
% It suffices to show that each of the summands is upper bounded by $1$. Then their weighted average is also upper bounded by $1$.
% Applying Lemma~\ref{lem:bernstein} with:
% \begin{align}
%     V_i =~& (f_j(Z_i) - f_{\theta^*}(Z_i))^2\\
%     c =~& 4C_f^2\\
%     \lambda =~& \frac{16s^2(1-\nu)^2C_{Y,f}}{n^2}\\
%     c_0 =~& \frac{1}{4C_f^2}
% \end{align}
% we have that as long as:
% \begin{align}
%     \lambda := \frac{16s^2(1-\nu)^2C_{Y,f}}{n^2} \leq \frac{2c_0}{1+2c_0 c} =: \frac{2\cdot 1/4C_f^2}{1 + 2} = \frac{1}{6C_f^2}
% \end{align}
% Then:
% \begin{align}
%     \E_{Z_{1:n}}\left[ \exp\left\{\frac{16s^2\,(1-\nu)^2 C_{Y,f}}{n} \left(\left(P_n - P\right)(f_{j}(Z) - f_{\theta^*}(Z))^2 - \frac{1}{4C_f^2} P (f_{j}(Z) - f_{\theta^*}(Z))^4 \right)\right\}\right]\leq 1
% \end{align}
% Thus we need that:
% \begin{align}
%     s \leq \frac{n}{4\sqrt{6} C_f (1-\nu) \sqrt{C_{Y,f}}}
% \end{align}
% which is satisfied by our assumptions on $s$.
% \end{proof}

% \begin{lemma}\label{lem:z2}
% For any positive $s\leq \frac{2 \mu n}{\nu (\nu C_{\ell,f} + 2 \mu C_{\ell})}$ where $C_{\ell,f} = 8((C_Y+C_f)^2+\alpha^2C_f^2)$ and $C_{\ell} = (C_Y+C_f)^2+\alpha C_f^2$:
% \begin{align}
%     \Pr\left(Z_{n,2} \geq \frac{\log(1/\delta)}{s}\right) \leq \delta
% \end{align}
% \end{lemma}
% \begin{proof}
% It suffices to show that:
% \begin{align}
%     I := \E\left[\exp\left\{s Z_{n,2})\right\}\right] \leq 1
% \end{align}
% % First note that:
% % \begin{align}
% %     I = \E\left[\exp\left\{s Z_{n,2}\right\} \exp\left\{-\log(M)\right\}\right] = \frac{1}{M}  \E\left[\exp\left\{s Z_{n,2}\right\}\right]
% % \end{align}
% By definition of $Z_{n,2}$:
% \begin{align*}
%     I =~& \E\left[\exp\left\{s \left( \nu(P-P_n) \sum_{j} (\hat{\theta}_j-\theta_j^*)\ell_{e_j}(Z) - \mu \hat{\theta}^\top H \theta^*  -\frac{1}{s}\sum_j\hat{\theta}_j \log(1/\pi_j) \right)\right\}\right]\\
%     =~& \E\left[\exp\left\{s \left( \nu(P-P_n) \sum_{j} \hat{\theta}_j \ell_{e_j}(Z) + \frac{1}{s}\sum_j\hat{\theta}_j \log(\pi_j)\right)\right\}\right.\\
%     ~&\quad \left.\exp\left\{s\left( - \nu (P-P_n)\sum_{j} \theta_j^*\ell_{e_j}(Z) - \mu \hat{\theta}^\top H \theta^* \right)\right\}\right]
% \end{align*}
% Note that:
% \begin{align*}
%     - \nu (P-P_n)\sum_{j} \theta_j^*\ell_{e_j}(Z) - \mu \hat{\theta}^\top H \theta^* =~& \sum_{j} \theta_j^* \left( -\nu (P-P_n) \ell_{e_j}(Z) - \mu \hat{\theta}^\top H_{\cdot, j}\right) \\
%     =~& \sum_{j} \theta_j^* \underbrace{\left( -\nu (P-P_n) \ell_{e_j}(Z) - \mu \sum_{k} \hat{\theta}_k \|f_k - f_j\|^2\right)}_{\Phi_j}
% \end{align*}
% Thus:
% \begin{align*}
%     I=~&  \E\left[\exp\left\{s \left( \nu(P-P_n) \sum_{j} \hat{\theta}_j \ell_{e_j}(Z)  +\frac{1}{s}\sum_j\hat{\theta}_j \log(\pi_j)\right)\right\} \exp\left\{\sum_{j}\theta_j^* s\Phi_j\right\}\right]
% \end{align*}
% By Jensen's inequality:
% \begin{align*}
%     \exp\left\{\sum_{j}\theta_j^* s\Phi_j\right\} \leq \sum_{j}\theta_j^* \exp\left\{s\Phi_j\right\}
% \end{align*}
% Thus:
% \begin{align*}
%     I\leq~& \sum_{j}\theta_j^*  \E\left[\exp\left\{s \left( \nu(P-P_n) \sum_{k} \hat{\theta}_k \ell_{e_k}(Z){+\frac{1}{s}\sum_k \hat{\theta}_k \log(\pi_k)}\right)\right\} \exp\left\{ s\Phi_j\right\}\right]\\
%     =~& \sum_{j}\theta_j^* \E\left[\exp\left\{s \left( \nu(P-P_n) \sum_{k} \hat{\theta}_k \ell_{e_k}(Z) + \Phi_j {+\frac{1}{s}\sum_k\hat{\theta}_k \log(\pi_k)}\right)\right\}\right]\\
%     =~& \sum_{j}\theta_j^*  \E\left[\exp\left\{s \left( \nu(P-P_n) \left(\sum_{k} \hat{\theta}_k \ell_{e_k}(Z) - \ell_{e_j}(Z)\right) - \mu \sum_{k} \hat{\theta}_k \|f_k - f_j\|^2{+\frac{1}{s}\sum_k\hat{\theta}_k \log(\pi_k)}\right)\right\}\right]\\
%     =~& \sum_{j}\theta_j^*  \E\left[\exp\left\{s \left( \nu(P-P_n) \sum_{k} \hat{\theta}_k \left(\ell_{e_k}(Z) - \ell_{e_j}(Z)\right) - \mu \sum_{k} \hat{\theta}_k \|f_k - f_j\|^2 + \frac{1}{s}\sum_k\hat{\theta}_k \log(\pi_k)\right)\right\}\right]\\
%     \leq~& \sum_{j}\theta_j^* \E\left[\exp\left\{s \max_{\theta} \left( \nu(P-P_n) \sum_{k} \theta_k \left(\ell_{e_k}(Z) - \ell_{e_j}(Z)\right) - \mu \sum_{k} \theta_k \|f_k - f_j\|^2 + \frac{1}{s}\sum_k\theta_k 
%     \log(\pi_k)\right)\right\}\right]
% \end{align*}
% Since the objective in the exponent is linear in $\theta$ it takes its value at one of the extreme points $e_1, \ldots, e_M$:
% \begin{align*}
%     I \leq~& \sum_{j}\theta_j^*  \E\left[\exp\left\{s \max_{k} \left( \nu(P-P_n) \left(\ell_{e_k}(Z) - \ell_{e_j}(Z)\right) - \mu \|f_k - f_j\|^2+\frac{1}{s} \log(\pi_k)\right)\right\}\right]\\
%     =~& \sum_{j}\theta_j^*  \E\left[\max_{k}{\pi_k}\exp\left\{s  \left( \nu(P-P_n) \left(\ell_{e_k}(Z) - \ell_{e_j}(Z)\right) - \mu \|f_k - f_j\|^2\right)\right\}\right]\\
%     \leq~& \sum_{j}\theta_j^* \E\left[\sum_{k}{\pi_k} \exp\left\{s  \left( \nu(P-P_n) \left(\ell_{e_k}(Z) - \ell_{e_j}(Z)\right) - \mu \|f_k - f_j\|^2\right)\right\}\right]\\
%     \leq~& \sum_{j}\theta_j^* \sum_{k} {\pi_k}\E\left[ \exp\left\{s  \left( \nu(P-P_n) \left(\ell_{e_k}(Z) - \ell_{e_j}(Z)\right) - \mu \|f_k - f_j\|^2\right)\right\}\right]
% \end{align*}
% By a the mean-value theorem, we have:
% \begin{align*}
%     |\ell_{e_k}(Z,X) - \ell_{e_j}(Z,X)|\leq 2(C_Y+C_f)|\hat{\Tcal}(f_{k}-f_{j})| + 2\alpha C_f|f_{\hat{\theta}}-f_{\theta^*}|
% \end{align*}
% Thus,
% \begin{align*}
%      \|\ell_{e_k} - \ell_{e_j}\|^2 \leq~& 8(C_Y+C_f)^2\|\hat{\Tcal}(f_{k}-f_{j})\|^2 + 8\alpha^2 C_f^2 \|f_k - f_j\|^2\\
%      \leq~& 8\left((C_Y+C_f)^2+\alpha^2C_f^2\right) \|f_k - f_j\|^2
% \end{align*}
% % By Lipschitzness of $\ell$, we have that:
% % \begin{align*}
% %     \|\ell_{e_k} - \ell_{e_j}\|^2 \leq C_b^2 \|f_k - f_j\|^2
% % \end{align*}
% Let $C_{\ell,f} = 8\left((C_Y+C_f)^2+\alpha^2C_f^2\right)$, thus:
% \begin{align*}
%     I \leq~& \sum_{j}\theta_j^* \sum_{k} {\pi_k}\E\left[ \exp\left\{s  \left( \nu(P-P_n) \left(\ell_{e_k}(Z) - \ell_{e_j}(Z)\right) - \frac{\mu}{C_{\ell,f}} \|\ell_{e_k} - \ell_{e_j}\|^2\right)\right\}\right] \\
%     =~& \sum_{j}\theta_j^* \sum_{k} {\pi_k}\E\left[ \exp\left\{s  \left( \nu(P-P_n) \left(\ell_{e_k}(Z) - \ell_{e_j}(Z)\right) - \frac{\mu}{C_{\ell,f}} P(\ell_{e_k}(Z) - \ell_{e_j}(Z))^2\right)\right\}\right]\\
%     =~& \sum_{j}\theta_j^* \sum_{k}{\pi_k} \E\left[ \exp\left\{s \nu \left((P-P_n) \left(\ell_{e_k}(Z) - \ell_{e_j}(Z)\right) - \frac{\mu}{\nu\, C_{\ell,f}} P(\ell_{e_k}(Z) - \ell_{e_j}(Z))^2\right)\right\}\right]
% \end{align*}
% Invoking Lemma~\ref{lem:bernstein} with:
% \begin{align*}
%     V_i =~& \ell_{e_j}(Z_i) - \ell_{e_k}(Z_i)\\
%     c_0 =~& \frac{\mu}{\nu C_{\ell,f}}\\
%     \lambda =~& \frac{s \nu}{n}\\
%     c=~& C_{\ell} :=(C_Y+C_f)^2+\alpha C_f^2\\ 
% \end{align*}
% Hence we need:
% \begin{align*}
%     ~&\lambda := \frac{s\nu}{n} \leq \frac{2c_0}{1 + 2c_0c} =: \frac{2 \mu}{\nu C_{\ell,f} (1 + 2 \frac{\mu}{\nu C_{\ell,f}} C_{\ell})} = \frac{2 \mu}{(\nu C_{\ell,f} + 2 \mu C_{\ell})} \implies s \leq \frac{2 \mu n}{\nu (\nu C_{\ell,f} + 2 \mu C_{\ell})}
% \end{align*}
% \end{proof}
\subsection{Proof of Oracle Inequality for Convex ERM}
\hlan{Probably don't need this section any more.}
\begin{theorem}[Convex aggregation via ERM]\label{convex-erm}
Assume that $Y$ is almost surely bounded by $b$, each candidate model $f_j$ is almost uniformly bounded in $[-b,b]$ almost surely. Moreover, assume that $g_0$ is bounded away from 0, i.e. exist $c_0>0$ such that $g_0\geq c_0$. Then, for any $\delta>0$, with probability$`>1-\delta$, the output of ERM using the loss function described above, $\hat{\theta}$, satisfies:
\begin{align*}
    R(\hat{\theta},g_0) - \min_{\theta\in\Theta}R(\theta,g_0) \leq O\left(\frac{M+\log(1/\delta)}{n} + \frac{1}{\alpha}\error(\hat{g})^2\right),
\end{align*} 
where $\error(\hat{g}) = 2(\frac{1}{c_0}+1)^2\E[H(\hat{g},g_0)^2]^\frac{1}{2}$. 
\end{theorem}

Before the proof, we introduce the shorthand notation:
\begin{align*}
    % \tilde{{\cal E}}(\hat{g}) :=~& \E\big[\tilde{\ell}_{\hat{\theta},g_0}(Z) - \tilde{\ell}_{\theta^*,g_0}(Z)\big] - \E\big[\tilde{\ell}_{\hat{\theta},\hat{g}}(Z) - \tilde{\ell}_{\theta^*,\hat{g}}(Z)\big]\\
    {\cal E}(\hat{g}) :=~& \E\big[{\ell}_{\hat{\theta},g_0}(Z) - {\ell}_{\theta^*,g_0}(Z)\big] - \E\big[{\ell}_{\hat{\theta},\hat{g}}(Z) - {\ell}_{\theta^*,\hat{g}}(Z)\big]
    % ~\leq~ \error(\hat{g})\, \|f_{\hat{\theta}} - f_{\theta^*}\|^{\gamma}
\end{align*}

\begin{lemma}\label{lemma:diffndiff}
    Suppose that $|Y|$ is almost surely bounded by some constant $C_Y$, and $|f(X)|$ is almost surely uniformly bounded by a constant $C_f$. Then, let $\hat{g}$ be the conditional density estimated using the standard MLE approach, then we have:
     \begin{align*}
    {\cal E}(\hat{g}) 
    \leq~& \left(|C_Y-C_f|\|\hat{f}-f^*\|_{2}+C_f\|\Tcal(\hat{f}-f^*)\|_2\right)\delta_{n,\Gcal}\\
    \end{align*}
    Moreover, if $\hat{g}$ is estimated using the $\chi^2$-MLE approach, then we have:
    \begin{align*}
    {\cal E}(\hat{g}) 
    \leq~& \left(|C_Y-C_f|\|\hat{f}-f^*\|_{\infty}+C_f\|\Tcal(\hat{f}-f^*)\|_{\infty}\right) \delta_{n,\Gcal}\\
    \leq~& \left(|C_Y-C_f|\|\hat{f}-f^*\|_{2}^{\gamma}+C_f\|\Tcal(\hat{f}-f^*)\|_{2}^{\gamma}\right) \delta_{n,\Gcal}\\
    \end{align*} where $\error(\hat{g}) =2(\frac{1}{c_0} + 1)^2\mathbb{E}\left[H(\hat{g},g_0)^2\right]^{\frac{1}{2}}$.
\end{lemma}
\begin{proof}
    \begin{align*}
    {\cal E}(\hat{g})
    =~& \E [-2Y (\Tcal-\hat{\Tcal})(\hat{f}-f^*) + (\Tcal(\hat{f}))^2-(\Tcal(f^*))^2 -(\hat{\Tcal}(\hat{f}))^2+(\hat{\Tcal}(f^*))^2]\\
    =~&\E[ -2Y (\Tcal-\hat{\Tcal})(\hat{f}-f^*) + (\Tcal-\hat{\Tcal})(\hat{f}-f^*)\hat{\Tcal}(\hat{f}+f^*)]\\
    ~&+ \E[\Tcal(\hat{f}-f^*)(\Tcal - \hat{\Tcal})(\hat{f}+f^*)]\\
    \leq~& \E[2|C_Y-C_f||(\Tcal-\hat{\Tcal})(\hat{f}-f^*)| + | \Tcal(\hat{f}-f^*)(\Tcal-\hat{\Tcal})(\hat{f}+f^*)|]
    % \leq~&\mathbb{E} \left[2C_Y \left(\int (f-f')^2g_0dx\right)^{\frac{1}{2}}\left(\int\left(\frac{g}{g_0}-1\right)^2g_0dx\right)^{\frac{1}{2}}\right] \tag{Here we need the true likelihood to be bounded away from $0$} \\
    % \leq~&\mathbb{E} \left\{2C_Y \mathbb{E}[(f-f')^2|Z]^{\frac{1}{2}}\mathcal{\chi}^2(g(X|Z),g_0(X|Z))^{\frac{1}{2}}\right\}\\
    % \leq~& 2C_Y \mathbb{E} \left[\mathbb{E}[(f-f')^2|Z] \right]^{\frac{1}{2}}\mathbb{E} \left[\mathcal{\chi}^2(g(X|Z),g_0(X|Z))\right]^{\frac{1}{2}}\\
    % =~&  2C_Y\|f-f'\|_{2} \mathbb{E} \left[\mathcal{\chi}^2(g(X|Z),g_0(X|Z))\right]^{\frac{1}{2}}\\
\end{align*}
First, we analyze the $(\Tcal-\hat{\Tcal})(\hat{f}-f^*)$ term:
\begin{align*}
    \mathbb{E}[|(\Tcal-\hat{\Tcal})(\hat{f}-f^*)|]\leq~&\mathbb{E}\left[\int |(\hat{f}-f^*)(g-g_0)|dx\right]\\
    \leq~&\mathbb{E}\left[\left(\int (\hat{f}-f^*)^2g_0dx\right)^{\frac{1}{2}}\left(\int\left(\frac{g}{g_0}-1\right)^2g_0dx\right)^{\frac{1}{2}}\right]\\
    =~& \mathbb{E}\left[\mathbb{E}[(\hat{f}-f^*)^2|Z]^{\frac{1}{2}}\mathcal{\chi}^2(g(X|Z),g_0(X|Z))^{\frac{1}{2}}\right]\\
    \leq~& \mathbb{E}\left[\mathbb{E}[(\hat{f}-f^*)^2|Z] \right]^{\frac{1}{2}}\mathbb{E} \left[\mathcal{\chi}^2(g(X|Z),g_0(X|Z))\right]^{\frac{1}{2}}\\
    =~& \|\hat{f}-f^*\|_{2} \mathbb{E} \left[\mathcal{\chi}^2(g(X|Z),g_0(X|Z))\right]^{\frac{1}{2}}.
\end{align*}

Next, we analyze the $\Tcal(\hat{f}-f^*)(\Tcal-\hat{\Tcal})(\hat{f}+f^*)$ term:
\begin{align*}
    \mathbb{E}\left[| \Tcal(\hat{f}-f^*)(\Tcal-\hat{\Tcal})(\hat{f}+f^*)|\right] \leq~& 
    \mathbb{E}\left[|\Tcal(\hat{f}-f^*)|\mathbb{E}[(\hat{f}+f^*)^2|Z]^{\frac{1}{2}}\mathcal{\chi}^2(g(X|Z),g_0(X|Z))^{\frac{1}{2}}\right]\\
    \leq~& 2C_f\mathbb{E}\left[|\Tcal(\hat{f}-f^*)|\mathcal{\chi}^2(g(X|Z),g_0(X|Z))^{\frac{1}{2}}\right]\\
    \leq~&2C_f\|\Tcal(\hat{f}-f^*)\|_2\mathbb{E}\left[\mathcal{\chi}^2(g(X|Z),g_0(X|Z))\right]^{\frac{1}{2}}.
\end{align*}

Since $g_o(X|Z)\geq c_0$, we have:
\begin{align*}
    \mathcal{\chi}^2(g(X|Z),g_0(X|Z)) =~&\int\left(\frac{g}{g_0}-1\right)^2g_0dx\\
    =~&\int\left(\frac{\sqrt{g}}{\sqrt{g_0}}+1\right)^2\left(\frac{\sqrt{g}}{\sqrt{g_0}}-1\right)^2g_0dx\\
    \leq~&\int\left(\frac{1}{c_0}+1\right)^2\left(\frac{\sqrt{g}}{\sqrt{g_0}}-1\right)^2g_0dx\\
    =~&\left(\frac{1}{c_0}+1\right)^2 H(g,g_0)^2.
\end{align*}

Thus, we get:
\begin{align*}
    {\cal E}(\hat{g}) \leq~& \left(2|C_Y-C_f|\|\hat{f}-f^*\|_{2}+2C_f\|_2\Tcal(\hat{f}-f^*)\|\right)(\frac{1}{c_0} + 1)^2\mathbb{E}\left[H(g,g_0)^2\right]^{\frac{1}{2}}\\
    \leq~& \left(|C_Y-C_f|\|\hat{f}-f^*\|_{2}+C_f\|\Tcal(\hat{f}-f^*)\|_2\right)\error(\hat{g}),
\end{align*} where $\error(\hat{g}) =2(\frac{1}{c_0} + 1)^2\mathbb{E}\left[H(\hat{g},g_0)^2\right]^{\frac{1}{2}}$.

% For $\tilde{\cal E}(\hat{g})$, we have that:
% \begin{align*}
%     \tilde{\cal E}(\hat{g}) =~& (1-\nu){\cal E}(\hat{g})+\E[-2Y(\Tcal -\hat{\Tcal})(\hat{f}-f^*) + \sum_{j} (\hat{\theta}_j-\theta_j^*)((\Tcal f_j)^2 -(\hat{\Tcal}f_j)^2)]\\
%     = ~& (1-\nu){\cal E}(\hat{g}) +\E[-2Y(\Tcal -\hat{\Tcal})(\hat{f}-f^*) + \sum_j (\hat{\theta}_j-\theta_j^*)( \Tcal f_j +\hat{\Tcal}f_j)(\Tcal f_j -\hat{\Tcal}f_j)]\\
%     \leq~&  (1-\nu){\cal E}(\hat{g}) + C_Y \|\hat{f}-f^*\|\error(\hat{g}) + 2\sum_j (\hat{\theta}_j+\theta_j^*) C_f\E[|(\Tcal-\hat{\Tcal})f_j|]\\
%     \leq~& (1-\nu){\cal E}(\hat{g}) + C_Y \|\hat{f}-f^*\|\error(\hat{g}) + 2 C_f^2\error(\hat{g}).
% \end{align*}

\end{proof}

The proof of the above oracle inequality follows from a modified proof of Theorem 10 from \citet{lan2023causal}, which we reproduce here for completeness:

\begin{proof}
Let $\Hcal_M$ denote the convex hull over the candidate functions $\{f_1, \dots,f_M\}$. For brevity, we also denote $\hat{h} = h_{\hat{\theta}}$, and $h^* = h_{\theta^*}$. By a second order Taylor expansion around $h_{\theta^*_{\alpha}}$, there exist $\bar{h}\in\text{Star}(\Hcal_M,h^*)$ such that:
\begin{align*}
    \E\big[\ell_{\hat{h},\hat{g}} - \ell_{h^*,\hat{g}}\big] - D_{h}\E\left[\ell(h^*,\hat{g})\right][\hat{h}-h^*] =~& \frac{1}{2}D_{h}^2 \E\left[\ell_{\bar{h},\hat{g}}\right][\hat{h} - h^*,\hat{h}-h^*]\\
     =~& \|\hat{\Tcal} (\hat{h}-h^*)\|^2 + \alpha\|\hat{h}-h^*\|^2\\
     % \geq~& \alpha\|\hat{h}-h^*\|^2
\end{align*}

Now to bound the first derivative term:
\begin{align*}
    |D_{h}\E\ell(h^*,\hat{g})[\hat{h}-h^*] &-D_{h}\E\ell(h^*,g_0)[\hat{h}-h^*] |\\
    &= \left|\E\left[ 2(Y-\hat{\Tcal}\bar{h})\hat{\Tcal}(\hat{h}-h^*) -2(Y-{\Tcal}\bar{h}){\Tcal}(\hat{h}-h^*) \right] \right|\\
    &= \left|\E\left[ 2Y(\hat{\Tcal}-\Tcal)(\hat{h}-h^*) - 2\hat{\Tcal}(\bar{h})\hat{\Tcal}(\hat{h}-h^*)+2{\Tcal}(\bar{h}){\Tcal}(\hat{h}-h^*)\right] \right|\\
    &= \left|\E\left[ 2Y(\hat{\Tcal}-\Tcal)(\hat{h}-h^*) - 2(\hat{\Tcal}-\Tcal)(\bar{h})\hat{\Tcal}(\hat{h}-h^*)+2{\Tcal}(\bar{h})(\Tcal-\hat{\Tcal})(\hat{h}-h^*)\right] \right|\\
    & \leq 2\left((C_Y+C_h)\|\hat{h}-h^*\| + C_h\|\hat{\Tcal}(\hat{h}-h^*)\|\right)\error(\hat{g})\\
    % ~& \leq 2 (C_Y+2C_h)\|\hat{h}-h^*\| \error(\hat{g})
\end{align*}

% Note we can also bound:
% \begin{align*}
%     |D_{h}\E\ell(h^*,\hat{g})[\hat{h}-h^*] &-D_{h}\E\ell(h^*,g_0)[\hat{h}-h^*] |\\
%     &= \left|\E\left[ 2Y(\hat{\Tcal}-\Tcal)(\hat{h}-h^*) - 2\hat{\Tcal}(\bar{h})\hat{\Tcal}(\hat{h}-h^*)+2{\Tcal}(\bar{h}){\Tcal}(\hat{h}-h^*)\right] \right|\\
%     &= \left|\E\left[ 2Y(\hat{\Tcal}-\Tcal)(\hat{h}-h^*) - 2\hat{\Tcal}(\bar{h})(\hat{\Tcal}-\Tcal)(\hat{h}-h^*)+2({\Tcal}-\hat{\Tcal})(\bar{h})\Tcal(\hat{h}-h^*)\right] \right|\\
%     & \leq 2\left((C_Y+C_h)\|\hat{h}-h^*\| + C_h\|\Tcal(\hat{h}-h^*)\|\right)\error(\hat{g})\\
% \end{align*}

Applying the AM-GM inequality, we have that:
\begin{align*}
     -D_{h}\E\ell(h^*,\hat{g})[\hat{h}-h^*] \leq~& - D_{h}\E\ell(h^*,g_0)[\hat{h}-h^*]+\frac{1}{2}\left(\alpha\|\hat{f}-f'\|^2 + \|\hat{\Tcal}(\hat{f}-f')\|^2\right)\\
     +~&\left(\frac{2(C_Y+2C_f)^2}{\alpha}+2C_h^2\right)\error(\hat{g})^2
\end{align*}
Rearranging, we get:
\begin{align*}
    \alpha\|\hat{h}-h^*\|^2 + \|\hat{\Tcal}(\hat{h}-h^*)\|^2\leq - 2 D_{h}\E\ell(h^*,g_0)[\hat{h}-h^*]+ 2(\E\ell_{\hat{h},\hat{g}} - \E\ell_{h^*,\hat{g}}) +O\left( \frac{1}{\alpha} \error(\hat{g})^2\right) \label{eqn:norm_f}
\end{align*}
By Lemma \ref{lem:local-concen}, we have that:
\begin{align*}
    \E\left(\ell_{\hat{h},\hat{g}} - \ell_{h^*,\hat{g}}\right) \leq~& (\E - \E_n)\left(\ell_{\hat{h},\hat{g}} - \ell_{h^*,\hat{g}}\right)\\
    \leq~& O\big(\delta_{n,\Hcal_M}(\alpha \|\hat{h}-h^*\| + \|\hat{\Tcal}(\hat{h}-h^*)\|) + \delta_{n,\Hcal_M}^2\big)
\end{align*} where $\delta_{n,\Hcal_M}$ is the critical radius of the function class $\text{Star}(\Hcal_M,h^*)$.

Again applying AM-GM, we get:
\begin{align*}
    \alpha\|\hat{h}-h^*\|^2 + \|\hat{\Tcal}(\hat{h}-h^*)\|^2\leq - 2 D_{h}\E\ell(h^*,g_0)[\hat{h}-h^*]  +O\left( \frac{1}{\alpha} \error(\hat{g})^2 + \delta_{n,\Hcal_M}^2 + \alpha\delta_{n,\Hcal_M}^2\right)
\end{align*}
Since we can assume $\alpha\leq 1$, we have:
\begin{align*}
        \alpha\|\hat{h}-h^*\|^2 + \|\hat{\Tcal}(\hat{h}-h^*)\|^2\leq - 2 D_{h}\E\ell(h^*,g_0)[\hat{h}-h^*]  +O\left( \frac{1}{\alpha} \error(\hat{g})^2 + \delta_{n,\Hcal_M}^2 \right)
\end{align*}
Taking Taylor expansion at $h^*$, we have that there exist $\bar{h}\in\text{Star}(\Hcal_M,h^*)$ such that:
\begin{align*}
    \E\ell_{\hat{h},g_0} -  \E\ell_{h^*,g_0} = ~& D_{h}\E\ell_{h^*,g_0}[\hat{h}-h^*] + \frac{1}{2}D_{h}^2\E\ell_{\bar{h},g_0}[\hat{h}-h^*,\hat{h}-h^*]\\ 
    =~& D_{h}\E\ell_{h^*,g_0}[\hat{h}-h^*] + \alpha\|\hat{h}-h^*\|^2 + \|\Tcal(\hat{h}-h^*)\|^2 \\
    =~& D_{h}\E\ell_{h^*,g_0}[\hat{h}-h^*] + \alpha\|\hat{h}-h^*\|^2 + \|\hat{\Tcal}(\hat{h}-h^*)\|^2 \\
    -~& \|\hat{\Tcal}(\hat{h}-h^*)\|^2 + \|\Tcal(\hat{h}-h^*)\|^2\\
\end{align*}
Now let's consider the difference in the weak metrics in $\Tcal$ and $\Tcal$:
\begin{align*}
    \|\Tcal(\hat{h}-h^*)\|^2 - \|\hat{\Tcal}(\hat{h}-h^*)\|^2 =~& \E\left[(\Tcal(\hat{h}-h^*))^2-(\hat{\Tcal}(\hat{h}-h^*))^2\right]\\
    =~&\E\left[(\Tcal+\hat{\Tcal})(\hat{h}-h^*)(\Tcal-\hat{\Tcal})(\hat{h}-h^*)\right]\\
    \leq~& 4C_h\E\left[\left|(\Tcal-\hat{\Tcal})(\hat{h}-h^*)\right|\right]\\
    \leq~& 4C_h \|(\hat{h}-h^*)\|\error(\hat{g})\\
    \leq~& \alpha \|(\hat{h}-h^*)\|^2 + O\big(\frac{1}{\alpha}\error(\hat{g})^2\big)
\end{align*}
Then we have:
\begin{align*}
    \E\ell_{\hat{h},g_0} -  \E\ell_{h^*,g_0} 
    =~& D_{h}\E\ell_{h^*,g_0}[\hat{h}-h^*] + 2\alpha\|\hat{h}-h^*\|^2 + \|\hat{\Tcal}(\hat{h}-h^*)\|^2 \\
    +~&O\big(\frac{1}{\alpha}\error(\hat{g})^2\big)\\
    \leq~& -D_{h}\E\ell_{h^*,g_0}[\hat{h}-h^*]+  O\left(\delta_{n,\Hcal_M}^2  + \frac{1}{\alpha}\error(\hat{g})^2\right)\\
\end{align*}


Since $h^*$ satisfies the first order condition,
$D_{h}\E\ell_{h^*,g_0}[\hat{h}-h^*] \geq 0$, we can conclude that:
\begin{align*}
     \E\ell_{\hat{h},g_0} -  \E\ell_{h^*,g_0}\leq O\left(\delta_{n,\Hcal_M}^2  + \frac{1}{\alpha}\error(\hat{g})^2\right)
\end{align*}
Now it remains to show that $\delta_{n,\Hcal_M}^2 = O\big(\frac{M}{n}\big)$.
Let $\{e_1, \dots, e_{M'}\}$, $M'\leq M$, be an orthogonal basis of the linear subspace, $S$, spanned by the function class $\text{Star}(\Hcal_M,h^*)$, then we have:
\begin{align*}
    \E\left[\sup_{Ph^2\leq\delta_n^2 , h \in S}\left|\frac{1}{n}\sum_{i=1}^n \epsilon_i h(X_i)\right|\right]
    \leq~& \E\left[\sup_{\beta\in \mathbb{R}^{M'} , \|\beta\|_2 \leq \delta_n}\left|\frac{1}{n}\sum_{i=1}^n \epsilon_i \left(\sum_{j=1}^{M'}\beta_j e_j(X_i)\right)\right|\right]\\ \leq~&\delta_n\E\left(\sum_{j=1}^{M'}\left(\frac{1}{n}\sum_{i=1}^{n}\epsilon_i e_j(X_i)\right)^2\right)^{\frac{1}{2}} \tag{By the Cauchy-Schwarz inequality}\\
    \leq~& \delta_n\sqrt{\frac{M'}{n}}
\end{align*}
This shows that $\delta_{n,\Hcal_M}^2 = O\big(\frac{M}{n}\big)$, and we conclude the proof.
\end{proof}


% \begin{assumption}[Separability Condition]\label{Separability}
% Let $G$ be a real-valued function class defined on $\mathcal{Z}$. There exist $G_0\subset G$ such that $G_0$ is countable and for any $g \in G$, there exist a sequence $(g_k)_k$ such that for any $z\in \mathcal{Z}$, $(g_k(z))_k \rightarrow g(z)$ as $k \rightarrow \infty$.
% \end{assumption}
% \begin{lemma}[\cite{bartlett2006empirical}]\label{concentration}
%     There exist an absolute constant $c_0>0$ such that the following holds. Let $G$ be a real-valued function class defined on $\mathcal{Z}$ such that Assumption \ref{Separability} holds. Further assume that for all $g\in G$, $Pg^2\leq B_1 Pg$ and $\sup_{g\in G}\|g\|_{L_\infty}= B_2$ for constants $B_1$ and $B_2$. Let $\lambda_n$ be such that:
%     \begin{align*}
%         \E\left[\sup_{Pg\leq\lambda_n, g\in \text{Star}(G)}|(P-P_n)g|\right]\leq \frac{1}{8}\lambda_n.
%     \end{align*}
%     For every $\delta>0$, with probability greater than $1-\delta$, for every $g\in G$, we have:
%     \begin{align*}
%         |(P-P_n)g|\leq \frac{1}{2}\max\{Pg,\rho_n(\delta)\},
%     \end{align*}where $\rho_n(\delta)=\max \{\lambda_n, \frac{c_0(B_1+B_2)\log(1/\delta)}{n}\}$
% \end{lemma}
% \begin{proof}[Proof of Theorem \ref{convex-erm}]
% Consider the loss function:
% \begin{align}
%     \ell_{f,g}(Y,Z,X) = \underbrace{\left(Y - \int f(x)g(x|Z)dx\right)^2}_{\ell^{0}_{f,g}} + \underbrace{\alpha f(X)^2}_{r_f}. \label{Eqn:q_agg_loss}
% \end{align}
% Let $\mathcal{L}_{\Theta,g} = \{\mathcal{L}_{\theta,g}:\theta\in \Theta\}$, where $\mathcal{L}_{\theta,g} :=\ell_{\theta,g} - \ell_{\theta^*_\alpha,g} $. In addition, we can decompose $\mathcal{L}_{\theta,g}$ into $\mathcal{L}_{\theta,g}^0:=\ell^0_{\theta,g} - \ell_{\theta^*_\alpha,g}^0$ and  $\mathcal{L}_{\theta}^r:=r_{\theta}-r_{\theta^*_\alpha}$, such that $\mathcal{L}_{\theta,g} = \mathcal{L}_{\theta,g}^0+\mathcal{L}_{\theta}^r$.
% \begin{align*}
%     R(\hat{\theta},g_0) - R(\theta^*_{\alpha}, g_0) =~& P\mathcal{L}_{\hat{\theta},g_0}\\
%     =~& P\mathcal{L}_{\hat{\theta},\hat{g}} - (P\mathcal{L}_{\hat{\theta},\hat{g}}-P\mathcal{L}_{\hat{\theta},g_0})\\
%     \leq~& P\mathcal{L}_{\hat{\theta},\hat{g}} + O(\error(\hat{g})) \tag{By Lemma \ref{lemma:diffndiff}}\\
%     =~& P\mathcal{L}^0_{\hat{\theta},\hat{g}} + P\mathcal{L}^r_{\hat{\theta}} + O(\error(\hat{g}))\\
%     \leq~& 2 P_n\mathcal{L}^0_{\hat{\theta},\hat{g}} + 2P_n\mathcal{L}^r_{\hat{\theta}} + \rho^0_n(\delta) + \rho^r_n(\delta) +O(\error(\hat{g})) \tag{by Lemma \ref{concentration}}\\
%     =~&2 P_n \mathcal{L}_{\hat{\theta},\hat{g}} + \rho^0_n(\delta) + \rho^r_n(\delta) +O(\error(\hat{g}))\\
%     \leq~& \rho^0_n(\delta) + \rho^r_n(\delta) +O(\error(\hat{g})).
% \end{align*}
% It remains to show that $\rho^0_n(\delta) $ and $ \rho^r_n(\delta)$ are on the order of $\frac{M}{n}$. First, we observe that both function classes $\mathcal{L}^{0}_{\Theta,\hat{g}}$ and $\mathcal{L}^{r}_{\Theta}$ satisfies the separability condition. Since $Y$ and $f$ are bounded by $C_Y$ and $C_f$ respectively, we have that $P(\mathcal{L}^{0}_{\theta,\hat{g}})^2\leq B_1 P\mathcal{L}^{0}_{\theta,\hat{g}}$ and $P(\mathcal{L}^{r}_{\theta})^2\leq B_2 P\mathcal{L}^{r}_{\theta}$ for constants $B_1$ and $B_2$.
% Let $\mathcal{L}_{\Theta}$ denote either $\mathcal{L}^{0}_{\Theta,\hat{g}}$ and $\mathcal{L}^{r}_{\Theta}$. Then by the peeling argument in \cite{bartlett2012}, for any $\lambda>0$:
% \begin{align*} 
% \left\{f\in\text{Star}(\mathcal{L}_{\Theta}), P(f)\leq \lambda\right\}\subset \cup_{i=0}^{\infty}\left\{\alpha f: 0\leq \alpha \leq 2^{-i}, f\in \mathcal{L}_{\Theta}, Pf\leq 2^{i+1}\lambda\right\},
% \end{align*}
% which implies
% \begin{align*}
% \E\left[\sup_{P\ell\leq\lambda, f \in \mathcal{L}_{\Theta}}|(P-P_n)f|\right]\leq \sum_{i=0}^{\infty}2^{-i}\E\left[\sup_{Pf\leq 2^{i+1}\lambda, f \in \mathcal{L}_{\Theta}}|(P-P_n)f|\right].
% \end{align*}
% Thus, if $\lambda_n/8$ upper bounds the RHS, then it also satisfies the condition in Lemma \ref{concentration}. 
% Let $\{e_1, \dots, e_{M'}\}$, $M'\leq M$, be an orthogonal basis of the linear subspace, $S$, spanned by the function class $\mathcal{F}$, then we have:
% \begin{align*}
%     \E\left[\sup_{Pf\leq\lambda, f \in \mathcal{L}_{\Theta}}|(P-P_n)f|\right] \leq~& \max\{1,\alpha\}8b \E\left[\sup_{Pf^2\leq\lambda , f \in S}\left|\frac{1}{n}\sum_{i=1}^n \epsilon_i f(X_i)\right|\right] \tag{By Symmetrization}\\
%     \leq~& \max\{1,\alpha\}8b \E\left[\sup_{\beta\in \mathbb{R}^{M'} , \|\beta\|_2 \leq \sqrt{\lambda}}\left|\frac{1}{n}\sum_{i=1}^n \epsilon_i \left(\sum_{j=1}^{M'}\beta_j e_j(X_i)\right)\right|\right]\\
%     \leq~& \max\{1,\alpha\}8b\sqrt{\lambda}\E\left(\sum_{j=1}^{M'}\left(\frac{1}{n}\sum_{i=1}^{n}\epsilon_i e_j(X_i)\right)^2\right)^{\frac{1}{2}} \tag{By the Cauchy-Schwarz inequality}\\
%     \leq~& \max\{1,\alpha\}8b\sqrt{\frac{M'\lambda}{n}}
% \end{align*}
% Thus, by the pealing argument, we showed that $\lambda_n = O\left(\frac{M}{n}\right)$.
% \end{proof}



\subsubsection{Auxiliary Lemmas}
\begin{lemma}\label{lem:exp-tail}
For any random variable $Z$, if $\mathbb{E}\left[\exp\left\{ Z\right\} \right] \leq 1$, then $\Pr(Z \geq \log(1/\delta)) \leq \delta$.
\end{lemma}
\begin{proof}
Suppose that one shows that:
\begin{align}
    \mathbb{E}\left[\exp\left\{ Z\right\} \right] \leq 1
\end{align}
Then by Markov's inequality:
\begin{align}
    \Pr(Z \geq x) = \Pr(\exp(Z) \geq \exp(x)) \leq \frac{\mathbb{E}[\exp(Z)]}{\exp(x)}\leq \frac{1}{\exp(x)} = \exp(-x)
\end{align}
Thus:
\begin{align}
    \Pr(Z\geq \log(1/\delta)) \leq \delta
\end{align}
\end{proof}

\begin{lemma}\label{lem:hoeffding}
Suppose $\epsilon_1, \ldots, \epsilon_n$ are independent Rademacher random variables drawn equiprobably in $\{-1, 1\}$ and that $a_1,\ldots, a_n$ are real numbers. Then:
\begin{align}
    \mathbb{E}_{\epsilon_{1:n}}\left[\exp\left\{\frac{1}{n} \sum_i \epsilon_i a_i \right\}\right] \leq \exp\left\{\frac{1}{2n^2} \sum_i a_i^2 \right\}
\end{align}
\end{lemma}
\begin{proof}
First note that for by the derivation in Example 2.2 of \url{https://www.stat.berkeley.edu/~mjwain/stat210b/Chap2_TailBounds_Jan22_2015.pdf}, we have that:
\begin{align}
    \mathbb{E}[\exp(\epsilon_i a_i/n)] \leq \exp(a_i^2/2n^2)
\end{align}
Invoking now independence of the $\epsilon_i$ we have:
\begin{align}
    \mathbb{E}_{\epsilon_{1:n}}\left[\exp\left\{\frac{1}{n}\sum_i \epsilon_i a_i \right\}\right] =~& \mathbb{E}_{\epsilon_{1:n}}\left[\prod_{i=1}^n \exp\left\{\epsilon_i a_i/n \right\}\right]\\
    =~& \prod_{i=1}^n \mathbb{E}_{\epsilon_{1:n}}\left[ \exp\left\{\epsilon_i a_i/n \right\}\right] \tag{independence}\\
    \leq~& \prod_{i=1}^n \exp\left(\frac{a_i^2}{2n^2}\right) = \exp\left\{\frac{1}{2n^2} \sum_i a_i^2 \right\}
\end{align}
\end{proof}

The following proposition can also be proven as an application of Bernstein's inequality.
\begin{lemma}[Proposition 1, \cite{Rigollet2014}]\label{lem:bernstein}
Let $V_1,\ldots, V_n$ be i.i.d. random variables, such that $|V_i|\leq c$ a.s.. Let $c_0>0$, then for any $0< \lambda < \frac{2c_0}{1+ 2c_0 c}$:
\begin{align}
    \mathbb{E}\left[\exp\left\{n\lambda \left( \frac{1}{n}\sum_i (V_i - \mathbb{E}[V_i] - c_0 \mathbb{E}[V_i^2]\right)\right\}\right] \leq 1  
\end{align}
\end{lemma}

\begin{lemma}[Symmetrization Inequality \citep{Rigollet2014}]\label{sym}
    let $\mathcal{F}$ be a function class, $A_f, f \in \mathcal{F}$ be a given function on $\mathcal{F}$ and $\Phi$ be a convex non-decreasing function then
    $$
    \mathbb{E} \Phi\left(\sup _{f \in \mathcal{F}}\left[\mathbb{E}[f(Z)]-\frac{1}{n}\sum_{i=1}^{n} f(Z_i)-A_f\right]\right) \leq \mathbb{E} \Phi\left(2 \sup _{f \in \mathcal{F}}\left[\frac{1}{n}\sum_{i=1}^{n} \varepsilon_i f(Z_i)-A_f\right]\right)
    $$
\end{lemma}

\begin{lemma}[Contraction Inequality, Theorem 4.12 in \citep{ledoux1991probability}]\label{contr}
    Let $F: \mathbb{R}_{+} \rightarrow \mathbb{R}_{+}$be convex and increasing. Let further $\varphi_i: \mathbb{R} \rightarrow \mathbb{R}, i \leq N$, be such that $|\varphi_i(s) - \varphi_i(t)|\leq |s-t|$ and $\varphi_i(0)=0$. Then, for any bounded subset $T$ in $\mathbb{R}^N$
    $$
    \mathbb{E} F\left(\frac{1}{2}\sup_{T}\sum_{i=1}^N \varepsilon_i \varphi_i\left(t_i\right)\right) \leq \mathbb{E} F\left(\sup_{T}\sum_{i=1}^N \varepsilon_i t_i\right)
    $$
\end{lemma}


\section{Remaining stuff}

(Should be done) 
\begin{itemize}
    \item With iteration (Done)
    \item Experiments (Proximal causal inference) 
    \item MLE analysis with critical radius when function classes are infinite  (done)
\end{itemize}

(Might be difficult)
\begin{itemize}
    \item Guarantees in terms of model selection 
    \item Lower bound 
\end{itemize}

\bibliographystyle{chicago}
\bibliography{rl}

\end{document} 
